% ------------------------------------------------------------------------
% ------------------------------------------------------------------------
% ICMC: Modelo de Trabalho Acadêmico (tese de doutorado, dissertação de
% mestrado e trabalhos monográficos em geral) em conformidade com 
% ABNT NBR 14724:2011: Informação e documentação - Trabalhos acadêmicos -
% Apresentação
% ------------------------------------------------------------------------
% ------------------------------------------------------------------------

% Opções: 
%   Qualificação          = qualificacao 
%   Curso                 = doutorado/mestrado
%   Situação do trabalho  = pre-defesa/pos-defesa (exceto para qualificação)
%   Versão para impressão = impressao
\documentclass[mestrado, pre-defesa]{packages/icmc}

% ---------------------------------------------------------------------------
% Pacotes Opcionais
% ---------------------------------------------------------------------------

\usepackage{rotating}           % Usado para rotacionar o texto
% Float rotacionado para o ambiente `quadro` (equivalente ao `sidewaystable`,
% preservando a numeração/nomenclatura de Quadros do abnTeX/ICMC).
\makeatletter
\newenvironment{sidewaysquadro}{\@rotfloat{quadro}}{\end@rotfloat}
\makeatother
\usepackage[all,knot,arc,import,poly]{xy}   % Pacote para desenhos gráficos
\usepackage{multirow}
\usepackage{booktabs}
% Incluir pacote para tcolorbox
% \usepackage{tcolorbox}
% Este pacote pode conflitar com outros pacotes gráficos como o ``pictex''
% Então é necessário usar apenas um dos pacotes conflitantes
\newcommand{\VerbL}{0.52\textwidth}
\newcommand{\LatL}{0.42\textwidth}

%\usepackage{easyReview} % for \comment, \alert, etc.

% ---------------------------------------------------------------------------


% Definicao de novos comandos e pacotes ===============================

\usepackage{pdflscape} % Orientacao horizontal da pagina

\usepackage{array} % Array para centralizacao de colunas
\newcolumntype{P}[1]{>{\centering\arraybackslash}p{#1}}
\newcolumntype{M}[1]{>{\centering\arraybackslash}m{#1}}

% Fim da secao ============================================

% ---
% Informações de dados para CAPA e FOLHA DE ROSTO
% ---
% Tanto na capa quanto nas folhas de rosto apenas a primeira letra da primeira palavra (ou nomes próprios) devem estar em letra maiúscula, todas as demais devem ser em letra minúscula.
\tituloPT{Técnicas de Otimização na Alocação de Médicos em Unidades de Diagnóstico por Imagem}
\tituloEN{Optimization Techniques for Physician Scheduling in Diagnostic Imaging Units}
\autor[Mariotti, T. R.]{Thiago Rafael Mariotti Claudio}
\genero{M} % Gênero do autor (M = Masculino / F = Feminino)
\orientador[Orientadora]{Profa. Dra.}{Maristela Oliveira dos Santos}
%\coorientador{Prof. Dr.}{Fulano de Tal}
\curso{CCMC}
\data{10}{06}{2026} % Data do depósito
\idioma{PT} % Idioma principal do documento (PT = português / EN = inglês)
% ---


% ---
% RESUMOS
% ---

% Resumo em PORTUGUÊS
% conter no máximo 500 palavras
% conter no mínimo 1 e no máximo 5 palavras-chave
\textoresumo[brazil]{
A geração de escalas em redes de clínicas de exames de diagnóstico por imagem é uma atividade combinatória custosa, marcada por demandas heterogêneas, disponibilidades irregulares e pela coordenação de múltiplas unidades autônomas que compartilham profissionais. Este trabalho propõe e avalia uma formulação de Programação Inteira Mista bi-objetiva para o problema, minimizando simultaneamente exames não atendidos e trocas de clínica ao longo do horizonte de planejamento. A modelagem incorpora uma variável de estado adaptada do \textit{General Lot-sizing and Scheduling Problem} para detectar trocas mesmo em agendas com lacunas, e a estrutura bi-objetiva é resolvida pelo método $\varepsilon$-restrito. Para tanto, três famílias metodológicas são parametrizadas e comparadas sobre dados reais de 250 médicos e 37 unidades operadas em 2022 e 2023: o método exato com o \textit{solver} Gurobi; matheurísticas Relax-and-Fix e Fix-and-Optimize, isoladas e em arquitetura híbrida; e o algoritmo evolutivo NSGA-II nas variantes Estendida e Compacta, com operadores customizados de amostragem com consciência de restrições, cruzamento por blocos diários e mutação por realocação de \textit{slot}. Avalia-se o desempenho em quatro horizontes operacionais (15, 30, 60 e 90 dias), sob protocolo de dois estágios com instâncias de parametrização e validação. Os resultados evidenciam transição de regime entre 15 e 30 dias: o método exato atinge otimalidade apenas no ciclo quinzenal, e a partir de 30 dias o crescimento combinatório inviabiliza a convergência sob orçamento de uma hora por ponto da fronteira. O híbrido Relax-and-Fix seguido de Fix-and-Optimize mantém hipervolume competitivo em horizontes curtos, enquanto o NSGA-II Compacto, ao reduzir a dimensionalidade do cromossomo para a associação diária médico-clínica, supera a variante Estendida e ultrapassa a fronteira do exato em horizontes longos, consolidando-se como a abordagem mais robusta sob restrições estritas de tempo. As soluções meta-heurísticas apresentam violação parcial da cobertura mínima, indicando direção natural para o desenvolvimento de operadores de reparo. O estudo contribui para a literatura de escalonamento médico ao tratar conjuntamente múltiplas unidades autônomas e minimização explícita de remanejamentos, configuração pouco explorada em trabalhos similares.}
{Escalonamento de Médicos, Otimização Bi-objetiva, Matheurísticas, MIP-Heuristicas, Modelagem Matemática}


% resumo em INGLÊS
% conter no máximo 500 palavras
% conter no mínimo 1 e no máximo 5 palavras-chave
\textoresumo[english]{
Scheduling physicians in networks of diagnostic imaging clinics is a costly combinatorial activity, marked by heterogeneous demand, irregular availabilities, and the coordination of multiple autonomous units that share specialists. This work proposes and evaluates a bi-objective Mixed Integer Programming formulation for the problem, simultaneously minimising unmet exams and clinic swaps across the planning horizon. The model incorporates a state variable adapted from the General Lot-sizing and Scheduling Problem to detect swaps even in agendas with gaps, and the bi-objective structure is solved through the $\varepsilon$-constraint method. Three method families are parametrised and compared on real-world data of 250 physicians and 37 units operated in 2022 and 2023: the exact method with the Gurobi solver; Relax-and-Fix and Fix-and-Optimize matheuristics, both standalone and in a hybrid architecture; and the NSGA-II evolutionary algorithm in Extended and Compact variants, with custom operators for constraint-aware sampling, day-block crossover, and slot-reassignment mutation. The evaluation spans four operational horizons (15, 30, 60, and 90 days), under a two-stage protocol with parametrisation and validation instance sets. Results show a regime transition between 15 and 30 days: the exact method achieves optimality only in the biweekly cycle, and from 30 days onwards combinatorial growth prevents convergence within one hour per Pareto-frontier point. The hybrid Relax-and-Fix followed by Fix-and-Optimize retains competitive hypervolume in short horizons, while the Compact NSGA-II, by reducing chromosome dimensionality to the daily physician-clinic assignment, outperforms the Extended variant and surpasses the exact frontier in longer horizons, establishing itself as the most robust approach under strict time budgets. Metaheuristic solutions exhibit partial violation of the minimum coverage constraint, indicating a natural direction for repair-operator development. The study contributes to the physician scheduling literature by jointly addressing multiple autonomous units and the explicit minimisation of reassignments, a setting rarely covered in similar works.}
{Physician Scheduling, Bi-objective Optimisation, Matheuristics, Matheuristics, Mathematical Modeling}


% ----------------------------------------------------------
% ELEMENTOS PRÉ-TEXTUAIS
% ----------------------------------------------------------

% Inserir a ficha catalográfica
\incluifichacatalografica{tex/1-pre-textual/ficha-catalografica.pdf}

% DEDICATÓRIA / AGRADECIMENTO / EPÍGRAFE
\textodedicatoria*{tex/1-pre-textual/dedicatoria}
\textoagradecimentos*{tex/1-pre-textual/agradecimentos}
\textoepigrafe*{tex/1-pre-textual/epigrafe}

% Inclui a lista de figuras
\incluilistadefiguras

% Inclui a lista de tabelas
\incluilistadetabelas

% % Inclui a lista de quadros
\incluilistadequadros

% % Inclui a lista de algoritmos
% \incluilistadealgoritmos

% % Inclui a lista de códigos
% \incluilistadecodigos

% % Inclui a lista de siglas e abreviaturas
% \incluilistadesiglas

% % Inclui a lista de símbolos
% \incluilistadesimbolos

% ----
% Início do documento
% ----
\begin{document}
% ----------------------------------------------------------
% ELEMENTOS TEXTUAIS
% ----------------------------------------------------------
\textual

\chapter{Introdução} % 1
\label{cap:introducao}

Exames de diagnóstico por imagem, como ultrassonografia e radiologia, são procedimentos comuns no cotidiano popular e na rotina de saúde, como evidenciado na notícia de \citeonline{MedicinaSA2024}, que aponta um total de 2,4 bilhões de exames diagnósticos realizados. O setor clínico de diagnósticos vêm se tornando um alicerce não apenas da saúde, como tendência médica em análises de caso, mas também um pilar econômico, com geração de receita bruta de R\$ 24 bilhões movimentados no ano de 2022 \cite{Abramed_2023}, sendo metade desse valor realizado pela saúde privada.
Autores como \citeonline{Yasui_Fujiwara_Okitsu_Yonemitsu_Ito_Yokosuka_2012} e \citeonline{OBRIEN_2007} avaliam a importância do uso de diagnóstico por exames para constatação precoce de doenças e elevação da qualidade no acompanhamento médico em detrimento ao custo relativo de tais operações, ponto comum de discussão, uma vez que os custos operacionais para uso e manutenção da aparelhagem necessária é elevado. Entretanto, como aponta \citeonline{CBDL_2024}, o custo da operação vêm caindo anualmente, como exemplificada pela queda de 22,9\% para 21\% no ano de 2023, influindo diretamente nos preços dos exames e estabilidade funcional.

Apesar de sua importância, profissionais capacitados nas áreas de radiologia e ultrassonografia estão em falta para operação da aparelhagem. Com uma estimativa de crescimento de 6\% para 18\% nas vagas não preenchidas por especialistas na área, espera-se um grande crescimento nas filas de atendimento dessa modalidade, devido o envelhecimento populacional global e as decorrentes necessidades de acompanhamento na saúde\cite{RSNA_2024}. No Brasil, verifica-se uma realidade de escassez generalizada para profissionais da área: para cada 1,99 milhões de indivíduos estimam-se apenas 1,51 mil médicos no ano de 2017, como apontam \citeonline{Liebel_Dias_Schneider_SaScielo2021}.

O excedente de demanda de exames pelos pacientes e a carência de disponibilidade de atendimento, causada pela escassez de mão de obra, geram custos na operação devido ao maquinário parado e ao afunilamento no atendimento. Além disso, segundo relatório público da \citeonline{ANS_2024}, embora o prejuízo fiscal da saúde suplementar tenha reduzido significativa, as margens operacionais continuam próximas ao limiar do adequado, pressionando as empresas do ramo a procurar formas de cortar custos.

Uma possibilidade para o controle de gastos, através do incremento na qualidade operacional, é o escalonamento da força de trabalho e a elaboração de escalas médicas otimizadas para atender altas demandas por equipes reduzidas. Autores como \citeonline{Rousseau_Pesant_Gendreau_2002} apontam que a geração de quadros de horário com boa qualidade podem aumentar significativamente a capacidade de atendimento de pacientes, promovendo a redução nas filas de espera, aumentando a satisfação do cliente, tudo isso enquanto mantém o tamanho da equipe baixa. Além disso, \citeonline{Robinson_Chen_2003} citam também a importância do processo para melhor aproveitamento do maquinário disponível, reduzindo tempo ocioso do equipamento e, impulsionando a operação. 
Entretanto, embora altamente benéfica, a confecção de escalas e quadros de horário, comumente realizada de maneira manual, é uma tarefa laboriosa.

Diante dos pontos discutidos, somados à crescente demanda populacional por serviços médicos de diagnóstico apontada por \citeonline{Liebel_Dias_Schneider_SaScielo2021}, o processo de modelagem matemática e desenvolvimento de algoritmos computacionais para confecção de escalas se torna evidente e necessário. Esta dissertação de mestrado visa estudar o problema de escalonamento da equipe médica em contextos de consultas pré-agendadas, onde se conhece previamente a demanda a ser atendida.

O referido estudo é baseado em dados reais disponibilizados por uma rede de clínicas de saúde suplementar do ramo de diagnóstico por imagem com especialidade em ultrassonografia localizada em uma grande cidade do estado de São Paulo. Observa-se conjunto de dados uma média de 37.000 horários semanais para consulta distribuídos entre 37 unidades, que, em média, disponibilizam até 4 salas individuais para exames, e 294 médicos especialistas associados disponíveis para alocação para realização dos exames. 

Nesse cenário, os profissionais têm uma disponibilidade de horas para realizar exames em um determinado período do dia, além da possibilidade de ausência programada, gerando a necessidade de realocar outro especialista (possivelmente de outra unidade) para suprir a demanda. Contudo, flexibilizar as trocas constantes de profissionais entre unidades para atender a demanda pode resultar em descontentamento da equipe devido a muitas mudanças ao longo do planejamento. Assim, este trabalho visa a alocação de médicos nas unidades de saúde em que atendem, buscando atender os exames e minimizar essas mudanças entre unidades. Todavia, não é comum na literatura acadêmica explorar e estudar modelagens referentes à minimização de trocas de médicos entre unidades em problemas de escalonamento médico, especialmente ao se considerar a possibilidade de múltiplas unidades compartilharem profissionais entre si. Desta forma, investiga-se o problema no contexto de otimização bi-objetivo, buscando alocar os profissionais de modo a minimizar os exames não atendidos enquanto flexibiliza as trocas entre clínicas. No entanto, também se busca minimizar as trocas entre unidades. Para lidar com o problema, foi proposto um modelo de otimização inteira, além de duas abordagens heurísticas baseadas na formulação matemática (Relax-and-Fix e Fix-and-Optimize) e uma meta-heurística evolucionária multiobjetivo (NSGA-II).

Assim, neste trabalho têm-se como objetivos principais: (1) a modelagem e representação matemática do escalonamento da equipe médica, (2) a confecção dos quadros de horário dos médicos por métodos exatos do cenário mediante demanda e disponibilidade já conhecidas através do método determinístico, (3) o estudo e implementação de técnicas de otimização bi-objetivo que explorem a minimização de remanejamentos e trocas de profissionais entre unidades, e, (4) contribuir para a literatura acadêmica e o conhecimento científico através do estudo de sistemas que abordem a alocação em múltiplas unidades ao mesmo tempo que se reduz a troca de clínicas para o profissional ao longo do horizonte de planejamento.

Espera-se que, abordando tais desafios, a dissertação contribua para o aprimoramento do saber científico e acadêmico, bem como melhorias significativas fornecendo modelos sistêmicos bem estruturados para gestão operacional de instituições médicas, tornando-as mais eficientes, impactando diretamente na eficiência do serviço prestado e na satisfação do médico e paciente. 

Para tal, o referido estudo encontra-se estruturado da seguinte maneira: o presente Capítulo \ref{cap:introducao} discorre o contexto a qual o problema está inserido e apresenta a proposta de intervenção pautada no cenário descrito. No Capítulo \ref{cap:revisao-bibliografica} é realizada a Revisão da Literatura em Pesquisa Operacional Aplicada à Saúde enquanto se exploram lacunas no conhecimento acadêmico e metodológico. O Capítulo \ref{cap:problema} descreve de maneira aprofundada o problema observado mediante, descrevendo a característica e classificação do tema estudado. No Capítulo \ref{cap:mip-heuristicas} descreve o modelo matemático proposto para descrever o assunto observado, bem como expandir as diversas técnicas computacionais empregadas para obtenção de resultados. O Capítulo \ref{cap:resultados} por sua vez, apresenta o conjunto de dados e explica o ambiente computacional utilizado para executar a metodologia proposta, ao mesmo tempo, em que expõe e analisa os resultados observados nos experimentos. Por fim, o \ref{cap:fechamento} tem por finalidade apresentar conclusões extraídas do estudo e sugerir melhorias para com o projeto, e possíveis cenários para estudo no tema.

\chapter{Revisão Bibliográfica} % 2
\label{cap:revisao-bibliografica}

Na literatura acadêmica, observa-se uma lacuna significativa em estudos que aplicam técnicas de Otimização ao contexto de Alocação de Médicos em Clínicas de Exames. Esta carência é especialmente evidente nos trabalhos de \citeonline{Abdalkareem2021_ac} e \citeonline{ERHARD20181}, que, ao analisarem a literatura existente, destacam a falta de contribuições focadas especificamente nesse problema, enquanto o trabalho de \citeonline{Ngoo_Goh_Sze_Sabar_Abdullah_Kendall_2022} reforça a tendência contemporânea de estudos em problemas de Alocação de Enfermeiras (os).

O estudo de \citeonline{Abdalkareem2021_ac} apresenta uma análise abrangente de 190 trabalhos publicados entre 1995 e 2016, classificados de acordo com diferentes contextos. Entre as categorias destacam-se: o Problema de Agendamento para Admissão de Pacientes (PASP, do inglês \textit{Patient Admission Scheduling Problem}), o Problema de Alocação de Enfermeiras(os) (NRP, do inglês \textit{Nurse Rostering Problem}), o Agendamento de Salas de Operações (ORSP, do inglês \textit{Operation Room Scheduling Problem}), e uma categoria genérica denominada Outros, que abrange problemas como Alocação de Médicos(as), Escalas em Telemedicina e Assistência Médica Domiciliar.

Para cada um desses problemas, os autores analisam as metodologias utilizadas, que variam entre métodos exatos de programação matemática, técnicas heurísticas e meta-heurísticas. Além disso, examinam o tipo de base de dados empregado nas pesquisas, sejam elas baseadas em dados reais, simulações ou instâncias artificiais. O estudo revela que, na categoria ``Outros'', apenas três trabalhos tratam especificamente do problema de Alocação de Médicos(as), e todos eles focam em contextos hospitalares, empregando como principal metodologia a programação inteira mista.

Em um contexto mais focado, o trabalho de \citeonline{Ngoo_Goh_Sze_Sabar_Abdullah_Kendall_2022} busca analisar e avaliar a vasta literatura dos problemas de Alocação de Enfermeiras(os), focando em 50 publicações entre os anos de 2012 a 2021 com o intuito de entender as tendências e metodologias aplicadas para a solução dos problemas levantados. Além da definição formal do NRP e observações quanto a modelagem básica do problema, dada por sua extensibilidade e complexidade em restrições que buscam conciliar preferência com demanda e restrições de tempo. Também é analisada a tendência em métodos de solução, principalmente para programação linear inteira como as  mat-heurísticas (do inglês \textit{matheuristics}) aplicadas em contextos de estágio único (planejamento de horizontes de eventos determinísticos) e multiestágio (planejamento de horizonte de eventos estocásticos).

O trabalho de \citeonline{ERHARD20181} também visa categorizar e avaliar estudos, dessa vez aplicados ao contexto de alocação de médicos(as) por meio de métricas quantitativas. O artigo diferencia os trabalhos conforme o horizonte de planejamento: problemas de alocação estratégica, que tratam de decisões de longo prazo, como a definição e quantificação da equipe necessária; problemas de escalonamento, que atuam no nível tático, criando escalas de trabalho para o médio prazo; e problemas de replanejamento, que abordam ajustes de curto prazo, com foco na adaptação de escalas em função de demandas não previstas. A revisão de \citeonline{ERHARD20181} abrange 68 publicações entre 1985 e 2016.

O estudo também aponta que a modelagem matemática desses problemas, fora do escopo puramente financeiro, apresenta algumas características comuns. Dentre elas destaca-se o atendimento da demanda, a definição de limites inferiores e superiores na carga horária (referido no texto como corredor de carga horária), a garantia de descanso mínimo e o equilíbrio na distribuição da carga de trabalho entre a equipe médica. Além disso, os autores identificam a quantidade limitada de pesquisas focadas na área médica hospitalar, especialmente em comparação com os estudos sobre alocação de enfermeiras (os), uma tendência já observada no trabalho de \citeonline{Abdalkareem2021_ac}. Por fim, \citeonline{ERHARD20181} ressaltam algumas lacunas na literatura existente, como a falta de flexibilidade nas escalas e a pouca atenção dada à integração de turnos flexíveis e pausas, abrindo espaço para futuras investigações e aprimoramentos na área.

\section{Sobre a Literatura de Alocação de Médicos(as)}
\label{2rev-sec1}

Problemas de Alocação de Médicos(as) (do inglês \textit{Physician Allocation Problems}, abreviado como PAP) podem ser datados na literatura acadêmica desde o século passado, com trabalhos como o de \citeonline{Canter_1977} que realiza simulação para agendamento paciente-médico(a). Apesar disso, investigações voltadas a métodos de otimização, na maioria das vezes, enfatizavam contextos amplos, como a composição (\textit{rostering}) de equipes hospitalares, ou em alocações de enfermeiras (os), como o trabalho de \citeonline{Fitzpatrick_Farrell_Richter_Zeunik_1987}, com foco em suporte à tomada de decisão auxiliada por computador, com uma breve menção a possibilidade de implementação de métodos de otimização para solução da rotação de alocação de leitos.

\citeonline{Bruni_Detti_2014} avaliam o uso de modelos exatos por meio da implementação de um MIP (Programação Inteira Mista, do inglês \textit{Mixed Integer Programming}) no contexto da saúde, buscando formalizar uma solução generalista aplicável em hospitais, clínicas e outros. Embora a natureza de problemas de alocação de agentes sejam considerados NP-difícil \cite{Even_Itai_Shamir_1976}, o modelo proposto apresenta flexibilidade o suficiente para aplicação em diferentes contextos sem quebra de restrições obrigatórias (também referidas como \textit{hard constraints} na literatura) para sua solução mat-heurística baseada em \textit{branch-and-cut} para gerar soluções com melhoras significativas em relação à alocação manual. Testado como aplicação real no hospital universitário Umberto I, em Roma, a solução proposta melhorou consideravelmente a equidade nas alocações, reduzindo cargas de trabalho injustas (como múltiplas alocações do mesmo funcionário à suscetivos turnos noturnos) com eficiente boa eficiência computacional para aplicação em horizontes de planejamento extensos em poucos segundos.

A produção de \citeonline{Beaulieu_Ferland_Gendron_Michelon_2000} é uma das primeiras a tratar especificamente dos PAPs por meio de métodos de programação matemática e otimização. Ao conduzir um estudo de caso no hospital \textit{Sacré-Coeur}, a tarefa de gerenciar e escalar os horários dos 20 médicos(as) associados(as) se mostrou complexa devido a restrições obrigatórias complexas, tais como limites de turnos consecutivos, ergonomia de trabalho (como um período mínimo de descanso), períodos de férias, e restrições flexíveis, como preferências individuais por horários e dias de trabalho, embora a classificação das restrições dependa diretamente do hospital, dificultando a generalização do modelo. O trabalho propõe como metodologia o uso de programação linear inteira  multiobjetivo, resolvido através do método \textit{branch-and-bound}, buscando minimizar as violações ao cronograma, respeitando tanto as restrições rígidas quanto tentando atender as flexíveis.

Similarmente \citeonline{Topaloglu_2006} propõe um modelo exato baseado em \textit{goal programming} (método de programação linear multiobjetivo com foco em solucionar objetivos conflitantes) para solucionar o problema de escala em alas hospitalares emergenciais, focando especificamente na população médica residente, etapa necessária para obtenção do título de especialista em medicina. O estudo foi desenvolvido em parceria com um hospital universitário em sua EMR (ala de medicina residente emergencial, do ingles \textit{emergency medicine resident}). Devido sua natureza, o funcionamento da ala é ininterrupto, com grande fluxo de pacientes, gerando um desafio difícil para o controle e escalonamento de profissionais quando considerado regras ergonômicas de descanso, períodos de trabalho, considerações estudantis e preferências. Para gerir essas restrições foi implementado o uso de AHP (Processos Analíticos de Hierarquia, do inglês \textit{Analytical Hierarchy Process}) para gerir e priorizar o uso e atendimento de restrições.
Aplicado em uma base de dados de 15 meses, a abordagem proposta apresentou bons resultados em termos de equidade na escala de trabalho e conformidade quanto as restrições obrigatórias, principalmente quando comparados ao método manual na maior instância disponibilizada (5633 restrições e 23 residentes para um período de planejamento de 9 meses), embora o modelo não atingisse otimalidade nas especificações computacionais, especialmente tempo.

\citeonline{Beaulieu_Ferland_Gendron_Michelon_2000} e \citeonline{Topaloglu_2006}, quando comparadas ao escalonamento manual durante o período de implementação, tenham resultado em melhorias significativas na distribuição de horários e na redução do descumprimento de restrições, é importante destacar que o contexto do trabalho se aplica a salas de emergência, onde a demanda estática — isto é, demanda previamente conhecida e fixada no horizonte de planejamento — é irrelevante. Ou seja, o(a) médico(a) deve ser alocado(a) à sala independentemente de haver pouca demanda na ala emergencial, pois a chegada de pacientes é imprevisível.

Passando para um novo contexto, o estudo de \citeonline{Gunawan_Lau_2013} aborda o MPSP (Problema de Alocação do(a) Médico(a) Chefe, do inglês \textit{Master Physician Scheduling Problem}), que envolve a atribuição de médicos(as) a uma ampla gama de funções, como cirurgias, clínicas, tarefas administrativas e plantões, ao longo de um determinado horizonte de planejamento.
Para lidar com a abundante quantidade de tarefas a serem planejadas, os autores propuseram 4 modelos de programação linear inteira mista (MIP), sendo tais: \textit{modelo I} que visa minimizar tarefas não-alocadas, de maneira básica, desconsiderando pontos como preferência; \textit{modelos IIa e IIb}, que, respectivamente priorizam restrições flexíveis de preferência dos especialistas (chamado de agenda ideal) enquanto mantém o objetivo do primeiro modelo, e, procurar reduzir fadiga médica forçando restrições referentes a corredores de tempo de descanso e finais de semana (similar a definição de \citeonline{ERHARD20181}); enquanto o \textit{modelo III} procura um meio-termo entre as preferências dos modelos anteriores, por soluções ótimas de Pareto. Apesar de o uso de métodos heurísticos em instâncias maiores gerar soluções com desvios da otimalidade, os resultados eram melhores do que os gerados manualmente pela chefia de setor sem o auxílio de tomadores de decisão, concluindo a eficácia de MIPs na modelagem de problemas da área.

Os autores também observam que durante a execução do projeto, a restrição de disponibilidade de recursos estava diretamente atrelada a qualidade da solução, isto é, quando se limitava recursos nas restrições nos testes, aumentava-se o número de atividades não programadas. Apesar do que foi observado, a solução apresentada foi robusta o bastante para diversas instâncias.

\section{O Contexto Clínico}
\label{2rev-sec2}

Embora os esforços da comunidade cientifica, trabalhos como os de \citeonline{Beaulieu_Ferland_Gendron_Michelon_2000} e \citeonline{Topaloglu_2006}, assumem um contexto de sistema fechado, hospitalar, cujo intuito é garantir atendimento contínuo e ininterrupto ao paciente, respeitando questões interpessoais do profissional. Nesse contexto, um paciente já admitido pode ser atendido por vários profissionais durante sua estadia, remetendo a capacidade da estrutura hospitalar. Embora o estudo de \citeonline{Bruni_Detti_2014} seja generalizável para contextos da saúde, como o contexto ambulatorial (atendimentos rápidos e sem prévio agendamento), assume-se que o(a) especialista está associado(a) a uma única unidade (hospital). Nessa situação basta conhecer a demanda desta única facilidade.

Alguns autores, como \citeonline{ChenLinPeng2016} estendem as possibilidades do uso de modelagem matemática, aplicando otimização em dois estágios ao contexto hospitalar ambulatorial voltado ao funcionamento dos setores de exames clínicos, como radiologia e diagnóstico por imagem, na premissa de integrar os problemas de dimensionamento da equipe médica e alocação da equipe médica, enquanto \citeonline{BentayebLahrichiRousseau2023} estuda, de maneira similar, a integração de geração do quadro de consulta do paciente simultâneo à alocação de horários da equipe técnica em radiologia e exames por imagem para exames clínicos hospitalares. Em ambos os casos foram observadas melhorias significativas interpessoais da população alvo, como economia operacional devido à redução do dimensionamento da equipe necessária, além da aceleração no tempo de computação e confecção das escalas antecipadas. Somado a isso, a redução de violações de preferência, o trabalho de \citeonline{ChenLinPeng2016} apresenta desenvolvimento importante na literatura ao resolver o problema integrado, aumentando eficiência operacional e qualidade de vida, enquanto para \citeonline{BentayebLahrichiRousseau2023}, fora observada um aumento de 20\% no atendimento dos pacientes, reduzindo ociosidade de máquinas e encurtando custos operacionais correlatos. Os autores apontam, no entanto, possibilidades no futuro acadêmico de averiguar a robustez do sistema de dimensionamento de equipes, bem como a necessidade de generalização do estudo, que se limita ao contexto hospitalar médico. Também é amplamente discutido a hipótese de que a melhoria na eficiência de algoritmos também permitiria a expansão do estudo para toda a unidade, sem necessidade de segmentação, bem como a estabilidade da aplicação \cite{BentayebLahrichiRousseau2023}.


Para \citeonline{Edwards_Clague_Barlow_Clarke_Reed_Rada_1994}, foi necessário analisar a estrutura na qual os serviços se fundamentam de modo a identificar a origem dos problemas ambulatoriais de tempo de espera do paciente e fluxo de pacientes, métricas diretamente atreladas a satisfação do usuário final. Os autores averíguam por meio de simulações baseadas em dados reais os resultados obtidos pelo uso de sistemas de filas únicas - onde todos os pacientes são alocados a uma fila unitária - e filas quase-paralelas - onde são feitas triagem e criação de múltiplas filas de acordo com preferência e necessidade. Modelando tais sistemas como uma relação de agente (paciente) buscando por recurso limitado (médico(a) associado(a) a uma sala), foram observadas melhorias de 25\% no tempo de espera quando aplicadas filas quase-paralelas, devido sua eficiência e dinamismo, permitindo alterações e correções pontuais para o modelo, que por consequência aumentou a satisfação.
Em próprio texto íntegro, os autores comentam das lacunas no trabalho, especialmente na dificuldade de aplicação em contexto clínicos multidisciplinares (diferentes especializações) e a integração de um agenda de pacientes dinâmica, com reagendamentos, trocas de horários e faltas.

Outro trabalho que se dispôs similarmente a utilizar técnicas de simulação para alocar atendimentos em clínicas ambulatoriais foi \citeonline{Rau_Tsai_Liang_Tan_Syu_Jheng_Ciou_Jaw_2013}, exercido no contexto de sessões fisioterápicas, onde além das métricas de tempo de espera e satisfação, fora considerado a duração do exame por paciente, e consequentemente, a carga horária de trabalho do terapeuta. Com o uso de simulação de eventos discretos (do inglês, \textit{Discrete Event Simulation}, abreviada para DES), os autores obtiveram menores tempo de espera e um aumento nominal de 12 pacientes no fluxo diário, e especialmente à introdução de \textit{pooling} - na qual o paciente não é mais associado a um profissional preferencial, e pode ser atendido por outros profissionais da clínica - reduziu em 45\% os tempos de espera. Os autores \citeonline{Rau_Tsai_Liang_Tan_Syu_Jheng_Ciou_Jaw_2013} não consideram, no entanto, fatores estocásticos inerentes ao sistema, como falta de pacientes e desgaste dos equipamentos, limitando-se a estimativas baseadas em valores históricos.

Enquanto ambos os trabalhos citados não aplicam técnicas de otimização exata e modelo matemática para descriminar o problema observado, \citeonline{Alizadeh_Rezaeian_Abedi_Chiong_2020} aplicam MIP para construção de agendas de uso em exames não-emergenciais com necessidade de maquinário específico (referido na literatura como NOASP). Visando reduzir penalidades como: agendamentos em dias e horários indesejados, tempo de maquinário ocioso e trocas de médico(a), o modelo foca na preferência do paciente para criar agendas ótimas. Os autores apontam, no entanto, que o uso de métodos exatos se tornou inviável devido à natureza combinatorial do problema, quando escaladas as instâncias, levando-os a criar um algoritmo genético (abreviado para GA), que, embora não apresentasse soluções comprovadamente ótimas, se gerava agendamentos de boa qualidade, 10\% do desvio relativo.

\subsection{O Problema Multiclínicas}
\label{2rev-sec2-sub1}

Evidencia-se na Seção \ref{2rev-sec2} que  a maioria das investigações discorridas tem como principal motivação o estudo do agendamento em condições ambulatoriais (referido no inglês como \textit{outpatient}). Como não existem agendamentos prévios para esse sistema, trabalhos como o de \citeonline{ChenLinPeng2016} focam em equilibrar fatores estocásticos como incerteza de demanda por meio de otimização robusta, enquanto \citeonline{Alizadeh_Rezaeian_Abedi_Chiong_2020} consideram o(a) médico(a) como recurso a ser disponibilizado ao agente (paciente) para modelagem do sistema matemático; já os demais autores utilizam de artifícios de simulação com o(a) médico(a) fixo(a) para solução do problema. Os autores mencionados restringem-se a uma ótica unitária, onde há somente uma clínica, que, portanto, não é parte do modelo, mas sim delimitador do conjunto. Ao não considerar a possibilidade múltiplas unidades como parte do problema, abre-se mão da possibilidade de estudar o comportamento de modelos em situação onde um profissional pode ser alocado em diferentes ambientes, bem como suas implicações ergonômicas e legais.

O trabalho desenvolvido por \citeonline{Lai_Lu_Chen_2024} aborda o conceito de múltiplas clínicas em um hospital através do uso de \textit{frameworks} voltados a métodos de mineração de dados, como redes neurais artificiais e árvores de decisão para determinar mediante classificação os profissionais mais adequados a um grupo de pacientes. Como principal proposta do estudo, investigou-se a natureza de consultas que carecem de múltiplos segmentos, como encaminhamento para exames de ressonância magnética, os quais também compartilham recursos de outros setores do hospital, podem ser modeladas como grupos, que podem ser designados a combinação de equipes para melhor eficácia nos atendimentos.

\section{Heurísticas e Meta-Heurísticas em Escalonamento na Saúde}
\label{2rev-sec3}

Embora os métodos exatos de Programação Linear Inteira Mista discutidos nas seções anteriores ofereçam garantias de otimalidade, sua aplicabilidade é limitada quando a instância cresce em número de profissionais, períodos ou unidades, dada a natureza NP-difícil dos problemas de escalonamento \cite{Even_Itai_Shamir_1976}. Por essa razão, parcela significativa da literatura recente em alocação de médicos(as) e escalonamento de enfermeiras(os) recorre a heurísticas, meta-heurísticas e mat-heurísticas, abordagens que abrem mão da prova de otimalidade em troca de tempos computacionais viáveis em escala operacional, tendência também apontada nas revisões de \citeonline{Abdalkareem2021_ac} e \citeonline{Ngoo_Goh_Sze_Sabar_Abdullah_Kendall_2022}. Esta seção apresenta trabalhos primários representativos dessa linha, organizados em métodos heurísticos e mat-heurísticos, e na subseção subsequente em métodos evolutivos baseados em NSGA-II.

\citeonline{Vermuyten_Namorado_Rosa_Marques_Belien_Barbosa_Povoa_2018} abordam o escalonamento integrado de equipes em um serviço médico de emergência, considerando simultaneamente médicos(as), enfermeiros(as) e despachadores ao longo do horizonte de planejamento. A metodologia combina uma formulação em Programação Inteira com duas heurísticas baseadas em decomposição: \textit{diving} e Variable Neighborhood Decomposition Search (VNDS), aplicadas após relaxação parcial do modelo. Avaliada sobre instâncias do Instituto Nacional de Emergência Médica de Portugal, a heurística VNDS supera tanto o método \textit{diving} quanto um \textit{solver} comercial em qualidade de solução para horizontes realistas, e foi posteriormente integrada em sistema de apoio à decisão com interface gráfica para uso operacional. O trabalho evidencia a viabilidade de heurísticas estruturalmente alinhadas ao MIP em contextos onde a otimização exata se torna computacionalmente proibitiva.

\citeonline{Turhan_Bilgen_2020} propõem uma mat-heurística híbrida para o NRP, combinando os métodos Fix-and-Relax (F\&R), Fix-and-Optimize (F\&O) e Simulated Annealing (SA). O F\&R decompõe o horizonte em janelas semanais resolvidas como subproblemas MIP até a otimalidade, fornecendo soluções iniciais ao SA, que executa movimentos de vizinhança sobre escalas completas; o F\&O é invocado periodicamente para otimizar novamente subconjuntos de variáveis quando o SA estagna, alternando exploração global e refinamento exato local. Em testes sobre 24 instâncias \textit{benchmark} públicas, a abordagem híbrida supera os métodos do estado da arte na maioria dos casos e estabelece sete novos melhores resultados conhecidos. Os autores ilustram um arquétipo de mat-heurística no qual o MIP cuida da estrutura local enquanto o componente meta-heurístico responde pela diversificação, padrão metodológico diretamente relacionado à abordagem adotada neste estudo.

\citeonline{Ceschia_Di_Gaspero_Mazzaracchio_Policante_Schaerf_2023} introduzem uma formulação de NRP em hospitais italianos, considerando habilidades distintas, turnos especializados, limites de carga horária, tipos diversos de folgas, preferências individuais e incompatibilidades interpessoais. A solução combina um modelo MIP de referência com uma busca local guiada por SA que alterna dois movimentos: realocação enfermeiro-dia e troca de padrões de turno. Avaliada sobre 34 instâncias reais do norte da Itália, a metodologia produziu escalas aceitas pelas instituições e é distribuída como parte de \textit{software} comercial de \textit{rostering}. O trabalho exemplifica como uma meta-heurística orientada à prática absorve restrições idiossincráticas mantendo tempos de computação tratáveis para uso operacional cotidiano.

\citeonline{Santos_Toffolo_Gomes_Ribas_2016} apresentam técnicas de Programação Inteira para o NRP no contexto da \textit{International Nurse Rostering Competition} (INRC), combinando uma formulação compacta com geração de cortes e heurísticas primais baseadas em decomposição temporal. O método fixa as variáveis fora de uma janela deslizante de dias e reotimiza apenas o subproblema interno via MIP, padrão posteriormente reconhecido pela literatura como mat-heurística do tipo Fix-and-Optimize. A abordagem provou otimalidade na maioria das instâncias da competição, melhorou em até 15\% as melhores soluções conhecidas e venceu a INRC-2010 em todas as categorias. O trabalho consolidou a janela deslizante de F\&O como referência em escalonamento de enfermeiras(os).

\citeonline{Devesse_Santos_Toledo_2017} aplicam mat-heurísticas Relax-and-Fix e Fix-and-Optimize a um problema de escalonamento de médicos(as) em prontos-socorros, com formulação MIP \textit{minimax} orientada à equidade na distribuição de turnos indesejáveis sob restrições da legislação trabalhista brasileira. As heurísticas particionam as variáveis binárias por semanas e por tipos de turno, resolvendo cada subproblema com CPLEX e aceitando melhorias incrementais sobre a escala incumbente. Avaliada em instâncias com 10 a 40 médicos(as) e horizontes de 4 a 36 semanas, a abordagem mostra que R\&F e F\&O produzem soluções competitivas com o MIP em uma fração do tempo, embora o MIP completo ainda dominasse instâncias menores. O trabalho é particularmente relevante por aplicar diretamente o par R\&F + F\&O ao escalonamento médico, contexto idêntico ao tratado no Capítulo \ref{cap:mip-heuristicas}.

\citeonline{Wickert_Kummer_Neto_Boniatti_Buriol_2021} propõem um modelo de Programação Inteira e uma mat-heurística Fix-and-Optimize para o escalonamento de médicos(as) no Hospital de Clínicas de Porto Alegre, considerando preferências individuais, restrições legais, requisitos hospitalares e equilíbrio de carga de trabalho. O método constrói uma solução factível via MIP e iterativamente reotimiza vizinhanças definidas por dia, médico(a), semana ou turno, mantendo as demais variáveis fixadas. Em instâncias sintéticas com 50, 100 e 150 médicos(as) ao longo de 4 semanas, a mat-heurística F\&O sistematicamente supera o MIP em horizontes de 10 a 30 minutos, particularmente nas instâncias maiores em que o solver enfrenta dificuldade de convergência. O estudo documenta um precedente direto, em contexto brasileiro, do uso de F\&O para escalonamento de médicos(as) em larga escala.

Os trabalhos apresentam padrão recorrente, evidenciado em escalonamento de enfermeiras(os) e médicos(as): a estrutura combinatorial do MIP é preservada em subproblemas locais, enquanto componentes meta-heurísticos respondem pela diversificação no espaço de busca. Essa hibridização é particularmente útil em horizontes longos e em contextos com múltiplas restrições conflitantes, condições que também caracterizam o problema vigente, motivando a adoção de métodos análogos no Capítulo \ref{cap:mip-heuristicas}.

\subsection{Abordagens Multiobjetivo em Escalonamento}
\label{2rev-sec3-sub1}

Problemas de escalonamento na saúde combinam objetivos tipicamente conflitantes, como cobertura da demanda, equidade na distribuição de turnos e atendimento a preferências individuais, motivando o tratamento explícito da estrutura multiobjetivo da formulação. A literatura organiza-se em duas grandes famílias metodológicas: o escalonamento, que agrega múltiplos critérios em uma única função objetivo via custos ponderados ou programação por metas, e a busca direta de soluções Pareto-eficientes, suportada tanto por métodos exatos de decomposição quanto por meta-heurísticas baseadas em dominância. As subseções a seguir percorrem trabalhos representativos dessas duas famílias, culminando no NSGA-II, algoritmo evolutivo adotado como base no Capítulo \ref{cap:mip-heuristicas}.

\citeonline{Parr_Thompson_2007} estudam o NRP com função de custo ponderada, tendo a cobertura como restrição rígida e os demais critérios (preferências, padrões de turno, equidade) agregados em penalidades soft. A metodologia compara duas estratégias de busca local sobre a função ponderada: SAWing, que ajusta dinamicamente os pesos durante o \textit{simulated annealing}, e Noising-SA, que perturba os custos com ruído para escapar de ótimos locais. O trabalho demonstra que Noising-SA produz escalas com penalidade total consistentemente menor, sendo menos sensível à calibração de pesos, evidenciando uma das principais dificuldades práticas do escalonamento.

\citeonline{Maenhout_Vanhoucke_2010} formulam o NRP de médio prazo com múltiplos objetivos e prioridades distintas, modelando-o como um problema de \textit{set partitioning} resolvido por \textit{branch-and-price} exato. A contribuição central está nas estratégias de \textit{branching} e \textit{pruning} adaptadas ao contexto multiobjetivo, capazes de preservar a eficiência da geração de colunas enquanto direcionam a busca para escalas de alta qualidade. Os experimentos computacionais demonstram que regras de \textit{branching} bem projetadas viabilizam a otimização exata de instâncias realistas, oferecendo um contraponto exato à dominância de meta-heurísticas no escalonamento multiobjetivo.

\citeonline{Burke_Li_Qu_2012} propõem busca Pareto-direta para o NRP em uma UTI holandesa, com formulação que possui nove restrições suaves distintas. O método combina geração aleatória de padrões legais por enfermeiro, heurística adaptativa para alocação inicial, reparo de cobertura e \textit{simulated annealing} multiobjetivo com dois modos de avaliação: soma ponderada e dominância de Pareto. A variante baseada em dominância produz dezenas a milhares de escalas não dominadas em cada instância, em cenário onde a Programação Inteira completa não encontrou solução factível, ilustrando a vantagem da busca Pareto-direta em formulações fortemente restritas.

A linha evolutiva multiobjetivo passou a dominar a literatura subsequente em escalonamento na saúde a partir do NSGA-II, formalizado por \citeonline{Deb_Pratap_Agarwal_Meyarivan_2002}, cuja contribuição metodológica é detalhada a seguir.

O algoritmo NSGA-II foi proposto por \citeonline{Deb_Pratap_Agarwal_Meyarivan_2002} para superar três limitações dos algoritmos evolutivos multiobjetivo anteriores: complexidade computacional elevada, ausência de elitismo e dependência do parâmetro de \textit{sharing}. A metodologia introduz a ordenação não-dominada rápida com complexidade \(O(MN^{2})\), a distância de \textit{crowding} para preservação de diversidade, e a seleção elitista baseada na união entre populações pai e filha, além de uma regra de dominância modificada para problemas com restrições. Em \textit{benchmarks} de teste, o NSGA-II superou os algoritmos SPEA e PAES em convergência e em distribuição da fronteira aproximada de Pareto. O trabalho consolidou o NSGA-II como o algoritmo evolutivo multiobjetivo de referência em Pesquisa Operacional, sendo adotado ou utilizado como linha de base em grande parte da literatura subsequente em escalonamento.

\citeonline{Hamid_Tavakkoli_Moghaddam_Golpaygani_Vahedi_Nouri_2020} formulam um problema multiobjetivo de escalonamento de enfermeiras(os) que enfatiza fatores humanos, com três objetivos: minimização do custo total de pessoal, minimização da incompatibilidade entre estilos de tomada de decisão na mesma escala e maximização da satisfação com os turnos atribuídos. Para resolução, comparam três meta-heurísticas multiobjetivo - NSGA-II, algoritmo Keshtel multiobjetivo e Busca Tabu multiobjetivo - e utilizam Análise Envoltória de Dados (DEA) para classificar as soluções da fronteira de Pareto. Em estudo de caso em hospital iraniano, todos os algoritmos geram escalas factíveis com diferentes compromissos entre os critérios, sendo o NSGA-II competitivo em termos de qualidade da fronteira. Os autores evidenciam que tratar fatores humanos como objetivos explícitos altera estruturalmente o conjunto Pareto-eficiente, com pequenos aumentos de custo gerando ganhos significativos em compatibilidade e satisfação.

\citeonline{Lavygina_Welsh_Crispin_2019} aplicam o NSGA-II ao escalonamento de médicos(as) residentes no Reino Unido, sob regras nacionais de horas trabalhadas e descanso. A novidade do estudo é tratar a equidade como segundo objetivo, formalizada como o desvio padrão das penalidades acumuladas por médico(a), de modo que soluções equitativas distribuem características indesejáveis da escala uniformemente entre os profissionais. Comparado contra variantes de Pareto Simulated Annealing e algoritmos genéticos mono-objetivo em dez instâncias, o NSGA-II produziu fronteiras com hipervolume consistentemente superior e convergência mais rápida. A aplicação direta a médicos(as) e a inclusão da equidade como objetivo explícito o aproximam metodologicamente da estrutura de objetivos conflitantes adotada no problema vigente.

\citeonline{Khalili_Shahrezaei_Abri_2020} propõem um modelo bi-objetivo de escalonamento de enfermeiras(os) que incorpora o fator fadiga, modelado em função dos intervalos de descanso e dos padrões de turnos. Para a maior instância do estudo de caso (90 profissionais no hospital iraniano Labbafinejad), o modelo é resolvido pelo NSGA-II e por MOPSO, com cromossomos especificamente projetados, e o método TOPSIS combinado a ponderação por entropia é empregado para selecionar a solução final na fronteira de Pareto. O NSGA-II encontrou escalas de melhor qualidade em menor tempo computacional, com soluções que indicam padrões de descanso noturno e intermediário capazes de reduzir fadiga sem comprometer significativamente a estrutura de custos. O trabalho ilustra a combinação NSGA-II com técnicas de seleção a posteriori, padrão útil quando o tomador de decisão precisa de uma única recomendação a partir do conjunto de Pareto.

O padrão recorrente nas aplicações evolutivas multiobjetivo aqui descritas envolve representação direta das variáveis de alocação no cromossomo, operadores de cruzamento e mutação adaptados ao contexto e, frequentemente, a integração com técnicas auxiliares de seleção ou refinamento local. Esses elementos motivam diretamente as variantes propostas no Capítulo \ref{cap:mip-heuristicas}, nas quais o NSGA-II base é estendido com operadores customizados e adaptações estruturais para o problema de alocação de médicos(as) em múltiplas clínicas.

\section{Lacunas de Pesquisa}

Não é comum observar na literatura acadêmica investigações acerca da natureza de múltiplas unidades de um mesmo segmento clínico cujas operações sejam completamente autônomas entre si, como filiais de redes laboratoriais de exames. Nesse contexto, a unidade tem recursos (salas, equipamentos, médicos), disponibilidades e demanda intrínsecos a si mesma, e a relação de \textit{vários-para-vários} (como exemplificado na imagem da Figura \ref{fig:rel-n-to-n}) entre profissional e unidade causa uma explosão combinatória nas possibilidades de escalonamento, uma vez que existem múltiplas possibilidades de alocações possíveis, o que inevitavelmente gera situações de remanejamento, onde um profissional outrora alocado em uma unidade é escalado para outra. Além disso, as demandas e disponibilidades, bem como fluxo operacional, podem não ser iguais, dificultando a modelagem quando aplicada a um panorama de encontrar boas soluções para todas as clínicas de uma só vez.

Não é considerado também o contexto de especialistas associados em regime de trabalho de pessoa jurídica. Nessa situação as normas ergonômicas de trabalho e questões salariais são completamente diferentes do normalmente observado na literatura, assim como não foram constatados estudos acerca do remanejamento dos profissionais em relação a diversas unidades ao longo do horizonte de planejamento e suas implicações quanto a qualidade da solução.

\newpage

\begin{figure}[h]

	\begin{subfigure}{0.5\textwidth}
		\includegraphics[scale=0.1]{images/clinic-one-to-m.png}
		\caption{Relação vários para um}
		\label{fig:rel-n-to-1}
	\end{subfigure}
	\begin{subfigure}{0.5\textwidth}
		\includegraphics[scale=0.1]{images/clinic-many-to-many.png}
		\caption{Relação vários para vários}
		\label{fig:rel-n-to-n}
	\end{subfigure}

	\caption{Exemplificação da diferença entre os contextos de relacionamento apresentados.}
	\label{fig:exemplo-relacionamentos}
	\fautor
\end{figure}

Para este estudo de mestrado considera-se como problema a alocação de médicos de um único segmento clínico em múltiplas unidades, considerando que tanto o profissional com sua agenda de disponibilidade para atendimentos, e a clínica com sua demanda de exames a cumprir devem ser respeitadas, não sendo possível, portanto realizar a alocação em períodos onde a unidade ou o profissional não estão disponíveis. Aplica-se também a situação na qual a folga de um profissional deve ser coberta. Além disso, propõe-se a investigação acerca da possibilidade de redução nas trocas de clínicas trabalhadas por um profissional ao longo do horizonte de planejamento e suas implicações quanto ao atendimento de demanda.

Ao avaliar o sistema na totalidade, buscando atender critérios operacionais de maneira global, para todas as unidades simultaneamente, têm-se um sistema complexo de alta sensibilidade, onde a alocação de um profissional (agente) em uma clínica pode afetar outra, assemelhando-se a definição de escalonamento (do inglês \textit{rostering}) proposta por \cite{Wren_1996}.

O Quadro \ref{tab:taxonomia-estudos} de avaliação taxonômica ressalta de maneira sucinta aspectos-chave dos principais trabalhos a serem considerados na arguição deste estudo quanto às lacunas de pesquisa e descreve de maneira sucinta as contribuições deste estudo. Nele, o campo \textit{Escalas Acíclicas} indica se o trabalho considera escalas em que os profissionais não seguem um padrão fixo e repetitivo de turnos ao longo do horizonte de planejamento — característica presente no contexto estudado, onde cada médico(a) possui disponibilidades individuais e variáveis entre os dias.


\begin{quadro}[h]
	\centering
	\resizebox{\textwidth}{!}{%
		\begin{footnotesize}
			\begin{tabular}{|M{2cm}|M{2cm}|M{1.5cm}|M{1.5cm}|M{1.5cm}|M{1.5cm}|M{1.8cm}|M{3cm}| }
				\hline
				\textbf{Estudo}                                                        & \textbf{Ambientação} & \textbf{Múltiplas Unidades}    & \textbf{Alocação de Médicos} &
				\textbf{Escalas Acíclicas}                                             & \textbf{Dados Reais} & \textbf{Minimização de Trocas} & \textbf{Metodologia}                                                                                                               \\
				\hline
				\citeonline{Edwards_Clague_Barlow_Clarke_Reed_Rada_1994}               & Ambulatório          &                                &                              &            & \textbf{X} &            & Simulação                                                    \\
				\hline
				\citeonline{BentayebLahrichiRousseau2023}                              & \textbf{Exames}  &                                & \textbf{X}                   & \textbf{X} &            &            & \textbf{Modelagem Matemática}                                \\
				\hline
				\citeonline{ChenLinPeng2016}                                           & Ambulatório          &                                & \textbf{X}                   &            &            &            & \textbf{Modelagem Matemática}                                \\
				\hline
				\citeonline{Bruni_Detti_2014}                                          & Hospitalar           &                                & \textbf{X}                   &            &            &            & \textbf{Modelagem Matemática}                                \\
				\hline
				\citeonline{Gunawan_Lau_2013}                                          & Hospitalar           &                                & \textbf{X}                   &            &            &            & Heurística                                                   \\
				\hline
				\citeonline{Topaloglu_2006}                                            & Emergencial          &                                &                              & \textbf{X} & \textbf{X} &            & \textbf{Modelagem Matemática}                                \\
				\hline
				\citeonline{Alizadeh_Rezaeian_Abedi_Chiong_2020}                       & Ambulatório          &                                & \textbf{X}                   &            &            &            & Heurística                                                   \\
				\hline
				\citeonline{Rau_Tsai_Liang_Tan_Syu_Jheng_Ciou_Jaw_2013}                & Ambulatório          &                                &                              &            & \textbf{X} &            & Simulação                                                    \\
				\hline
				\citeonline{Lai_Lu_Chen_2024}                                          & Hospitalar           & \textbf{X}                     &                              &            &            &            & Mineração de Dados                                           \\
				\hline
				\citeonline{Vermuyten_Namorado_Rosa_Marques_Belien_Barbosa_Povoa_2018} & Emergencial          &                                & \textbf{X}                   &            & \textbf{X} &            & Heurística                                                   \\
				\hline
				\citeonline{Turhan_Bilgen_2020}                                        & Hospitalar           &                                &                              &            &            &            & mat-heurística                                                \\
				\hline
				\citeonline{Lavygina_Welsh_Crispin_2019}                               & Hospitalar           &                                & \textbf{X}                   &            & \textbf{X} &            & Meta-heurística                                               \\
				\hline
				\citeonline{Hamid_Tavakkoli_Moghaddam_Golpaygani_Vahedi_Nouri_2020}    & Hospitalar           &                                &                              &            & \textbf{X} &            & Meta-heurística                                               \\
				\hline
				\citeonline{Khalili_Shahrezaei_Abri_2020}                              & Hospitalar           &                                &                              &            & \textbf{X} &            & Meta-heurística                                               \\
				\hline
				\citeonline{Ceschia_Di_Gaspero_Mazzaracchio_Policante_Schaerf_2023}    & Hospitalar           &                                &                              &            & \textbf{X} &            & Meta-heurística                                               \\
				\hline
				\citeonline{Parr_Thompson_2007}                                        & Hospitalar           &                                &                              &            &            &            & Meta-heurística                                               \\
				\hline
				\citeonline{Maenhout_Vanhoucke_2010}                                   & Hospitalar           &                                &                              &            &            &            & \textbf{Modelagem Matemática}                                \\
				\hline
				\citeonline{Burke_Li_Qu_2012}                                          & Hospitalar           &                                &                              &            & \textbf{X} &            & Meta-heurística                                               \\
				\hline
				\citeonline{Santos_Toffolo_Gomes_Ribas_2016}                           & Hospitalar           &                                &                              &            &            &            & mat-heurística                                                \\
				\hline
				\citeonline{Devesse_Santos_Toledo_2017}                                & Emergencial          &                                & \textbf{X}                   &            &            &            & mat-heurística                                                \\
				\hline
				\citeonline{Wickert_Kummer_Neto_Boniatti_Buriol_2021}                  & Hospitalar           &                                & \textbf{X}                   &            & \textbf{X} &            & mat-heurística                                                \\
				\hline
				\textbf{Este estudo}                                                   & \textbf{Exames}   & \textbf{X}                     & \textbf{X}                   & \textbf{X} & \textbf{X} & \textbf{X} & \textbf{Modelagem Matemática, Heurísticas e Meta-Heurísticas} \\
				\hline
			\end{tabular}
		\end{footnotesize}}
	\caption{Análise taxonômica de estudos em Pesquisa Operacional relacionados à Alocação e Escalonamento na saúde.}
	\fautor
	\label{tab:taxonomia-estudos}
\end{quadro}

\chapter{Definição do Problema e Modelagem Matemática} % 3
\label{cap:problema}

Neste capítulo, apresenta-se a definição do problema em estudo, a partir de uma contextualização detalhada do cenário real, seguida da modelagem matemática do problema, contemplando as principais restrições e objetivos identificados. A modelagem é estruturada de maneira a refletir as particularidades operacionais e administrativas da rede de clínicas, garantindo que as soluções geradas sejam aplicáveis e relevantes para o contexto em questão. A seguir, discorre-se sobre cada etapa dessa formalização, iniciando pela contextualização do problema e detalhando posteriormente a modelagem matemática, incluindo as funções objetivo e as restrições associadas.

\section{Contextualização do Problema}
\label{sec:contexto-problema}
O problema considerado neste estudo consiste em alocar médicos profissionais da área de ultrassonografia e diagnóstico por imagem em uma rede de clínicas baseado em um problema real de uma grande cidade do Brasil com o intuito de atender a demanda de exames solicitada para o período conforme a disponibilidade de atendimentos por cada profissional. Tanto disponibilidade quanto demanda são valores discretos, que, quando não conhecidos previamente, são estimados a partir do tempo médio de uma consulta, com tempo médio de 15 minutos para esse tipo de exame, conforme aponta o \citeonline{conselho_regional_de_medicina_do_estado_do_parana_2022}.

O contrato de trabalho dos profissionais permite livre acordo, para todos os efeitos, com o corpo gerencial das unidades, garantindo flexibilidade, pois o profissional pode negociar os dias da semana e os turnos em que atuará. A empresa, por sua vez, contrata o profissional com base na previsão de exames em suas unidades e disponibilidade de  salas.
Não existem restrições administrativas referentes a remanejar um profissional entre unidades, desde que o mesmo esteja associado contratualmente a nova clínica, e, essa troca não ocorra entre turnos de um mesmo dia (ver Figura \ref{fig:ex-sched-medico}).

\begin{figure}[h]
	\centering
	\includegraphics[width=10cm]{images/exemplo-schedule-medico.png}
	\caption{Exemplo de alocação para um médico qualquer.}
	\label{fig:ex-sched-medico}
	\fautor
\end{figure}

A empresa atua com horizontes de planejamento de diferentes tamanhos, com a finalidade de gerar escalas de trabalho (horizontes de até 90 dias), embora também apliquem de maneira flexibilizada outros horizontes menores, como 60 e 30 dias, comum para pequenos ajustes de escala. Tais dados de demanda e atendimento são obtidos a partir de registros históricos de sistemas de gestão hospitalar utilizados pela rede de clínicas, permitindo observar a sazonalidade. Juntamente com o desenvolvimento de métodos de aprendizado de máquina para regressão e captação de agendamentos ativos, isso permite à empresa gerar uma estimativa robusta de exames, embora não seja o objetivo deste trabalho.

Os demais dados e registros referentes a atuação de profissionais clínicos associados à rede são mantidos em sistemas internos de gestão, permitindo a verificação de disponibilidade dos profissionais, que também informam suas disponibilidades com antecedência ao horizonte de planejamento. O controle de salas e unidades associadas a cada médico também é incluso no sistema interno.

Isso permite que a empresa possua um histórico, possibilitando a geração de previsões confiáveis para o horizonte de planejamento vigente. Portanto, os registros de demanda e disponibilidade empregados para a elaboração deste referido trabalho foram previamente obtidos em comum acordo com a empresa, tratando-se de dados estáticos, imutáveis quanto ao escopo e desenvolvimento deste mestrado.

O processo de alocação dos médicos é feito manualmente pelo setor  de operações, buscando respeitar normas contratuais e operacionais, associado um profissional a um horário, utilizando como principais métricas de qualidade da solução a quantidade total de exames não atendidos e a cobertura das clínicas, que permitem visualizar e implementar possíveis ajuste. Porém, como essa tarefa é realizada manualmente, demanda tempo do setor e gera soluções sem garantia de otimalidade.

Na Tabela \ref{tab:hipoteses-modelagem} apresentam-se as hipóteses levantadas para o desenvolvimento da modelagem a partir das regras de funcionamento e conformidades administrativas utilizadas pela gestão da rede. Considera-se para este trabalho, portanto, um problema de escalonamento no qual várias unidades com demandas e disponibilidades de equipamentos diferentes devem ser atendidas simultaneamente.

% usar tabela restrições
\begin{table}[h]
	\centering
	\caption{Hipóteses levantadas referentes à natureza do problema.}
	\begin{tabular}{cM{12cm}}
		\toprule
		\textbf{Hipótese} & \textbf{Descrição}                                                                                                                \\
		\midrule
		\textit{H1}       & Todas as clínicas devem ser atendidas, mesmo que com cobertura mínima                                                             \\
		\textit{H2}       & Um médico deve ser associado a somente uma unidade  por dia                                                                       \\
		\textit{H3}       & Cada médico deve estar associado a uma sala da clínica                                                                            \\
		\textit{H4}       & Um médico realiza uma determinada quantidade de exames da demanda da unidade.                                                     \\
		\textit{H5}       & A diferença entre atendimentos e demanda é denominada de exames não atendidos                                                     \\
		\textit{H6}       & Dada a agenda irregular dos médicos, não existe limite para total de remanejamentos/trocas entre clínicas                         \\
		\textit{H7}       & Assumindo uma condição inicial de alocação precedente ao horizonte de planejamento vigente (histórico), minimizam-se possíveis remanejamentos \\
		\bottomrule
	\end{tabular}
	\label{tab:hipoteses-modelagem}
	\fautor
\end{table}

Para tal atendimento, portanto, é necessário alocar um médico a uma sala, sem sobreposição ou choque profissionais. Ao ser alocado a uma unidade no dia, o profissional não pode ser remanejado para outra clínica, mesmo que tenha disponibilidade para trabalhar no contraturno da alocação. Essas escalas devem ser elaboradas de modo a minimizar o número de exames não atendidos, bem como evitar a transferência de profissionais entre unidades, conforme discutido nas seções a seguir.

\section{Modelagem Matemática}

O problema é modelado como uma otimização inteira, em que cada médico pode ser alocado a no máximo, uma unidade por turno, sem a possibilidade de trabalhar em mais de uma unidade no mesmo período, respeitando a disponibilidade do mesmo.  Define-se por período o  \(\mbox{dia}\) da semana, sendo que esse é dividido em  \(\mbox{turnos}\). Cada dia da semana possui dois turnos (manhã e tarde), incluindo o domingo, que, para fins de modelagem, mesmo que não possua demanda vespertina, ainda é incorporado no modelo com demandas e disponibilidades \(0\). Dessa forma é possível sempre representar o período a partir da tupla \(d \times t\). Além disso, a quantidade de profissionais alocados não pode exceder a capacidade da unidade, definida pelo número de salas disponíveis para exames. Por fim, sempre que alguma unidade tiver demanda em um período e turno, impreterivelmente, ao menos um médico deverá ser alocado a ela para garantir cobertura mínima da unidade.

\newpage

\begin{small}
	\begin{tabular}{ll}
		\multicolumn{2}{l}{\textbf{Índices, conjuntos e subconjuntos}} \\
		\(i \in I\)       \quad & Médicos                                \\
		\(u \in U\)       \quad & Unidades                               \\
		\(U_i\)           \quad & Unidades atendidas pelo médico \(i\)     \\
		\(I_u\)           \quad & Médicos que atendem na unidade \(u\)     \\
		\(d \in D\)       \quad & Dias                                   \\
		\(t \in T\)       \quad & Turnos                                 \\
	\end{tabular}
\end{small}



\begin{small}
	\begin{tabular}{ll}
		\multicolumn{2}{l}{\textbf{Parâmetros}}                                                                                    \\
		\(e_{udt}\)       \quad & Demanda por exames na unidade \(u\) no dia \(d\) e turno \(t\)                                           \\
		\(s_{ut}\)       \quad  & Salas disponíveis para exames na unidade \(u\) no dia \(d\) e turno \(t\)                                \\
		\(q_{idt}\)      \quad  & Quantidades de exames que o médico \(i\) pode realizar dia \(d\) e turno \(t\)                           \\
		\(a_{idt}\)      \quad  & \(\begin{cases} 1, & \text{se o médico $i$ estiver disponível para trabalho no dia $d$ e turno $t$} \\
              0, & \text{caso contrário.}\end{cases} \) \\
		\(b_{udt}\)      \quad  & \(\begin{cases} 1, & \text{se a unidade $u$ tem demanda no dia $d$ e turno $t$} \\
              0, & \text{caso contrário.}\end{cases} \)
	\end{tabular}
\end{small}

O parâmetro \(q_{idt}\) é estimado com base na quantidade de horas do turno e no fato de que cada exame, gasta-se um tempo médio de 15 minutos, conforme levantando \citeonline{conselho_regional_de_medicina_do_estado_do_parana_2022}.
Salienta-se que parâmetros auxiliares como \(b_{udt}\) ou \(a_{idt}\), que indicam a existência de demanda ou disponibilidade, respectivamente, tem por finalidade facilitar a linearização do modelo, uma vez que suas informações são diretamente baseadas em outros parâmetros, como \(e_{udt}\) e \(q_{idt}\), respectivamente. Dessa forma, sabemos que caso o valor de \(e_{udt}\) seja maior que \(0\), necessariamente, o valor de \(b_{udt}\) será \(1\), indicando a existência de demanda. De maneira similar, caso o valor de \(q_{idt}\) seja maior que \(0\), o valor de \(a_{idt}\) será \(1\), indicando a disponibilidade do profissional para atendimento.

\begin{small}
	\begin{tabular}{ll}
		\multicolumn{2}{l}{\textbf{Variáveis}}                                                                                                                   \\
		\(x_{iudt}\)       \quad & \( \begin{cases} 1, & \text{se o médico $i$ está alocado na unidade $u$ no dia $d$ e turno $t$} \\    0, & \text{caso contrário.}
			                           \end{cases} \) \\
		\(w_{iud}\)       \quad  & \( \begin{cases} 1, & \text{se o médico $i$ está alocado na unidade $u$ durante o dia $d$} \\    0, & \text{caso contrário.}
			                           \end{cases} \)      \\
		\(z_{id}\)       \quad   & \( \begin{cases} 1, & \text{se o médico $i$ trocou de unidade no dia $d$} \\    0, & \text{caso contrário.} \end{cases} \)        \\
	\end{tabular}
\end{small}

Discorre-se neste trabalho a modelagem do problema observado de maneira genérica, sem a inclusão de objetivos particulares, com a finalidade de apresentar a estrutura básica do problema e validar as hipóteses levantadas no quadro \ref{tab:hipoteses-modelagem}. A seguir, apresenta-se a modelagem de restrições e variáveis:

\begin{alignat}{2}
	\sum_{u \in U_i} x_{iudt} \leq a_{idt}  \quad & \forall i \in I, d \in D, t \in T \label{eq:r1}            \\
	\sum_{u \in U_i} w_{iud} \leq 1  \quad        & \forall i \in I, d \in D \label{eq:r2}                     \\
	x_{iudt} \leq w_{iud} \quad                   & \forall i \in I, u \in U, d \in D, t \in T_d \label{eq:r3} \\
	w_{iud} = 0 \quad                             & \forall i \in I, u \notin U_i, d \in D \label{eq:r4}       \\
	\sum_{i \in I_u} x_{iudt} \leq s_{udt} \quad  & \forall u \in U, d \in D, t \in T \label{eq:r5}            \\
	\sum_{i \in I_u}x_{iudt} \geq b_{udt} \quad   & \forall u \in U, d \in D, t \in T \label{eq:r6}            \\
	x_{iudt}, w_{iud} \in \{0,1\} \quad           & \forall i \in I, u \in U, d \in D, t \in T \label{eq:dom2} \\
\end{alignat}

O conjunto de restrições (\ref{eq:r1}) assegura que o médico só será alocado para uma unidade qualquer em um determinado dia e turno, se e somente se o mesmo houver disponibilidade para atendimento nesse período. As restrições (\ref{eq:r2}-\ref{eq:r3}) asseguram que o profissional será alocado apenas a uma unidade por dia, independente de quantos turnos trabalhe, enquanto \ref{eq:r4} assegura que o mesmo não poderá trabalhar em clínicas as quais não está associado. A quantidade de médicos alocados por turno é limitada pela quantidade de salas disponibilizadas pela unidade, mediante restrições (\ref{eq:r5}). Em contraparte, as restrições (\ref{eq:r6}) indicam que, se houver demanda por trabalho na unidade, ao menos um profissional deve ser alocado para garantir cobertura mínima da clínica. As restrições (\ref{eq:dom2}) definem o domínio das variáveis.

\subsection{Funções Objetivo}
Mediante a necessidade de atender os problemas observados pela rede de clínicas, a investigação das hipóteses levantadas na Tabela \ref{tab:hipoteses-modelagem} proporcionaram para o estudo objetivos distintos a serem contemplados. A nível operacional, têm-se a cobertura dos exames clínicos, tratando da minimização de exames não atendidos, e consequentemente o fluxo de operações e caixa dessa modalidade. No nível ergonômico, as realocações de médicos ao longo do horizonte de planejamento têm por objetivo minimizar o número de trocas de clínica, evitando o excesso de remanejamentos nas agendas médicas que reflete na rotatividade de funcionários e relação com o paciente. Ambos os casos têm implicações diretas na modelagem e solução do problema. Decorre-se a seguir, detalhadamente, a modelagem de cada função objetivo contemplada, suas justificativas e implicações para o problema em questão.

\subsubsection{\(Z_1\) - Minimizar Exames Não Atendidos}
\label{subsec:fo1}

Conforme discorrido no início da Seção \ref{cap:problema}, é conhecida, mesmo que por estimativa, uma demanda \(e_{udt}\) de exames a serem realizados para uma unidade em um período qualquer, da mesma maneira que se conhece a disponibilidade de atendimento \(a_{idt}\) e a quantidade de exames possíveis \(q_{idt}\) de um médico para diariamente e períodos. Considera-se por exame não atendimento, portanto, o excedente de exames que não puderam ser atendidos pelos médicos disponíveis na ocasião, introduzido ao modelo as restrições (\ref{eq:fo1-r1}) que computam a diferenças, para cada clínica, em todos os períodos. Para modelar essa situação, consideram-se as novas variáveis \(c_{udt}\) que representam a quantidade  de exames não atendidos na unidade  \(u\), no turno \(t\), no dia \(d\).

\begin{alignat}{2}
	c_{udt} \geq e_{udt} - \sum_{i \in I_u}(q_{iut} \cdot x_{iudt}) \quad & \forall u \in U, d \in D, t \in T \label{eq:fo1-r1}            \\
	c_{udt} \in \mathbb{Z}^{*}_{+} \quad                                  & \forall i \in I, u \in U, d \in D, t \in T \label{eq:c-domain}
\end{alignat}

Dessa forma, o sistema pode ser modelado nessa ocasião como um problema de minimização, procurando obter a menor quantidade possível de exames não atendidos ao longo de todo o horizonte de planejamento.

\begin{equation}
	\min \sum_{u \in U}\sum_{d \in D}\sum_{t \in T} c_{udt} \label{eq:fo1-fo}
\end{equation}


\subsubsection{\(Z_2\) - Minimizar Trocas de Clínicas}
\label{subsec:fo2}

Para este problema, define-se por troca toda ação de remanejar um profissional de uma clínica para outra durante a transição entre dias. Tomando como exemplo a Figura \ref{fig:ex-troca-medico}, a verificação de remanejamento é binária (o médico trocou ou não de clínica entre dias), devido às restrições (\ref{eq:r2}) e (\ref{eq:r3}), que impedem o profissional de trabalhar em mais de uma unidade no dia.

\begin{figure}[h]
	\centering
	\includegraphics[width=10cm, height=5cm]{images/troca-examplo.png}
	\caption{Exemplo de trocas para um médico qualquer. }
	\label{fig:ex-troca-medico}
	\fautor
\end{figure}

Entretanto, dada a irregularidade da agenda dos profissionais, é possível que existam lacunas (dias sem trabalho) entre alocações consecutivas. A simples comparação de variáveis \(w\) entre dias consecutivos não é suficiente para detectar trocas que ocorram através dessas lacunas. Por exemplo, se um médico atua na clínica \(u\) nos dias 1 e 2, não trabalha nos dias 3 e 4, e retorna à clínica \(u' \neq u\) no dia 5, a troca de \(u\) para \(u'\) deve ser detectada. Para contornar essa limitação, introduz-se a variável de estado \(\gamma_{iud}\), baseada nos estudos em \textit{GLSP - General Lot-sizing and Scheduling Problem} de \citeonline{Fleischmann_Meyr_1997}, que registra a qual clínica o profissional está conceitualmente associado em cada dia, independentemente de estar efetivamente alocado. Esse estado persiste durante dias sem alocação, viabilizando a detecção correta de trocas ao longo de lacunas na agenda.
Dessa forma, ao invés de dinamicamente computar trocas entre unidades, é possível reduzir o problema a uma comparação estática entre o estado do profissional em dias consecutivos, mesmo que haja lacunas entre eles, garantindo a correta contabilização de trocas.

Introduz-se também o parâmetro \(CI\), vetor de tamanho \(|I|\), representando a clínica à qual o profissional estava alocado em um período anterior ao horizonte de planejamento. Na ausência de condição inicial, todos os valores do vetor são \(0\), indicando que o médico não estava previamente associado a nenhuma clínica.

\begin{small}
	\begin{tabular}{ll}
		\multicolumn{2}{l}{\textbf{Parâmetros}}                                                                    \\
		\(CI_{i}\) \quad & Unidade à qual o profissional \(i\) estava alocado no dia precedente a \(d=1\) (no dia \(d=0\)) \\
	\end{tabular}
\end{small}

\begin{small}
	\begin{tabular}{ll}
		\multicolumn{2}{l}{\textbf{Variáveis}}                                                                                                                                    \\
		\(\gamma_{iud}\)       \quad & \( \begin{cases} 1, & \text{se o estado do médico $i$ é a clínica $u$ no dia $d$ (mesmo sem alocação ativa)} \\    0, & \text{caso contrário.}
			                               \end{cases} \) \\
	\end{tabular}
\end{small}

As restrições (\ref{eq:gamma-def}--\ref{eq:gamma-sum}) definem o comportamento da variável de estado. A restrição (\ref{eq:gamma-def}) garante que, ao ser alocado a uma clínica no dia \(d\), o estado do profissional é atualizado para essa clínica. A restrição (\ref{eq:gamma-sum}) assegura que exatamente uma clínica permanece como estado corrente do profissional em cada dia, preservando o estado durante dias sem alocação.

\begin{alignat}{2}
	\gamma_{iud} \geq w_{iud} \quad         & \forall i \in I, u \in U_i, d \in D \label{eq:gamma-def} \\
	\sum_{u \in U_i} \gamma_{iud} = 1 \quad & \forall i \in I, d \in D \label{eq:gamma-sum}            \\
	\gamma_{iud} \in \{0,1\} \quad          & \forall i \in I, u \in U, d \in D \label{eq:gamma-dom}
\end{alignat}

A condição inicial do estado é definida conforme a alocação prévia ao horizonte de planejamento:

\begin{alignat}{2}
	\gamma_{i,CI_i,0} = 1 \quad & \forall i \in I \label{eq:gamma-init}
\end{alignat}

As restrições (\ref{eq:fo2-r0}) realizam a verificação de trocas entre o primeiro dia do horizonte (\(d=1\)) e a condição inicial. Computa-se troca somente caso o estado do médico no primeiro dia difira de \(CI_i\).

\begin{alignat}{2}
	z_{i1} \geq \gamma_{iu1} + \gamma_{i,CI_i,0} - 1 \quad & \forall i \in I, u \in U_i, u \neq CI_i \label{eq:fo2-r0}
\end{alignat}

As restrições (\ref{eq:fo2-r1}) computam as trocas para os demais dias. Detecta-se troca quando o estado do profissional no dia \(d\) (clínica \(u\)) difere de seu estado no dia \(d-1\) (qualquer clínica \(u' \neq u\)). Como \(\gamma\) preserva o estado durante dias sem alocação, as trocas são corretamente detectadas mesmo na presença de lacunas na agenda.

\begin{alignat}{2}
	z_{id} \geq \gamma_{iud} + \sum_{\substack{u' \in U_i \\ u' \neq u}} \gamma_{i,u',d-1} - 1 \quad & \forall i \in I, u \in U_i, d \in D - \{1\} \label{eq:fo2-r1}
\end{alignat}

O objetivo final da função é, portanto, \textbf{minimizar} a quantidade total de trocas entre clínicas para todos os profissionais durante todo o horizonte de planejamento, dada uma condição inicial de alocação, como em (\ref{eq:fo2-fo}). A variável de estado \(\gamma\) garante que as trocas sejam corretamente contabilizadas mesmo quando os profissionais possuem agendas irregulares com lacunas.

\begin{equation}
	\min \sum_{i \in I}\sum_{d \in D} z_{id} \label{eq:fo2-fo}
\end{equation}

É importante destacar a necessidade de um estado pré-existente ao utilizar computações de estados. Dessa forma, assume-se, mesmo que por estimativa (ou inicialização aleatória), uma clínica inicial a qual o médico está alocado, implicando na infactibilidade do modelo caso a representativa \(\gamma\) associada ao estado inicial do médico seja \(0\) para todas as clínicas, ou seja, o médico não esteja alocado a nenhuma clínica no período precedente ao horizonte de planejamento. Essa condição é necessária para que o modelo possa computar corretamente as trocas, mesmo que haja lacunas na agenda dos profissionais.


\chapter{Abordagens de Solução} % 4
\label{cap:mip-heuristicas}

Este capítulo detalha as abordagens de solução para o problema. Na Seção \ref{biojetivo} apresenta uma abordagem  Bi-objetivo para o problema, a abordagem $\epsilon$-restrito.  A Seção \ref{mipheuristicas}  é dedicada às MIP Heurísticas, onde são apresentados os métodos Relax-and-Fix e Fix-and-Optimize. A Seção \ref{genetico} foca em uma abordagem de Meta-heurística, explorando um Algoritmo Genético baseado no NSGA-2. Serão discutidos os operadores customizados, como cruzamento e mutação, que foram criados para adaptar o algoritmo às especificidades do problema, bem como as variações propostas, como as versões Estendida e Compacta do NSGA-2.

\section{Abordagem Bi-objetivo}\label{biojetivo}

Observa-se, pela natureza do problema, um \textit{trade-off} durante o estágio de otimização do sistema proposto: a minimização de exames não atendidos, visando à cobertura máxima, pode resultar em um número excessivo de remanejamentos de profissionais entre clínicas. Por outro lado, minimizar as trocas pode aumentar significativamente o número de exames não atendidos, comprometendo a qualidade do serviço. Essa situação reflete o desafio típico da otimização multiobjetivo, onde é necessário equilibrar múltiplos critérios conflitantes.
Tome como exemplo o cenário em que um médico com baixa disponibilidade é alocado a uma clínica com alta demanda para minimizar exames não atendidos, mas isso gera uma troca de clínica para esse médico, aumentando o número total de remanejamentos. Em contrapartida, se o mesmo médico for mantido em sua clínica atual para minimizar trocas, a clínica com alta demanda pode ficar sem cobertura adequada, aumentando os exames não atendidos.

Essa dicotomia entre os objetivos caracteriza uma abordagem multiobjetivo (neste caso, bi-objetivo) para o problema, onde é necessário encontrar um equilíbrio entre a minimização de exames não atendidos e a minimização de trocas de clínicas.

Existem diversas formas de se abordar um problema caracterizado multiobjetivo na literatura, como o escalonamento clássico por soma ponderada que de \citeonline{CLARK1983259} sobre otimização com critérios de desempenho de valor não escalar, bem como Programação de Objetivos (do inglês \textit{Goal Programming} ou GP), introduzida por \citeonline{charnes1957management} que modela múltiplas metas e minimiza os desvios em relação a elas. Outros métodos notáveis incluem o de Tchebycheff, proposto por \citeonline{Zelany_1974}, que avalia soluções pela desviação ponderada máxima em relação a um ponto ideal, e o $\epsilon$-restrito aumentado, que combina ambos. Neste trabalho, adota-se o método $\epsilon$-restrito, discutido na Seção \ref{subsec:eps-restrito-model} a seguir.

\subsection{\texorpdfstring{$\epsilon$}{epsilon}-Restrito}
\label{subsec:eps-restrito-model}

O método $\epsilon$-restrito \cite{Haimes_Wismer_1971} transforma um problema multiobjetivo em uma sequência de subproblemas mono-objetivo, ao converter os objetivos secundários em restrições limitadas por um parâmetro $\epsilon$. Variando sistematicamente $\epsilon$, obtêm-se soluções eficientes que aproximam a fronteira de Pareto, inclusive em suas regiões não convexas \cite{Saez-Aguado_Trandafir_2018}.

Formalmente, dado um problema bi-objetivo $\min\{Z_1(x), Z_2(x)\}$ com conjunto de soluções factíveis $X$, o método $\epsilon$-restrito seleciona $Z_1$ como objetivo primário e converte $Z_2$ em restrição: minimiza-se $Z_1(x)$ sujeito a $Z_2(x) \leq \varepsilon$ e às demais restrições do modelo. O valor de $\varepsilon$ é variado parametricamente para gerar diferentes pontos da fronteira de Pareto.

Neste trabalho, $Z_1$ (exames não atendidos) é a função objetivo e $Z_2$ (trocas de clínica) é convertida em restrição, pois $Z_2$ assume apenas valores inteiros não-negativos, tornando o intervalo de $\varepsilon$ finito e enumerável. Os limites do intervalo são determinados por otimização lexicográfica: o ponto Lex-$Z_1$, obtido minimizando exclusivamente $Z_1$, fornece $\varepsilon_{\min}$; o ponto Lex-$Z_2$, obtido minimizando exclusivamente $Z_2$, fornece $\varepsilon_{\max}$. A grade é composta por sete valores uniformemente espaçados em $[\varepsilon_{\min}, \varepsilon_{\max}]$, gerando sete subproblemas independentes cuja solução conjunta constitui a fronteira de Pareto aproximada de cada instância.

O modelo proposto para a abordagem bi-objetivo é apresentado a seguir, onde o objetivo primário é minimizar os exames não atendidos ($Z_1$), enquanto a minimização de trocas de clínicas ($Z_2$) é tratada como uma restrição, limitando o número de trocas a um valor máximo $\varepsilon$.
\begin{alignat}{2}
    \min \sum_{u \in U}\sum_{d \in D}\sum_{t \in T} c_{udt} \quad         &                                              \\
    \text{sujeito a:} \quad                                               &                                              \\
    \sum_{u \in U_i} x_{iudt} \leq a_{idt}  \quad                         & \forall i \in I, d \in D, t \in T            \\
    \sum_{u \in U_i} w_{iud} \leq 1  \quad                                & \forall i \in I, d \in D                     \\
    x_{iudt} \leq w_{iud} \quad                                           & \forall i \in I, u \in U, d \in D, t \in T_d \\
    w_{iud} = 0 \quad                                                     & \forall i \in I, u \notin U_i, d \in D       \\
    \sum_{i \in I_u} x_{iudt} \leq s_{udt} \quad                          & \forall u \in U, d \in D, t \in T            \\
    \sum_{i \in I_u}x_{iudt} \geq b_{udt} \quad                           & \forall u \in U, d \in D, t \in T            \\
    c_{udt} \geq e_{udt} - \sum_{i \in I_u}(q_{idt} \cdot x_{iudt}) \quad & \forall u \in U, d \in D, t \in T            \\
    \gamma_{iud} \geq w_{iud} \quad                                       & \forall i \in I, u \in U_i, d \in D          \\
    \sum_{u \in U_i} \gamma_{iud} = 1 \quad                               & \forall i \in I, d \in D                     \\
    \gamma_{i,CI_i,0} = 1 \quad                                           & \forall i \in I                              \\
    z_{i1} \geq \gamma_{iu1} + \gamma_{i,CI_i,0} - 1 \quad                & \forall i \in I, u \in U_i, u \neq CI_i      \\
    z_{id} \geq \gamma_{iud} + \sum_{\substack{u' \in U_i                                                                \\ u' \neq u}} \gamma_{i,u',d-1} - 1 \quad & \forall i \in I, u \in U_i, d \in D - \{1\} \\
    \sum_{i \in I}\sum_{d \in D} z_{id} \leq \epsilon \quad               &                                              \\
    x_{iudt} \in \{0,1\} \quad                                            & \forall i \in I, u \in U, d \in D, t \in T   \\
    w_{iud}, \gamma_{iud} \in \{0,1\} \quad                               & \forall i \in I, u \in U, d \in D            \\
    z_{id} \in \{0,1\} \quad                                              & \forall i \in I, d \in D                     \\
    c_{udt} \in \mathbb{Z}^{*}_{+} \quad                                  & \forall u \in U, d \in D, t \in T
\end{alignat}

\section{MIP Heurísticas}\label{mipheuristicas}

mat-heurísticas são estratégias que exploram a estrutura do modelo MIP para guiar a busca, combinando a formulação exata com a escalabilidade de procedimentos heurísticos. Em problemas de escalonamento em saúde, esse tipo de abordagem é bem utilizada, sobretudo em cenários nos quais o horizonte de planejamento é longo e o número de variáveis binárias torna a resolução exata onerosa. Revisões recentes destacam que heurísticas, meta-heurísticas e mat-heurísticas têm papel central em problemas de \textit{staff scheduling} e \textit{rostering} médico \cite{Ngoo_Goh_Sze_Sabar_Abdullah_Kendall_2022,Abdalkareem2021_ac}.

Neste trabalho, as mat-heurísticas são acopladas ao método $\varepsilon$-restrito discutido no Seção~\ref{biojetivo}, atuando como resolvedores internos para cada instância mono-objetivo gerada por um valor de $\varepsilon$. Dessa forma, cada ponto da fronteira de Pareto é obtido sob o mesmo custo computacional, garantindo comparabilidade entre métodos.

\subsection{Relax-and-Fix}

No problema estudado, o Relax-and-Fix decompõe o horizonte em períodos atômicos $(d,t)$ — combinações de dia $d \in D$ e turno $t \in T$, que constituem a menor unidade temporal da formulação —, mantendo a heurística compatível com a estrutura temporal do modelo. Em cada janela, apenas as variáveis associadas às decisões correntes permanecem inteiras, enquanto as de períodos futuros são temporariamente relaxadas para o domínio contínuo, reduzindo a dificuldade do subproblema.

\begin{algoritmo}
    \caption{Relax-and-Fix para alocação de médicos}
    \label{alg:rnf}
    \begin{algorithmic}[1]
        \Require Modelo MIP, \texttt{window\_size\_pct}, \texttt{overlap}, \texttt{total\_time\_limit}, \texttt{polish\_ratio}, \texttt{mip\_gap}
        \Ensure Solução factível refinada

        \State Construir a lista ordenada de períodos $(d,t)$
        \State Definir o tamanho da janela $W$ e particionar os períodos em janelas com sobreposição
        \State Reservar parte do tempo total para o \textit{polish} final

        \For{cada janela}
        \State Restaurar a integralidade das variáveis da janela corrente
        \State Relaxar as variáveis associadas a períodos posteriores
        \State Resolver o subproblema da janela
        \If{uma solução viável for encontrada}
        \State Fixar as variáveis binárias estruturais da janela
        \Else
        \State Encerrar o procedimento
        \EndIf
        \EndFor

        \State Restaurar os tipos originais de todas as variáveis
        \State Utilizar a solução construída como \textit{warm-start}
        \State Resolver o modelo completo com o tempo remanescente
        \State \Return melhor solução encontrada
    \end{algorithmic}
\end{algoritmo}

Formalmente, define-se $Q$ como o conjunto de índices dos períodos atômicos $(d,t)$ do horizonte de planejamento. Particiona-se $Q$ em $R$ conjuntos disjuntos $Q_1, \ldots, Q_R$, correspondentes às janelas temporais sucessivas. A heurística tenta encontrar uma solução resolvendo subproblemas de forma iterativa e sequencial: inicialmente, todas as variáveis binárias são relaxadas; na iteração $r=1$, resolve-se $\text{MIP}^1$ mantendo inteiras as variáveis dos períodos em $Q_1$; nas iterações $r=2,\ldots,R$, resolve-se $\text{MIP}^r$ mantendo inteiras as variáveis de $Q_r$ e fixando as variáveis estruturais $w_{iud}$ dos períodos em $Q_{r-1} \setminus U_{r-1}$ com os valores de $\text{MIP}^{r-1}$, caso factível; caso contrário, o procedimento é interrompido. A sobreposição entre janelas consecutivas é definida por $U_r = Q_r \cap Q_{r-1} \neq \emptyset$, com $r=1,\ldots,R-1$, de modo que os períodos em $U_r$ são mantidos no domínio original tanto na janela $r-1$ quanto na janela $r$, antes de serem fixados.

Adota-se como parâmetros do Algoritmo~\ref{alg:rnf} os seguintes componentes:

\begin{itemize}
    \item \texttt{window\_size\_pct} — fração do horizonte que define o tamanho de cada janela $|Q_r|$;
    \item \texttt{overlap} — número de períodos compartilhados entre janelas consecutivas ($|U_r|$);
    \item \texttt{total\_time\_limit} — tempo limite total do procedimento;
    \item \texttt{polish\_ratio} — fração do tempo total reservada para a etapa de refinamento com \textit{warm-start};
    \item \texttt{mip\_gap} — tolerância de otimalidade imposta ao \textit{solver} em cada subproblema.
\end{itemize}
    
A cada iteração, a heurística fixa principalmente as variáveis $w_{iud}$, pois elas representam a decisão estrutural de associação diária entre médico e clínica. Ao final, a solução parcial construída ao longo das janelas é utilizada como ponto de partida para uma heurística de refinamento, permitindo recuperar qualidade sem abandonar a estratégia de decomposição temporal \cite{Turhan_Bilgen_2020}. 

\subsection{Fix-and-Optimize}

O \textit{Fix-and-Optimize} utiliza a mesma decomposição temporal em períodos $(d,t)$, mas com finalidade distinta da do \textit{Relax-and-Fix}. Em vez de construir a solução progressivamente, o método parte de uma solução inicial viável e resolve sucessivos subproblemas locais, mantendo livre apenas uma vizinhança de períodos por vez.

\begin{algoritmo}
    \caption{Fix-and-Optimize para alocação de médicos}
    \label{alg:fno}
    \begin{algorithmic}[1]
        \Require Modelo MIP com solução inicial, \texttt{neighborhood\_size\_pct}, \texttt{max\_iterations}, \texttt{total\_time\_limit}, \texttt{mip\_gap}, \texttt{no\_improvement\_limit}
        \Ensure Solução refinada

        \State Construir a lista ordenada de períodos $(d,t)$
        \State Definir o tamanho da vizinhança
        \State Inicializar a melhor solução com a solução inicial disponível

        \For{cada iteração enquanto houver tempo e não houver estagnação}
        \State Selecionar uma vizinhança de períodos
        \State Desafixar todas as variáveis
        \State Fixar as variáveis fora da vizinhança com base na melhor solução corrente
        \State Resolver o subproblema restrito
        \If{uma solução melhor for encontrada}
        \State Atualizar a melhor solução e reiniciar o contador de estagnação
        \Else
        \State Incrementar o contador de estagnação
        \EndIf
        \EndFor

        \State Desafixar todas as variáveis
        \State Restaurar a melhor solução como \textit{warm-start}
        \State \Return melhor solução encontrada
    \end{algorithmic}
\end{algoritmo}

Na prática, o procedimento mais eficiente é utilizar R\&F como fase construtiva e F\&O como fase de refinamento. Primeiro, obtém-se uma solução factível rapidamente; depois, aplica-se busca local orientada por \textit{MIP} para melhorar a qualidade final. Essa estratégia híbrida é compatível com o método $\varepsilon$-restrito e com os limites de tempo adotados em cada instância \cite{Turhan_Bilgen_2020}.

Do ponto de vista metodológico, essa escolha se justifica porque o problema apresenta estrutura temporal clara, inúmeras variáveis binárias e múltiplas restrições operacionais, favorecendo decomposição e refinamento iterativo. Além disso, revisões recentes em scheduling em saúde apontam que abordagens híbridas são particularmente adequadas para problemas reais, nos quais a escalabilidade é tão importante quanto a qualidade da solução \cite{Vermuyten_Namorado_Rosa_Marques_Belien_Barbosa_Povoa_2018}.

\section{Meta-Heurísticas - Algoritmo Genético}\label{genetico}


Além das MIP-heurísticas discutidas anteriormente, este trabalho também considera uma abordagem baseada em algoritmos genéticos para tratar o problema de alocação de médicos em múltiplas clínicas. A motivação para essa escolha decorre do fato de que o problema estudado apresenta elevada complexidade combinatória, grande número de variáveis de decisão e objetivos conflitantes, características que dificultam a exploração eficiente do espaço de soluções por métodos exclusivamente exatos. Nesse contexto, a busca evolutiva oferece um mecanismo natural para produzir e refinar, ao longo do processo de otimização, múltiplas escalas candidatas com perfis distintos de compromisso entre cobertura da demanda e estabilidade das alocações \cite{Hamid_Tavakkoli_Moghaddam_Golpaygani_Vahedi_Nouri_2020,Tan_Gan_Ren_2019}.

No escopo deste estudo, um algoritmo genético deve ser entendido como um procedimento iterativo de busca populacional em que cada indivíduo representa uma solução completa para o horizonte de planejamento. Em vez de construir uma única escala por vez, o método trabalha simultaneamente com uma população de soluções candidatas, permitindo explorar diferentes regiões do espaço factível em paralelo. Essa característica é particularmente importante no problema em questão, pois pequenas mudanças nas decisões de alocação médico--clínica--dia--turno podem produzir efeitos significativos sobre os objetivos, e diferentes estruturas de agenda podem apresentar desempenhos semelhantes sob perspectivas distintas.

A aplicação do algoritmo genético é diretamente ancorada na modelagem matemática apresentada no Capítulo~\ref{cap:problema}. Em particular, as soluções manipuladas pela meta-heurística refletem a mesma estrutura de decisão associada às variáveis de alocação por turno $x_{iudt}$ e de associação diária $w_{iud}$. Com isso, a busca meta-heurística não atua sobre uma abstração desconectada do problema real, mas sobre representações compatíveis com a lógica operacional da rede de clínicas. Essa aderência é importante tanto para garantir coerência metodológica entre modelagem e solução quanto para facilitar a interpretação prática dos resultados.

Do ponto de vista operacional, uma solução candidata deve representar, para cada médico, sua possível alocação ao longo do horizonte de planejamento, respeitando a disponibilidade informada, as clínicas às quais o profissional pode ser associado, os limites de salas por unidade e a necessidade de cobertura mínima quando existe demanda. Em função disso, o espaço de busca do problema é extremamente restrito e esparso: embora o número de combinações possíveis seja muito elevado, apenas uma fração reduzida dessas combinações resulta em escalas plausíveis. Consequentemente, a aplicação do algoritmo genético neste trabalho exige mecanismos que reduzam a geração de soluções inviáveis e orientem a evolução para regiões estruturalmente promissoras.

A primeira etapa do procedimento evolutivo consiste na geração de uma população inicial. Em problemas altamente restritos como o presente, essa fase é decisiva para o desempenho do algoritmo. Caso os indivíduos iniciais sejam produzidos de forma puramente aleatória, grande parte do esforço computacional subsequente será desperdiçada na exploração de soluções de baixa qualidade ou inviáveis. Por isso, a metodologia adotada neste estudo privilegia a construção de indivíduos iniciais com algum nível de coerência estrutural, de modo que a evolução tenha como ponto de partida escalas compatíveis com a dinâmica do problema.

Uma vez definida a população inicial, cada indivíduo é avaliado conforme os critérios de desempenho do estudo. Como o problema é bi-objetivo, a avaliação considera simultaneamente a quantidade total de exames não atendidos e o número total de trocas de clínicas realizadas ao longo do horizonte. O primeiro critério mede a efetividade operacional da escala em absorver a demanda prevista. O segundo mede o grau de estabilidade da agenda dos profissionais, evitando remanejamentos excessivos entre unidades. Dessa forma, a qualidade de uma solução não depende de um único escalar agregado, mas do balanço entre dois critérios conflitantes e igualmente relevantes para a operação.

Em termos metodológicos, isso significa que o algoritmo genético não deve simplesmente procurar uma solução ``melhor'' em sentido absoluto, mas sim explorar um conjunto de soluções potencialmente competitivas segundo diferentes compromissos entre os objetivos. Algumas escalas podem favorecer cobertura mais agressiva da demanda às custas de maior mobilidade dos médicos; outras podem preservar maior estabilidade logística, aceitando níveis mais altos de exames não atendidos. A função da busca evolutiva é justamente navegar entre essas possibilidades e construir uma população que represente adequadamente esse espectro de alternativas.

A etapa seguinte consiste na evolução propriamente dita da população. Em linhas gerais, o algoritmo seleciona indivíduos promissores, combina partes de suas estruturas por meio de recombinação e, em seguida, aplica pequenas modificações por mutação. A seleção privilegia soluções mais competitivas, a recombinação permite transferir padrões favoráveis entre indivíduos distintos e a mutação introduz perturbações controladas capazes de ampliar a exploração do espaço de busca. A repetição desse ciclo ao longo de várias gerações faz com que a população tenda a concentrar soluções progressivamente mais adequadas aos objetivos definidos.

No contexto da alocação de médicos, a recombinação não pode ser encarada como simples mistura arbitrária de bits ou genes. As agendas apresentam forte dependência temporal e estrutural: alocações em dias consecutivos afetam trocas de clínicas, e decisões em um turno podem repercutir sobre a coerência da alocação diária. Assim, a combinação de duas soluções deve preservar, tanto quanto possível, blocos semânticos relevantes da agenda. De forma análoga, a mutação também precisa ser cuidadosamente controlada, pois alterações excessivamente agressivas podem destruir propriedades desejáveis da solução corrente e conduzir o algoritmo a regiões inviáveis do espaço.

Outro ponto importante é que, em problemas reais de escalonamento, a busca genética raramente se beneficia de uma separação rígida entre viabilidade e qualidade. Em muitos casos, soluções muito promissoras do ponto de vista dos objetivos podem apresentar pequenas inconsistências, enquanto soluções plenamente viáveis podem ser pouco competitivas. Por essa razão, a metodologia adotada neste estudo prioriza a produção de soluções plausíveis por construção e por verificação local, reduzindo a dependência de reparos extensivos e evitando que a evolução se torne excessivamente penalizada por inviabilidades estruturais. Isso melhora a eficiência do processo e torna a busca mais aderente ao problema prático.

A atualização da população ao final de cada geração cumpre dois papéis simultâneos. O primeiro é preservar os indivíduos mais competitivos encontrados até o momento, garantindo que a evolução não perca padrões favoráveis já descobertos. O segundo é manter diversidade populacional suficiente para evitar convergência prematura para uma única região do espaço de busca. Essa tensão entre intensificação e diversificação é central na aplicação de algoritmos genéticos ao presente problema: se a busca se tornar excessivamente conservadora, pode estagnar cedo em soluções medianas; se for excessivamente dispersa, pode desperdiçar tempo computacional em regiões pouco promissoras.

De modo geral, a estrutura de aplicação do algoritmo genético neste estudo pode ser sintetizada pelo Algoritmo~\ref{alg:ga-geral}. O pseudocódigo resume a lógica comum às abordagens evolutivas empregadas no trabalho, independentemente de suas especializações posteriores.

\begin{algoritmo}
    \caption{Estrutura do algoritmo genético para alocação de médicos}
    \label{alg:ga-geral}
    \begin{algorithmic}[1]
        \Require Instância do problema, parâmetros do algoritmo, critério de parada
        \Ensure Conjunto de soluções finais de boa qualidade

        \State Gerar população inicial de soluções candidatas
        \State Avaliar cada indivíduo segundo os objetivos do problema
        \State Definir a população corrente $P$

        \While{critério de parada não for satisfeito}
        \State Selecionar indivíduos de $P$ para reprodução
        \State Aplicar recombinação para gerar descendentes
        \State Aplicar mutação nos descendentes
        \State Avaliar os novos indivíduos gerados
        \State Atualizar a população corrente preservando qualidade e diversidade
        \EndWhile

        \State \Return conjunto final de soluções obtidas
    \end{algorithmic}
\end{algoritmo}

O pseudocódigo acima descreve o esqueleto metodológico da abordagem genética. A geração da população inicial define o ponto de partida da busca; a avaliação quantifica o desempenho de cada escala candidata; a recombinação e a mutação promovem exploração e refinamento; e a atualização populacional regula o avanço da busca ao longo das gerações. O critério de parada pode ser definido por número de gerações, número de avaliações ou limite de tempo, a depender do protocolo experimental adotado.

No problema de alocação de médicos, essa estrutura é particularmente útil porque permite evoluir escalas completas sem depender de uma sequência rígida de decisões construtivas. Em vez de fixar progressivamente partes do horizonte, como nas MIP-heurísticas, a abordagem genética revisita continuamente soluções inteiras, recombinando padrões de alocação e ajustando escolhas locais iterativamente. Isso amplia a capacidade de exploração global e aumenta a chance de identificar soluções competitivas em cenários de alta complexidade combinatória.

Em síntese, o papel do algoritmo genético neste trabalho é fornecer um mecanismo de busca robusto para explorar o espaço de escalas factíveis e quase factíveis sob dois objetivos conflitantes. Sua utilização é especialmente adequada quando se deseja ir além de uma solução única e investigar diferentes padrões de compromisso entre atendimento da demanda e estabilidade operacional. A subseção a seguir descreve a instância específica desse paradigma evolutivo empregada no estudo, a saber, o NSGA-II.

\subsection{NSGA-2}

Para a resolução meta-heurística bi-objetivo do problema, adotou-se o NSGA-II \cite{Deb_Pratap_Agarwal_Meyarivan_2002} como baseline evolutivo. No contexto deste estudo, sua função é construir uma aproximação da fronteira de Pareto formada pelas soluções não dominadas em termos de exames não atendidos e trocas de clínicas. Diferentemente de uma abordagem que reduz o problema a uma única função escalar, o NSGA-II mantém explicitamente o caráter bi-objetivo da formulação, permitindo que diferentes soluções concorrentes coexistam e evoluam simultaneamente ao longo do processo.

A aplicação do NSGA-II parte da mesma estrutura decisória formalizada no Capítulo~\ref{cap:problema}. Na formulação base empregada no projeto, cada indivíduo é representado por um vetor binário concatenando as variáveis $x_{iudt}$ e $w_{iud}$. Em outras palavras, o cromossomo codifica diretamente, ao nível de detalhe fino, tanto as alocações de médicos por clínica, dia e turno quanto a associação diária do profissional à unidade. Essa escolha torna a representação plenamente compatível com a modelagem exata, ao custo de uma dimensionalidade elevada do espaço de busca.

A avaliação de cada indivíduo é feita com base nos dois objetivos do problema. O primeiro objetivo, denotado por $F_1$, corresponde à demanda não atendida total ao longo do horizonte. O segundo, $F_2$, corresponde ao número de trocas de clínica por médico ao longo dos dias. Dessa forma, a qualidade de uma solução é determinada por um vetor bi-objetivo, e não por um único escalar. Em termos práticos, isso implica que a comparação entre indivíduos deve respeitar relações de dominância: uma solução é preferível a outra quando não é pior em nenhum objetivo e é estritamente melhor em pelo menos um deles.

O mecanismo central do NSGA-II para lidar com isso é a ordenação não-dominada da população. Após a avaliação dos indivíduos, a população é particionada em frentes de Pareto: a primeira frente reúne as soluções não dominadas, a segunda reúne aquelas dominadas apenas pelas da primeira frente, e assim sucessivamente. Esse procedimento permite organizar a população por níveis de qualidade multiobjetivo, sem impor previamente pesos ou prioridades entre os dois objetivos. No presente problema, isso é especialmente importante porque não se deseja privilegiar, de antemão, apenas cobertura da demanda ou apenas estabilidade das agendas, mas sim explorar diferentes balanços entre ambos.

\begin{figure}[htb]
    \centering
    \includegraphics[width=0.77\linewidth]{images/04_nsga-02_ordenacao.png}
    \caption{Ordenação não-dominada.}
    \label{fig:nsga-ordenacao}
    \fautor
\end{figure}

Além da ordenação não-dominada, o NSGA-II emprega a distância de crowding para promover diversidade dentro de cada frente. No contexto do estudo, esse mecanismo evita que a população convirja rapidamente para uma região muito estreita da fronteira aproximada. Como existem múltiplas escalas potencialmente interessantes sob perspectivas diferentes, a preservação de diversidade é essencial para que o algoritmo represente adequadamente a variedade de compromissos possíveis entre os objetivos. Sem esse mecanismo, a evolução poderia concentrar soluções em um pequeno subconjunto da fronteira, comprometendo a utilidade gerencial dos resultados.

\begin{figure}[htb]
    \centering
    \includegraphics[width=0.77\linewidth]{images/04_nsga-03_crowding-distance.png}
    \caption{Distância \textit{crowding} entre indivíduos.}
    \label{fig:nsga-crowding}
    \fautor
\end{figure}

A execução do NSGA-II neste trabalho inicia-se com a construção de um objeto de problema compatível com a instância de alocação de médicos. Esse objeto recebe os dados da instância, codifica a lógica de avaliação dos objetivos e incorpora a gestão das restrições relevantes. Em seguida, são inicializados os operadores de amostragem, recombinação e mutação definidos para o contexto do problema. A população inicial é então gerada e avaliada, constituindo a primeira aproximação do conjunto de soluções.

A partir desse ponto, o algoritmo entra em um ciclo iterativo de evolução. Em cada geração, selecionam-se indivíduos da população corrente para reprodução, produzindo descendentes por meio da aplicação combinada de \textit{crossover} e mutação. Esses descendentes são então avaliados da mesma forma que os indivíduos parentais. Na sequência, pais e filhos são reunidos em uma população temporária ampliada, sobre a qual se aplica novamente a ordenação não-dominada. A nova população da geração seguinte é construída a partir das melhores frentes, utilizando a distância de crowding para decidir quais indivíduos sobrevivem quando uma frente precisa ser truncada.

Esse mecanismo de sobrevivência elitista é particularmente relevante. Ao formar a nova população a partir da união entre pais e filhos, o algoritmo garante que boas soluções já encontradas não sejam perdidas apenas por flutuações estocásticas da reprodução. Em problemas complexos como o presente, nos quais soluções competitivas podem ser difíceis de gerar, preservar explicitamente essas estruturas favoráveis aumentam a estabilidade da evolução e contribui para a construção progressiva de uma fronteira de Pareto de melhor qualidade.

No estudo desenvolvido, a configuração base do NSGA-II utiliza população de tamanho 100 e critérios de parada definidos por número de gerações, número de avaliações ou limite temporal, dependendo do experimento executado. Essa flexibilidade é importante porque diferentes horizontes de planejamento e diferentes tamanhos de instância podem demandar orçamentos computacionais distintos \cite{Lavygina_Welsh_Crispin_2019,Khalili_Shahrezaei_Abri_2020}. Em especial, quando o horizonte temporal cresce, a avaliação das soluções e a manipulação dos cromossomos tornam-se mais custosas, exigindo cuidado na definição do esforço computacional disponível.

Outro aspecto metodológico relevante é que o problema tratado apresenta espaço factível altamente esparso. Em consequência, a simples aplicação do NSGA-II em sua forma totalmente genérica poderia gerar excesso de indivíduos de baixa utilidade. Por isso, a aplicação adotada neste trabalho utiliza uma versão base já especializada para o problema de escalonamento médico, com operadores estruturados para produzir soluções plausíveis e reduzir a destruição da coerência temporal das agendas. Ainda assim, nesta subseção o foco permanece no funcionamento global do NSGA-II, e não nos detalhes específicos dos operadores, tratados em partes posteriores do capítulo.

Sob a ótica da solução do problema, o papel do NSGA-II é duplo. Primeiro, ele atua como mecanismo de otimização, empurrando a população em direção a soluções com menor demanda não atendida e menor número de trocas. Segundo, ele atua como mecanismo de exploração da fronteira eficiente, produzindo simultaneamente múltiplas escalas não dominadas que podem ser comparadas entre si. Esse segundo papel é particularmente valioso no problema estudado, pois diferentes gestores podem preferir soluções distintas conforme a prioridade operacional do período analisado: em alguns cenários, atender mais demanda pode ser fundamental; em outros, reduzir remanejamentos pode ser mais desejável.

\newpage

\begin{figure}[!htb]
    \centering
    \includegraphics[width=0.42\linewidth]{images/04_nsga-04_torneio.png}
    \caption{Seleção por torneio binário.}
    \label{fig:nsga-torneio}
    \fautor
\end{figure}

A lógica de execução do NSGA-II empregada neste trabalho pode ser sintetizada pelo Algoritmo~\ref{alg:nsga2-base}. O pseudocódigo abaixo descreve a estrutura principal do método no contexto da alocação de médicos.

\begin{algoritmo}
    \caption{NSGA-II base para alocação de médicos}
    \label{alg:nsga2-base}
    \begin{algorithmic}[1]
        \Require Instância do problema, parâmetros dos operadores, tamanho da população, critério de parada, semente aleatória
        \Ensure Fronteira de Pareto aproximada e histórico de execução

        \State Carregar a instância e construir o problema de alocação
        \State Inicializar os operadores evolutivos
        \State Criar o algoritmo NSGA-II com tamanho populacional predefinido
        \State Gerar a população inicial $P_0$
        \State Avaliar os objetivos de todos os indivíduos em $P_0$
        \State Definir $P \leftarrow P_0$

        \While{critério de parada não for satisfeito}
        \State Selecionar indivíduos de $P$ para reprodução
        \State Gerar descendentes $Q$ por crossover e mutação
        \State Avaliar os indivíduos de $Q$
        \State Construir $R \leftarrow P \cup Q$
        \State Aplicar ordenação não-dominada em $R$
        \State Construir a nova população $P$ por frentes e distância de crowding
        \EndWhile

        \State Extrair as soluções não dominadas de $P$
        \State \Return conjunto aproximado de soluções de Pareto
    \end{algorithmic}
\end{algoritmo}

\begin{figure}[htb]
    \centering
    \includegraphics[width=\linewidth]{images/04_nsga-01_laco-principal.png}
    \caption{Fluxo principal do NSGA-II.}
    \label{fig:nsga-laco}
    \fautor
\end{figure}

O Algoritmo~\ref{alg:nsga2-base} resume o funcionamento do método tal como aplicado neste trabalho. A geração da população inicial fornece um ponto de partida distribuído no espaço de busca. O conjunto de descendentes $Q$ introduz novas combinações e pequenas perturbações sobre as estruturas já conhecidas. A união $R = P \cup Q$ materializa o elitismo do método, enquanto a ordenação não-dominada e a distância de \textit{crowding} organizam a sobrevivência multiobjetiva da população. Ao final, o conjunto de soluções não dominadas remanescente representa a aproximação obtida para a fronteira de Pareto sob o orçamento computacional disponível.

Do ponto de vista interpretativo, cada solução final produzida pelo NSGA-II corresponde a uma escala médica completa e operacionalmente analisável. Isso significa que os resultados do algoritmo não se limitam a valores abstratos de função objetivo: cada ponto da fronteira aproximada pode ser traduzido em um plano concreto de alocação de médicos por clínicas, dias e turnos. Essa propriedade é importante para a aplicação prática do método, pois permite ao decisor inspecionar diretamente as consequências operacionais de diferentes \textit{trade-offs} entre cobertura da demanda e remanejamento de profissionais.

Em síntese, o NSGA-II é utilizado neste estudo como a principal base meta-heurística para exploração direta do problema bi-objetivo. Sua estrutura é adequada ao cenário analisado porque combina comparação por dominância, preservação de diversidade e elitismo, elementos que favorecem a obtenção de múltiplas soluções competitivas em espaços de busca complexos. A partir dessa base, as demais variantes do projeto introduzem adaptações adicionais para melhorar diversidade, escalabilidade e robustez computacional, mas o núcleo conceitual da abordagem evolutiva permanece o descrito nesta subseção.

\subsection{Operadores Genéticos}
\label{subsec:custom-operators}

A aplicação do NSGA-2 ao problema de alocação de médicos exigiu a construção de operadores especializados. Embora o mecanismo central do algoritmo permaneça baseado em ordenação não-dominada, elitismo e \textit{crowding distance}, a qualidade prática da busca depende fortemente da forma como os indivíduos são gerados e transformados ao longo das gerações. No problema estudado, o espaço de soluções é altamente restrito, de modo que operadores genéricos tenderiam a produzir grande volume de soluções inviáveis ou estruturalmente inconsistentes. Por essa razão, foram empregados operadores desenhados para respeitar a semântica da agenda médica e preservar, tanto quanto possível, a coerência temporal das decisões.

De forma geral, os operadores customizados cumprem três papéis distintos. O primeiro é produzir uma população inicial com maior plausibilidade estrutural, reduzindo o custo de iniciar a evolução em regiões improdutivas do espaço de busca. O segundo é recombinar soluções preservando blocos de agenda relevantes, em vez de misturar genes de maneira excessivamente destrutiva. O terceiro é introduzir perturbações locais controladas, capazes de melhorar ou diversificar a população sem romper a estrutura básica do indivíduo. Em conjunto, esses mecanismos tornam o NSGA-2 mais aderente ao problema real e aumentam a eficiência da busca multiobjetivo.

\subsubsection{Geração da População inicial}

A geração da população inicial é feita por uma rotina de amostragem construtiva considerando as  restrições. Em vez de sortear indivíduos inteiramente ao acaso, o método constrói soluções de maneira incremental, priorizando candidatos factíveis e verificando restrições locais ao longo do processo. Essa escolha é importante porque, no problema estudado, a simples aleatoriedade produziria muitas escalas incompatíveis com a disponibilidade médica, com a elegibilidade das clínicas ou com a capacidade das unidades.

A construção inicial combina diferentes estratégias, cada uma com uma ênfase específica. Algumas soluções priorizam demanda; outras, cobertura mínima; outros, equilíbrios de carga, antecipação temporal ou construção gulosa com randomização. Além disso, o método pode reservar uma fração da população para soluções extremas, isto é, mais orientadas à minimização de exames não atendidos ou à minimização de trocas de clínicas. Essa combinação melhora a diversidade inicial e ajuda a distribuir a população em diferentes regiões da fronteira de Pareto desde o início da execução.

O pseudocódigo da amostragem é apresentado no Algoritmo~\ref{alg:nsga2-sampling}. Sua lógica é simples: definir quantos indivíduos serão gerados por estratégias centrais e quantos por estratégias extremas; para cada solução, escolher uma estratégia, construir a agenda incrementalmente, reforçar cobertura mínima e verificar duplicidade antes de aceitar o indivíduo na população. A presença de um mecanismo do tipo \textit{tabu} evita repetição excessiva de soluções muito semelhantes, o que é particularmente útil em espaços factíveis estreitos.

\begin{algoritmo}
    \caption{Geração da População Inicial}
    \label{alg:nsga2-sampling}
    \begin{algorithmic}[1]
        \Require Problema, tamanho da população, estratégias construtivas
        \Ensure População inicial $P_0$

        \State Definir quantidades de soluções centrais e extremas
        \State Inicializar $P_0 \leftarrow \emptyset$

        \While{$|P_0|$ for menor que o tamanho desejado}
        \State Selecionar uma estratégia de construção
        \State Gerar candidatos válidos prioritários
        \State Construir uma solução por inserção gulosa com verificações locais
        \State Executar reforço de cobertura mínima
        \If{a solução não for duplicata}
        \State Inserir solução em $P_0$
        \Else
        \State Tentar nova construção com randomização controlada
        \EndIf
        \EndWhile

        \State \Return $P_0$
    \end{algorithmic}
\end{algoritmo}

\begin{figure}[htb]
    \centering
    \includegraphics[width=\linewidth]{images/01_sampling.png}
    \caption{Etapas do operador de amostragem construtiva: (a) identificação de candidatos, (b) seleção da estratégia, (c) validação e (d) codificação}
    \label{fig:nsga-sampling}
    \fautor
\end{figure}

A principal vantagem do Algoritmo~\ref{alg:nsga2-sampling} é deslocar parte do esforço de viabilidade para a própria fase de inicialização. Em vez de depender apenas da seleção natural da meta-heurística para eliminar indivíduos ruins, o método já inicia a busca com soluções mais coerentes, o que melhora a qualidade média da população e acelera a evolução nas primeiras gerações.

\subsubsection{Cruzamento por blocos diários}

O operador de cruzamento utilizado no estudo é baseado em blocos diários. A ideia central é decodificar ambos os pais, percorrer os dias do horizonte e, com uma determinada probabilidade, trocar entre eles todos os genes associados àquele dia. Em vez de recombinar posições isoladas do cromossomo, o método troca blocos temporais inteiros, preservando estruturas locais de agenda.

A justificativa para essa escolha é direta. No problema de alocação de médicos, o dia é uma unidade de decisão fortemente estruturada, pois a clínica diária do médico organiza as decisões por turno e afeta diretamente o cálculo de trocas entre unidades. Um cruzamento genérico, ao nível de gene isolado, poderia facilmente destruir coerência entre $x_{iudt}$ e $w_{iud}$, gerando descendentes inconsistentes. O cruzamento por dia reduz esse risco, pois mantém juntas as decisões que, do ponto de vista operacional, fazem parte do mesmo bloco lógico.

O funcionamento do operador é mostrado no Algoritmo~\ref{alg:nsga2-crossover}. Primeiro, os pais são decodificados para estruturas explícitas de alocação. Em seguida, percorre-se cada dia do horizonte e decide-se, estocasticamente, se haverá troca daquele bloco. Quando a troca ocorre, intercambiam-se simultaneamente as alocações por turno e a associação diária à clínica. Ao final, os descendentes são re-codificados no formato usado pelo algoritmo.

\begin{algoritmo}
    \caption{Cruzamento por blocos diários}
    \label{alg:nsga2-crossover}
    \begin{algorithmic}[1]
        \Require Dois indivíduos pais $p_1$ e $p_2$, probabilidade de troca diária
        \Ensure Dois descendentes $f_1$ e $f_2$

        \State Decodificar $p_1$ em $\textit{var\_x}_1$ e $\textit{var\_w}_1$
        \State Decodificar $p_2$ em $\textit{var\_x}_2$ e $\textit{var\_w}_2$

        \For{cada dia $d \in D$}
        \If{ocorrer troca para o dia $d$}
        \State Trocar os blocos de $x$ do dia $d$ entre os pais
        \State Trocar os blocos de $w$ do dia $d$ entre os pais
        \EndIf
        \EndFor

        \State Re-codificar as estruturas geradas
        \State \Return $f_1$ e $f_2$
    \end{algorithmic}
\end{algoritmo}

\begin{figure}[htb]
    \centering
    \includegraphics[width=0.85\linewidth]{images/02_crossover.png}
    \caption{Operação de cruzamento por blocos diários.}
    \label{fig:nsga-crossover}
    \fautor
\end{figure}

A leitura do Algoritmo~\ref{alg:nsga2-crossover} mostra que o operador atua como um mecanismo de recombinação estrutural, e não apenas combinatória. Seu papel é transferir padrões diários de agenda entre soluções promissoras, preservando a coerência temporal do indivíduo e reduzindo o caráter destrutivo da recombinação.

\subsubsection{Mutação por realocação de slot}

A mutação adotada no NSGA-2 customizado segue a lógica de realocação de slot. Em vez de alterar arbitrariamente vários genes do indivíduo, o operador seleciona uma alocação ativa e tenta movê-la para uma clínica alternativa viável no mesmo dia e turno. Com isso, a mutação passa a atuar como um mecanismo de ajuste fino, adequado a um problema em que pequenas mudanças locais podem produzir efeitos relevantes sobre os dois objetivos.

O procedimento começa com a decodificação do indivíduo. Em seguida, o operador escolhe uma atribuição ativa $(i,u,d,t)$ e constrói o conjunto de clínicas candidatas para o mesmo médico, dia e turno. Essa construção já considera disponibilidade, elegibilidade da clínica, capacidade da unidade e consistência da associação diária. Se houver pelo menos uma alternativa válida, o movimento é realizado e o indivíduo é codificado novamente; caso contrário, novas tentativas são feitas até um limite máximo.

O Algoritmo~\ref{alg:nsga2-mutation} resume esse processo. Seu comportamento é conservador por construção: ele não força uma alteração a qualquer custo, mas procura uma mudança local plausível. Se nenhum movimento válido for encontrado após o número máximo de tentativas, o indivíduo original é preservado. Essa decisão evita degradações artificiais da população por mutações excessivamente agressivas.

\begin{algoritmo}
    \caption{Mutação por realocação de slot}
    \label{alg:nsga2-mutation}
    \begin{algorithmic}[1]
        \Require Indivíduo $s$, número máximo de tentativas
        \Ensure Indivíduo mutado $s'$

        \State Decodificar $s$ em $\textit{var\_x}$ e $\textit{var\_w}$

        \For{tentativa $=1$ até o máximo permitido}
        \State Selecionar uma alocação ativa $(i,u,d,t)$
        \State Construir o conjunto de clínicas alternativas viáveis
        \If{o conjunto de alternativas não for vazio}
        \State Escolher uma nova clínica $u'$
        \State Realocar a atribuição para $(i,u',d,t)$
        \State Atualizar a associação diária correspondente
        \State Re-codificar e retornar o indivíduo modificado
        \EndIf
        \EndFor

        \State \Return o indivíduo original
    \end{algorithmic}
\end{algoritmo}

\begin{figure}[htb]
    \centering
    \includegraphics[width=0.95\linewidth]{images/03_mutacao.png}
    \caption{Mutação mista: (a) realocação de \textit{slot}, (b) adição ou remoção de atribuição e (c) troca de turnos.}
    \label{fig:nsga-mutacao}
    \fautor
\end{figure}

A principal contribuição do Algoritmo~\ref{alg:nsga2-mutation} é permitir exploração local sem exigir reparos globais complexos. Em termos práticos, ele funciona como um refinamento pontual sobre soluções já razoáveis, ajudando a escapar de regiões estagnadas da população e a produzir novos compromissos entre cobertura e estabilidade das agendas.

\newpage

\subsection{NSGA-2 Extendido}
\label{subsec:nsga2-extended}

A variante NSGA-2-Estendido foi desenvolvida para lidar com limitações observadas na aplicação direta do NSGA-2 básico ao problema estudado. Em especial, verificou-se que, para instâncias maiores ou mais restritivas, a população podia se concentrar rapidamente em regiões centrais da fronteira aproximada, com perda de diversidade e baixa cobertura dos extremos. Para contornar esse comportamento, a versão estendida preserva o núcleo do NSGA-2, mas incorpora mecanismos complementares de diversidade, exploração orientada e refinamento.

A primeira modificação diz respeito à inicialização enviesada para extremos. Em vez de construir toda a população inicial de forma homogênea, uma fração dos indivíduos pode ser gerada com maior orientação à minimização de um dos objetivos. Isso aumenta a probabilidade de que a fronteira aproximada já comece a execução com melhor cobertura de soluções extremas. Em seguida, o algoritmo pode empregar controle adaptativo de diversidade, ajustando parâmetros internos quando detecta excesso de concentração populacional.

Outra extensão relevante é a evolução em dois estágios. Em uma fase inicial, o algoritmo favorece exploração e amplitude de cobertura; em uma fase posterior, enfatizando refinamento das soluções já encontradas. Além disso, a implementação permite reinjeção periódica de extremos, manutenção de um arquivo de soluções-limite e busca local sobre a elite da população. Em conjunto, esses mecanismos procuram reduzir convergência prematura e melhorar a qualidade final da fronteira aproximada.

O pseudocódigo geral da variante é dado no Algoritmo~\ref{alg:nsga2-extended}. Sua leitura deve ser feita como uma camada adicional sobre o NSGA-2 básico. A estrutura principal do algoritmo permanece a mesma, mas antes, durante e após o ciclo evolutivo padrão são ativados módulos complementares conforme a configuração experimental adotada.

\newpage

\begin{algoritmo}
    \caption{NSGA-2-Estendido para alocação de médicos}
    \label{alg:nsga2-extended}
    \begin{algorithmic}[1]
        \Require Problema, operadores customizados, parâmetros das extensões, critério de parada
        \Ensure Fronteira de Pareto aproximada e estatísticas de aprimoramento

        \State Configurar módulos opcionais de extensão
        \If{a inicialização extrema estiver habilitada}
        \State Gerar população inicial com viés para extremos
        \Else
        \State Gerar população inicial padrão
        \EndIf
        \State Avaliar a população inicial $P$

        \While{critério de parada não for satisfeito}
        \State Determinar o estágio atual da evolução
        \If{o controle adaptativo de diversidade estiver habilitado}
        \State Medir diversidade e ajustar parâmetros
        \EndIf
        \State Executar o passo padrão do NSGA-2
        \If{a reinjeção periódica estiver habilitada}
        \State Reinserir indivíduos extremos na população
        \EndIf
        \If{o arquivo de fronteira estiver habilitado}
        \State Atualizar soluções-limite armazenadas
        \EndIf
        \If{a busca local de elite estiver habilitada}
        \State Refinar uma fração dos melhores indivíduos
        \EndIf
        \EndWhile

        \State Extrair as soluções não dominadas finais
        \State \Return fronteira aproximada e estatísticas
    \end{algorithmic}
\end{algoritmo}

O ponto mais importante do Algoritmo~\ref{alg:nsga2-extended} é que ele não substitui a dinâmica do NSGA-2 básico; ele a reforça. A inicialização extrema melhora a cobertura inicial da fronteira, o controle adaptativo reage à perda de diversidade, a reinjeção e o arquivo de fronteira preservam amplitude de soluções, e a busca local sobre a elite intensifica a qualidade em regiões promissoras. No contexto do problema estudado, essa arquitetura mostrou-se adequada quando a formulação base apresentava dificuldade em manter diversidade ou em explorar adequadamente soluções extremas.

\subsection{NSGA-2 Compacto}
\label{subsec:nsga2-compact}

A motivação para a formulação compacta decorre diretamente da dimensionalidade da representação básica. Quando o cromossomo codifica explicitamente todas as variáveis $x_{iudt}$ e $w_{iud}$, o tamanho do indivíduo cresce rapidamente com o número de médicos, clínicas, dias e turnos. Em horizontes longos, esse crescimento compromete o número de gerações possíveis dentro de um tempo computacional fixo. Para reduzir esse custo, foi adotada uma variante compacta em que a decisão genética principal passa a ser a associação diária médico--clínica.

\begin{figure}[htb]
    \centering
    \includegraphics[width=0.9\linewidth]{images/05_encoding.png}
    \caption{Comparação dos \textit{encoders} NSGA-ii-e e NSGA-II-c.}
    \label{fig:nsga-encoding}
    \fautor
\end{figure}

Na formulação compacta, cada gene representa um par médico--dia. O valor do gene indica a clínica atribuída naquele dia ou ausência de alocação. Com isso, o cromossomo passa a codificar diretamente a estrutura diária da agenda, isto é, uma representação análoga à variável $w$, sem carregar explicitamente todas as alocações por turno no nível genético. A principal consequência é uma redução expressiva da dimensionalidade do espaço de busca.

Essa mudança exige uma etapa de avaliação mais elaborada. Como o cromossomo não contém a forma completa da solução, é necessário decodificar a atribuição diária e expandi-la para uma representação operacional compatível com o cálculo dos objetivos. A partir dessa expansão, calcula-se a demanda não atendida por clínica, dia e turno e, em seguida, o número de trocas de clínicas por médico ao longo do horizonte. Assim, a compactação altera a forma de representar a solução, mas não altera os objetivos que orientam a busca.

O Algoritmo~\ref{alg:compact-eval} descreve a lógica da avaliação compacta. Primeiro, decodifica-se o cromossomo em uma matriz médico--dia. Depois, essa matriz é expandida para uma representação completa compatível com a disponibilidade e com a lógica operacional do problema. A seguir, calcula-se $F_1$ somando os exames não atendidos e $F_2$ contabilizando as trocas de clínica. Esse algoritmo é o elo entre a representação reduzida e a formulação bi-objetiva original.

\newpage

\begin{algoritmo}
    \caption{Avaliação de um indivíduo no NSGA-2-Compacto}
    \label{alg:compact-eval}
    \begin{algorithmic}[1]
        \Require Cromossomo compacto $y$
        \Ensure Vetor de objetivos $(F_1,F_2)$

        \State Decodificar $y$ em matriz de atribuição diária médico--clínica
        \State Expandir a atribuição para uma representação operacional
        \State Inicializar $F_1 \leftarrow 0$ e $F_2 \leftarrow 0$

        \For{cada unidade $u$, dia $d$ e turno $t$}
        \State Calcular a capacidade total disponível na unidade
        \State Atualizar $F_1$ com os exames não atendidos
        \EndFor

        \For{cada médico $i$}
        \State Percorrer os dias considerando o estado inicial
        \State Incrementar $F_2$ quando houver troca de clínica
        \EndFor

        \State \Return $(F_1,F_2)$
    \end{algorithmic}
\end{algoritmo}

\begin{figure}[htb]
    \centering
    \includegraphics[width=0.80\linewidth]{images/04_nsga-05_gama.png}
    \caption{Variável de estado $\gamma_{iud}$.}
    \label{fig:nsga-gama}
    \fautor
\end{figure}

A interpretação do Algoritmo~\ref{alg:compact-eval} é central para entender a formulação compacta. O ganho computacional vem da redução do cromossomo, mas isso só é útil porque a etapa de avaliação consegue reconstruir adequadamente a solução necessária para medir os objetivos. Em outras palavras, a compactação desloca parte da complexidade da representação para a decodificação e a avaliação.

A lógica global da variante compacta é apresentada no Algoritmo~\ref{alg:nsga2-compact}. Em essência, o método repete a estrutura do NSGA-2, mas agora sobre a representação reduzida. A população inicial é gerada por operadores compactos, os descendentes são produzidos por crossover e mutação também compactos, e os indivíduos podem ainda passar por reparo antes da avaliação. Ao final de cada geração, sobrevivência, dominância e crowding seguem a mesma lógica do NSGA-2 tradicional.

\begin{algoritmo}
    \caption{NSGA-2-Compacto para alocação de médicos}
    \label{alg:nsga2-compact}
    \begin{algorithmic}[1]
        \Require Problema compacto, operadores compactos, parâmetros evolutivos, critério de parada
        \Ensure Fronteira de Pareto aproximada e estatísticas de execução

        \State Construir a representação compacta do problema
        \State Inicializar os operadores compactos
        \State Gerar a população inicial
        \State Avaliar a população inicial por decodificação e expansão
        \State Definir a população corrente $P$

        \While{critério de parada não for satisfeito}
        \State Gerar descendentes por seleção, crossover e mutação compactos
        \If{reparo estiver habilitado}
        \State Aplicar reparo nos descendentes
        \EndIf
        \State Avaliar os descendentes
        \State Unir pais e filhos
        \State Aplicar ordenação não-dominada e \textit{crowding distance}
        \State Atualizar a população sobrevivente
        \If{mecanismos estendidos estiverem habilitados}
        \State Aplicar diversidade adaptativa, reinjeção, arquivo ou busca local
        \EndIf
        \EndWhile

        \State Extrair as soluções não dominadas finais
        \State \Return fronteira aproximada e estatísticas
    \end{algorithmic}
\end{algoritmo}

O Algoritmo~\ref{alg:nsga2-compact} mostra que a principal diferença da variante compacta não está na lógica evolutiva, mas no nível de abstração da decisão codificada. Isso a torna particularmente útil para horizontes de 60 e 90 dias, em que a representação completa tende a se tornar pesada demais. Ao reduzir o tamanho dos indivíduos, o algoritmo consegue aumentar o número efetivo de gerações dentro do mesmo tempo limite e, com isso, explorar melhor a fronteira de Pareto. Assim o NSGA-2-Compacto foi empregado neste estudo como uma resposta ao problema de escalabilidade da formulação base. Sua principal vantagem é reduzir dimensionalidade sem abandonar a estrutura multiobjetivo da busca. Quando combinado aos mecanismos do NSGA-2-Estendido, ele forma a variante mais robusta para instâncias maiores, pois reúne representação mais enxuta, operadores especializados e mecanismos adicionais de diversidade e refinamento.

\chapter{Experimentos e Resultados} % 5
\label{cap:resultados}

Neste capítulo, são apresentados a base de dados e as instâncias utilizadas nos experimentos, a configuração do ambiente computacional, a análise de parametrização dos métodos propostos e os resultados comparativos obtidos. Os experimentos avaliam as três famílias de métodos introduzidas no Capítulo~\ref{cap:mip-heuristicas}: a abordagem exata por $\varepsilon$-restrito, as MIP-heurísticas (Relax-and-Fix e híbrido) e as meta-heurísticas NSGA-II (Estendido e Compacto).

\section{Base de Dados}
\label{cap:resultados:instancias}

O conjunto dispõe de um total de 28.445 registros de exames agendados em 2022, contando com 250 profissionais disponíveis para alocação em 37 unidades distintas. Na prática, o tomador de decisão gera a escala a partir da demanda clínica de até 90 dias.

O trabalho vigente baseia-se em dados reais históricos, disponibilizados pelo corpo gestor da rede de clínicas em estudo. Trata-se de registros dispostos em diferentes bases de dados com informações sobre diferentes aspectos do funcionamento das unidades: o conjunto base armazena os registros de cobertura e atendimentos de cada clínica em todos os períodos, o conjunto de escalas que apresenta a alocação de cada profissional em cada período do horizonte histórico, e o conjunto restrição que informa a qual clínica o médico está associado, podendo, portanto, só atender a esse grupo de unidades. O Quadro \ref{quad:registros-base-dados} descreve os atributos observados nas três bases de dados utilizada para construção das instâncias do modelo.

\begin{quadro}[h]
	\centering
	\caption{Registros e informações disponíveis na base de dados.}
	\begin{footnotesize}
		\begin{tabular}{|c|P{2cm}|M{10cm}|}
			\hline
			\textbf{Arquivo}             & \textbf{Dado} & \textbf{Descrição}                                    \\
			\hline\hline

			\multirow{7}{2cm}{Base}      & Data          & Data do registro (dia e turno)                        \\ \cline{2-3}
			                             & Clínica       & Registro da clínica observada                         \\ \cline{2-3}
			                             & Demanda       & Total de exames agendados                             \\ \cline{2-3}
			                             & Oferecidos    & Total de exames atendidos                             \\ \cline{2-3}
			                             & Salas         & Quantidade de salas disponíveis no período            \\	\cline{2-3}
			                             & Cobertura     & Relação entre Oferecidos e Demanda                    \\ \cline{2-3}
			                             & Buracos       & Salas ociosas no período                              \\ \cline{2-3}
			                             & Ocupação      & Relação de salas utilizadas e salas disponíveis       \\
			\hline \hline
			\multirow{5}{2cm}{Escala}    & Data          & Data do registro (dia e turno)                        \\ \cline{2-3}
			                             & Médico        & Registro do médico observado                          \\ \cline{2-3}
			                             & Situação      & Estado de alocação do médico (se foi ou não escalado) \\ \cline{2-3}
			                             & Clínica       & A clínica para qual foi escalado                      \\ \cline{2-3}
			                             & Ofertados     & Quantos exames o médico realizou no atendimento       \\ \cline{2-3}
			\hline \hline
			\multirow{2}{2cm}{Restrição} & Médico        & Registro do médico observado                          \\ \cline{2-3}
			                             & Clínica       & Unidade que está associado                            \\
			\hline
		\end{tabular}
	\end{footnotesize}
	\label{quad:registros-base-dados}
	\fautor
\end{quadro}

Devido ao volume de dados e grande possibilidade de geração de instâncias, dada pela combinação de diversos recortes e parâmetros, utilizou-se de técnicas de processamento e modelagem de dados, descritas e justificadas na seção \ref{cap:resultados:subsec-pipe}, para automatizar o processo de instanciação do problema.

\subsection{Esteira de Dados e Gerador de Instâncias}
\label{cap:resultados:subsec-pipe}

Para facilitar processamento computacional e uniformizar o processo, optou-se pela criação de uma esteira de dados (em inglês \textit{pipeline}) sequencial para extração de informações relevantes, ilustrada na imagem apresentada pela Figura \ref{fig:diag-pipeline-instancias}. A partir do conjunto de dados histórico inicia-se o processo de captação, que obtêm todos os valores, permitindo operações de transformação (como limpeza, separação, recorte e ratificação dos dados).

\begin{figure}[h]
	\centering
	\includegraphics[width=11cm]{images/Quali - ETL Base de Dados.png}
	\caption{\textit{Pipeline} da geração de instâncias.}
	\label{fig:diag-pipeline-instancias}
	\fautor
\end{figure}

O processo final, descrito por \textit{entrega}, divide as informações em dois grupos que abrangem o mesmo período: dados como a escala registrada foram montados em uma base de dados acessível somente na etapa de validação do modelo, quando há comparação direta entre o resultado do modelo e o registrado. Elementos que compõem a fundação para formulação de componentes do modelo foram movidos para um grupo de base de dados da mesma natureza, facilitando o acesso durante o fluxo de operações computacionais e reduzindo \textit{overhead} (custo adicional de computacional) para o processamento do material.

\section{Análise da Base de Dados}
\label{cap:resultados:sec-analise-base-dados}

A análise descritiva da base histórica é organizada em dois recortes complementares. A Seção~\ref{cap:resultados:base-geral} apresenta uma visão consolidada do conjunto completo (de 02/01/2022 a 31/07/2023) e duas razões \emph{oferta\,$\times$\,demanda} que servem de referência para todo o capítulo. A Seção~\ref{cap:resultados:base-instancias} caracteriza, individualmente, as 10 instâncias de cada um dos quatro horizontes de planejamento (15, 30, 60 e 90 dias) utilizadas nos experimentos. O detalhamento numérico por instância, juntamente com os gráficos descritivos, encontra-se no Apêndice~\ref{apendice-cap:analise-base-dados}.

\subsection{Visão Geral do Conjunto Histórico}
\label{cap:resultados:base-geral}

A base histórica disponibilizada pela rede de clínicas cobre 576 dias consecutivos entre 02/01/2022 e 31/07/2023, dos quais 493 são dias úteis (segunda a sábado) e 83 são domingos --- estes últimos com turno vespertino vazio por característica operacional. Ao todo, são 250 médicos distintos com pelo menos uma janela de disponibilidade registrada e 37 unidades clínicas com demanda observada no período. O Quadro~\ref{quad:base-overall} consolida os totais e médias diárias do conjunto completo, junto às duas razões agregadas \emph{oferta\,$\times$\,demanda} adotadas como referência ao longo do capítulo.

\begin{quadro}[h]
	\centering
	\caption{Indicadores agregados do conjunto histórico completo.}
	\begin{tabular}{|M{5cm}|M{5cm}|}
\hline
\textbf{Indicador} & \textbf{Valor} \\
\hline\hline
Per\'iodo coberto & 02/01/2022 a 31/07/2023 \\
\hline
Total de dias & 576 (\textbf{493} \^{u}teis, \textbf{83} domingos) \\
\hline
M\'edicos distintos & 250 \\
\hline
Unidades distintas & 37 \\
\hline
$\Sigma$ demanda (exames) & 1.478.480 \\
\hline
$\Sigma$ disponibilidade (exames) & 1.496.369 \\
\hline
$\Sigma$ salas $\times$ turnos & 65.045 \\
\hline
M\'edia di\'aria de demanda & 2567 \\
\hline
M\'edia di\'aria de disponibilidade & 2598 \\
\hline
M\'edia di\'aria de salas $\times$ turnos & 112{,}9 \\
\hline
Demanda / (salas $\times$ turnos) & 22{,}73 exames por sala-turno \\
\hline
Disponibilidade / Demanda & 1{,}01 \\
\hline
\end{tabular}

	\label{quad:base-overall}
	\fautor
\end{quadro}

A Figura~\ref{fig:line-demanda-dispo-all} mostra a evolução diária de demanda e disponibilidade ao longo dos 19 meses do conjunto histórico, junto às respectivas médias móveis de 14 dias. Observa-se ciclicidade semanal acentuada (com quedas regulares aos domingos, em razão da ausência do turno vespertino) e crescimento gradual de ambos os indicadores ao longo do período, com ampliação coordenada entre demanda e capacidade. A análise de dispersão entre salas, demanda, disponibilidade e quantidade de médicos, bem como a distribuição não-uniforme de salas entre as 37 unidades, é apresentada no Apêndice~\ref{apendice-cap:analise-base-dados}.

\begin{figure}[h]
	\centering
	\includegraphics[width=12cm]{images/line-disponibilidade-demanda.png}
	\caption{Demanda e disponibilidade diárias entre 02/01/2022 e 31/07/2023, com média móvel de 14 dias.}
	\label{fig:line-demanda-dispo-all}
	\fautor
\end{figure}

\noindent\textbf{Estimativa.}\, As razões \emph{Demanda / (Salas\,$\times$\,Turnos)}~$\approx 22{,}73$ exames por sala-turno e \emph{Disponibilidade / Demanda}~$\approx 1{,}01$ apresentadas no Quadro~\ref{quad:base-overall} devem ser interpretadas como \emph{limitantes ingênuos} da relação oferta\,$\times$\,demanda agregada. Elas \textbf{ignoram} as restrições \ref{eq:r1}--\ref{eq:r5} do modelo (qualificação médico--unidade, unicidade diária, capacidade por turno, cobertura mínima) e \textbf{não correspondem a uma alocação factível}: servem apenas para sinalizar, em ordem de grandeza, que o volume agregado de exames demandados é da mesma magnitude da disponibilidade médica registrada e que a saturação média por sala-turno é compatível com a operação observada.

\subsection{Caracterização das Instâncias}
\label{cap:resultados:base-instancias}

Cada instância utilizada nos experimentos corresponde a um recorte contíguo do conjunto histórico definido pelo seu horizonte (15, 30, 60 ou 90 dias) e pela data de início. Os critérios de seleção das datas de início, bem como a separação entre instâncias de parametrização e de validação, são descritos na Seção~\ref{cap:resultados:instancias-horizontes}. Os Quadros~\ref{quad:base-h15}--\ref{quad:base-h90} resumem, para cada um dos quatro horizontes, as 10 instâncias selecionadas, com os totais agregados de demanda, disponibilidade e capacidade de salas, e as duas razões \emph{oferta\,$\times$\,demanda} introduzidas na Seção~\ref{cap:resultados:base-geral}. As mesmas instâncias são detalhadas, com colunas adicionais (médias diárias, número de unidades ativas, dias úteis), no Apêndice~\ref{apendice-cap:analise-base-dados}.

A leitura conjunta dos Quadros~\ref{quad:base-h15}--\ref{quad:base-h90} evidencia variabilidade não-trivial entre instâncias de um mesmo horizonte: a razão \emph{Demanda / (Salas\,$\times$\,Turnos)} oscila tipicamente entre 20{,}8 e 23{,}5, refletindo períodos com saturação maior ou menor da capacidade de salas, enquanto a razão \emph{Disponibilidade / Demanda} permanece próxima de 1, com leves desvios em direção ao excesso (recortes de 2022) ou ao déficit (alguns recortes de 2023). \textbf{Estas razões mantêm o mesmo caráter de limitante ingênuo} discutido na Seção~\ref{cap:resultados:base-geral}: não substituem a verificação de factibilidade do modelo MIP, mas indicam, em primeira aproximação, o regime esperado de cada instância.

\begin{quadro}[h]
	\centering
	\caption{Caracterização agregada das 10 instâncias do horizonte de 15 dias.}
	\resizebox{0.8\textwidth}{!}{\begin{minipage}{\textwidth}
			\begin{tabular}{|M{2.6cm}|c|c|c|c|c|c|}
\hline
\textbf{Inst\^ancia} & \textbf{M\'ed.} & \textbf{$\Sigma$ Dem.} & \textbf{$\Sigma$ Disp.} & \textbf{$\Sigma$ Salas$\times$T} & \textbf{Dem./Oferta} & \textbf{Disp./Dem.} \\
\hline\hline
02/01/22 a 16/01/22 & 127 & 27.855 & 29.950 & 1.338 & 20{,}82 & 1{,}08 \\
\hline
01/03/22 a 15/03/22 & 136 & 36.397 & 37.374 & 1.575 & 23{,}11 & 1{,}03 \\
\hline
01/05/22 a 15/05/22 & 145 & 36.358 & 36.726 & 1.548 & 23{,}49 & 1{,}01 \\
\hline
01/07/22 a 15/07/22 & 143 & 35.136 & 37.459 & 1.648 & 21{,}32 & 1{,}07 \\
\hline
01/09/22 a 15/09/22 & 139 & 36.467 & 38.108 & 1.691 & 21{,}57 & 1{,}04 \\
\hline
01/11/22 a 15/11/22 & 147 & 36.316 & 37.454 & 1.609 & 22{,}57 & 1{,}03 \\
\hline
01/01/23 a 15/01/23 & 150 & 31.350 & 32.371 & 1.416 & 22{,}14 & 1{,}03 \\
\hline
01/03/23 a 15/03/23 & 162 & 44.210 & 43.319 & 1.920 & 23{,}03 & 0{,}98 \\
\hline
01/05/23 a 15/05/23 & 171 & 43.904 & 43.193 & 1.895 & 23{,}17 & 0{,}98 \\
\hline
01/07/23 a 15/07/23 & 171 & 44.332 & 44.019 & 2.046 & 21{,}67 & 0{,}99 \\
\hline
\end{tabular}

		\end{minipage}}
	\label{quad:base-h15}
	\fautor
\end{quadro}

\begin{quadro}[h]
	\centering
	\caption{Caracterização agregada das 10 instâncias do horizonte de 30 dias.}
	\resizebox{0.8\textwidth}{!}{\begin{minipage}{\textwidth}
			\begin{tabular}{|M{2.6cm}|c|c|c|c|c|c|}
\hline
\textbf{Inst\^ancia} & \textbf{M\'ed.} & \textbf{$\Sigma$ Dem.} & \textbf{$\Sigma$ Disp.} & \textbf{$\Sigma$ Salas$\times$T} & \textbf{Dem./Oferta} & \textbf{Disp./Dem.} \\
\hline\hline
02/01/22 a 31/01/22 & 141 & 57.553 & 61.723 & 2.676 & 21{,}51 & 1{,}07 \\
\hline
01/03/22 a 30/03/22 & 144 & 74.772 & 76.578 & 3.197 & 23{,}39 & 1{,}02 \\
\hline
01/05/22 a 30/05/22 & 154 & 74.067 & 75.659 & 3.218 & 23{,}02 & 1{,}02 \\
\hline
01/07/22 a 30/07/22 & 157 & 70.942 & 74.545 & 3.279 & 21{,}64 & 1{,}05 \\
\hline
01/09/22 a 30/09/22 & 146 & 73.809 & 76.854 & 3.367 & 21{,}92 & 1{,}04 \\
\hline
01/11/22 a 30/11/22 & 154 & 75.006 & 76.229 & 3.266 & 22{,}97 & 1{,}02 \\
\hline
01/01/23 a 30/01/23 & 163 & 70.932 & 72.476 & 3.146 & 22{,}55 & 1{,}02 \\
\hline
01/03/23 a 30/03/23 & 175 & 90.943 & 88.115 & 3.862 & 23{,}55 & 0{,}97 \\
\hline
01/05/23 a 30/05/23 & 178 & 90.052 & 87.432 & 3.830 & 23{,}51 & 0{,}97 \\
\hline
01/07/23 a 30/07/23 & 182 & 88.366 & 87.615 & 3.979 & 22{,}21 & 0{,}99 \\
\hline
\end{tabular}

		\end{minipage}}
	\label{quad:base-h30}
	\fautor
\end{quadro}

\begin{quadro}[h]
	\centering
	\caption{Caracterização agregada das 10 instâncias do horizonte de 60 dias.}
	\resizebox{0.8\textwidth}{!}{\begin{minipage}{\textwidth}
			\begin{tabular}{|M{2.6cm}|c|c|c|c|c|c|}
\hline
\textbf{Inst\^ancia} & \textbf{M\'ed.} & \textbf{$\Sigma$ Dem.} & \textbf{$\Sigma$ Disp.} & \textbf{$\Sigma$ Salas$\times$T} & \textbf{Dem./Oferta} & \textbf{Disp./Dem.} \\
\hline\hline
02/01/22 a 02/03/22 & 146 & 129.617 & 136.118 & 5.784 & 22{,}41 & 1{,}05 \\
\hline
01/02/22 a 01/04/22 & 149 & 147.997 & 151.902 & 6.333 & 23{,}37 & 1{,}03 \\
\hline
01/04/22 a 30/05/22 & 162 & 147.217 & 149.008 & 6.421 & 22{,}93 & 1{,}01 \\
\hline
01/06/22 a 30/07/22 & 169 & 141.915 & 149.923 & 6.571 & 21{,}60 & 1{,}06 \\
\hline
01/08/22 a 29/09/22 & 160 & 152.426 & 158.342 & 6.843 & 22{,}27 & 1{,}04 \\
\hline
01/10/22 a 29/11/22 & 158 & 148.527 & 152.139 & 6.542 & 22{,}70 & 1{,}02 \\
\hline
01/12/22 a 29/01/23 & 175 & 133.009 & 136.052 & 5.931 & 22{,}43 & 1{,}02 \\
\hline
01/02/23 a 01/04/23 & 184 & 174.722 & 171.118 & 7.552 & 23{,}14 & 0{,}98 \\
\hline
01/04/23 a 30/05/23 & 193 & 177.365 & 170.870 & 7.489 & 23{,}68 & 0{,}96 \\
\hline
01/06/23 a 30/07/23 & 193 & 179.589 & 177.056 & 7.905 & 22{,}72 & 0{,}99 \\
\hline
\end{tabular}

		\end{minipage}}
	\label{quad:base-h60}
	\fautor
\end{quadro}

\begin{quadro}[h]
	\centering
	\caption{Caracterização agregada das 10 instâncias do horizonte de 90 dias.}
	\resizebox{0.8\textwidth}{!}{\begin{minipage}{\textwidth}
			\begin{tabular}{|M{2.6cm}|c|c|c|c|c|c|}
\hline
\textbf{Inst\^ancia} & \textbf{M\'ed.} & \textbf{$\Sigma$ Dem.} & \textbf{$\Sigma$ Disp.} & \textbf{$\Sigma$ Salas$\times$T} & \textbf{Dem./Oferta} & \textbf{Disp./Dem.} \\
\hline\hline
02/01/22 a 01/04/22 & 154 & 205.550 & 213.625 & 9.009 & 22{,}82 & 1{,}04 \\
\hline
01/02/22 a 01/05/22 & 162 & 218.818 & 222.837 & 9.439 & 23{,}18 & 1{,}02 \\
\hline
01/04/22 a 29/06/22 & 175 & 218.683 & 224.851 & 9.728 & 22{,}48 & 1{,}03 \\
\hline
01/06/22 a 29/08/22 & 174 & 217.185 & 227.748 & 9.910 & 21{,}92 & 1{,}05 \\
\hline
01/08/22 a 29/10/22 & 162 & 228.215 & 236.405 & 10.217 & 22{,}34 & 1{,}04 \\
\hline
01/09/22 a 29/11/22 & 162 & 222.336 & 228.993 & 9.909 & 22{,}44 & 1{,}03 \\
\hline
01/11/22 a 29/01/23 & 180 & 208.015 & 212.281 & 9.197 & 22{,}62 & 1{,}02 \\
\hline
01/01/23 a 31/03/23 & 189 & 246.456 & 244.418 & 10.724 & 22{,}98 & 0{,}99 \\
\hline
01/03/23 a 29/05/23 & 198 & 267.207 & 257.980 & 11.323 & 23{,}60 & 0{,}97 \\
\hline
01/05/23 a 29/07/23 & 195 & 272.689 & 267.385 & 11.851 & 23{,}01 & 0{,}98 \\
\hline
\end{tabular}

		\end{minipage}}
	\label{quad:base-h90}
	\fautor
\end{quadro}

\clearpage

\section{Configuração Experimental}
\label{cap:resultados:config-experimental}

Os experimentos computacionais foram executados no cluster \textit{Euler} do Instituto de Ciências Matemáticas e de Computação (ICMC-USP), equipado com processadores Intel Xeon E5-2670v3 (12 núcleos a 2,5 GHz) e 128 GB de memória RAM por nó. Cada tarefa foi submetida como um \textit{job} PBS com alocação de 4 núcleos e 16 GB de memória RAM.

O ambiente de software utilizou Python 3.12.8, Gurobi 12.0.3 com licença acadêmica e a biblioteca pymoo 0.6.1 para o NSGA-II. O \textit{solver}
Gurobi foi executado com os parâmetros padrão, com exceção do limite de tempo de 3.600 segundos (1 hora) por ponto da fronteira de Pareto.

\subsection{Instâncias e Horizontes de Planejamento}
\label{cap:resultados:instancias-horizontes}

A escolha dos horizontes de planejamento avaliados neste trabalho segue o ciclo operacional observado na coordenação da rede de clínicas em estudo. O planejamento das escalas médicas é realizado de forma hierárquica e incremental: um horizonte primário de longo prazo é estabelecido para definição de contratos e capacidade, e horizontes sucessivamente mais curtos são utilizados para refinamento à medida que novas informações de demanda e disponibilidade se tornam disponíveis. Essa prática justifica a avaliação dos métodos em quatro horizontes aninhados, cujo significado operacional é resumido a seguir:

\begin{itemize}
	\item \textbf{90 dias (trimestral):} horizonte primário de planejamento. Corresponde à janela na qual a coordenação efetua a contratação e distribuição inicial de capacidade, estabelecendo a escala de referência do trimestre.
	\item \textbf{60 dias (bimestral):} janela de refinamento de médio prazo. É utilizada quando a estimativa de demanda esperada é atualizada com dados mais recentes, permitindo ajustar blocos expressivos da escala sem revisar contratos.
	\item \textbf{30 dias (mensal):} atualização operacional. Aproxima-se do horizonte em que remanejamentos e ajustes finos começam a ser factíveis sem perturbar compromissos contratuais com os profissionais.
	\item \textbf{15 dias (quinzenal):} ciclo mais curto de reprogramação, aplicado a alterações emergenciais de demanda ou disponibilidade e à validação da escala imediatamente antes de sua execução.
\end{itemize}

Além de espelharem a prática operacional, os horizontes escolhidos cobrem granularidades comumente adotadas na literatura de \textit{physician rostering} e impõem níveis crescentes de dificuldade computacional ao modelo MIP, como evidenciado pelos experimentos das seções subsequentes.

Para cada horizonte, foram selecionadas 10 instâncias distribuídas ao longo de 2022 e 2023, divididas em dois conjuntos disjuntos:

\begin{itemize}
	\item \textbf{Parametrização (2 instâncias por horizonte):} conjunto reduzido, utilizado exclusivamente para selecionar os hiperparâmetros de cada família de método (janela das MIP-heurísticas; tamanho de população, probabilidade de mutação e modo de extensão do NSGA-II). As instâncias foram escolhidas em trimestres bem separados do ano, de modo a capturar sazonalidade sem repetição de padrões de demanda.
	\item \textbf{Validação (8 instâncias por horizonte):} conjunto principal, utilizado para comparar os métodos após a seleção de hiperparâmetros. As instâncias cobrem todos os trimestres de 2022 e 2023, com distribuição balanceada ao longo do calendário para mitigar vieses sazonais. Os resultados discutidos na Seção~\ref{cap:resultados:comparacao} referem-se exclusivamente a este conjunto.
\end{itemize}

A separação em dois conjuntos disjuntos garante que o desempenho reportado na fase de validação não esteja condicionado aos mesmos dados usados para ajustar os hiperparâmetros, reduzindo o risco de sobreajuste aos perfis de demanda específicos usados na parametrização. Ao todo, são $4 \times (2 + 8) = 40$ pares (horizonte, instância) utilizados nos experimentos.

O Quadro~\ref{quad:instancias-experimentos} resume a estrutura dos experimentos, e o Quadro~\ref{quad:instancias-datas} cataloga os recortes específicos utilizados em cada conjunto.

\begin{quadro}[h]
	\centering
	\caption{Estrutura dos experimentos por horizonte de planejamento.}
	\resizebox{0.75\textwidth}{!}{\begin{minipage}{\textwidth}
			\begin{tabular}{|c|M{3.3cm}|M{2.5cm}|M{1.8cm}|M{6cm}|}
				\hline
				\textbf{Horizonte} & \textbf{Instâncias (Parametrização)} & \textbf{Instâncias (Validação)} & \textbf{Pontos $\varepsilon$} & \textbf{Métodos avaliados}      \\
				\hline
				15 dias            & 2                                    & 8                               & 7                             & RF, Híbrido, NSGA-IIe, NSGA-IIc \\
				\hline
				30 dias            & 2                                    & 8                               & 7                             & RF, Híbrido, NSGA-IIe, NSGA-IIc \\
				\hline
				60 dias            & 2                                    & 8                               & 7                             & RF, Híbrido, NSGA-IIe, NSGA-IIc \\
				\hline
				90 dias            & 2                                    & 8                               & 7                             & RF, Híbrido, NSGA-IIe, NSGA-IIc \\
				\hline
			\end{tabular}
		\end{minipage}}

	\label{quad:instancias-experimentos}
	\fautor
\end{quadro}

\begin{quadro}[h]
	\centering
	\caption{Catálogo de instâncias por horizonte e conjunto.}
	\resizebox{0.8\textwidth}{!}{\begin{minipage}{\textwidth}
			\begin{tabular}{|M{1.7cm}|M{4.7cm}|M{10cm}|}
				\hline
				\textbf{Horizonte} & \textbf{Parametrização \newline (2 instâncias)} & \textbf{Validação (8 instâncias)}                                                                     \\
				\hline
				15 dias            & 02/01/2022,\ 01/07/2022                & 01/03/2022,\ 01/05/2022,\ 01/09/2022,\ 01/11/2022,\ 01/01/2023,\ 01/03/2023,\ 01/05/2023,\ 01/07/2023 \\
				\hline
				30 dias            & 02/01/2022,\ 01/07/2022                & 01/03/2022,\ 01/05/2022,\ 01/09/2022,\ 01/11/2022,\ 01/01/2023,\ 01/03/2023,\ 01/05/2023,\ 01/07/2023 \\
				\hline
				60 dias            & 01/04/2022,\ 01/08/2022                & 02/01/2022,\ 01/02/2022,\ 01/06/2022,\ 01/10/2022,\ 01/12/2022,\ 01/02/2023,\ 01/04/2023,\ 01/06/2023 \\
				\hline
				90 dias            & 01/04/2022,\ 01/08/2022                & 02/01/2022,\ 01/02/2022,\ 01/06/2022,\ 01/09/2022,\ 01/11/2022,\ 01/01/2023,\ 01/03/2023,\ 01/05/2023 \\
				\hline
			\end{tabular}
		\end{minipage}}
	\label{quad:instancias-datas}
	\fautor
\end{quadro}

\subsection{Métricas de Avaliação}
\label{cap:resultados:metricas}

Devido à natureza multiobjetivo do problema, a avaliação dos métodos utiliza métricas que capturam tanto a qualidade da fronteira de Pareto aproximada quanto o desempenho computacional. O Quadro \ref{quad:metricas-resultados} descreve os indicadores empregados.

\begin{quadro}[h]
	\centering
	\caption{Métricas de avaliação dos experimentos.}
	\resizebox{0.9\textwidth}{!}{\begin{minipage}{\textwidth}
			\centering
			\begin{tabular}{|c|M{11.5cm}|}
				\hline
				\textbf{Métrica}      & \textbf{Descrição}                                                                                                                               \\
				\hline
				Hipervolume (HV)      & Volume do espaço dominado pela fronteira aproximada em relação a um ponto de referência. Quanto maior, melhor a qualidade da fronteira.          \\ \hline
				Amplitude $Z_1$       & Diferença entre o maior e o menor valor de exames não atendidos na fronteira ($Z_1^{\max} - Z_1^{\min}$). Mede a cobertura do eixo de cobertura. \\ \hline
				Amplitude $Z_2$       & Diferença entre o maior e o menor valor de trocas na fronteira ($Z_2^{\max} - Z_2^{\min}$). Mede a cobertura do eixo de estabilidade.            \\ \hline
				Gap de referência     & Distância relativa entre a fronteira heurística e a fronteira exata, calculada ponto a ponto.                                                    \\ \hline
				MIP Gap               & Desvio percentual entre o valor objetivo e o melhor limitante inferior reportado pelo Gurobi.                                                    \\ \hline
				Tempo de execução (s) & Tempo de parede para obtenção de todos os pontos da fronteira.                                                                                   \\ \hline
				Taxa de factibilidade & Proporção de soluções NSGA-II verificadas como factíveis pelo modelo Gurobi.                                                                     \\ \hline
			\end{tabular}
		\end{minipage}}
	\label{quad:metricas-resultados}
	\fautor
\end{quadro}

\section{Resultados do Método Exato}
\label{cap:resultados:exato}
Esta seção caracteriza, em isolado, o desempenho do método exato por $\varepsilon$-restrito. O método não possui hiperparâmetros sujeitos a sintonia experimental --- depende apenas das definições padrão do Gurobi 12.0.3 e do limite de tempo de 3.600~s por ponto $\varepsilon$ (Seção~\ref{cap:resultados:config-experimental}) ---, de modo que os resultados a seguir utilizam o conjunto completo das 10 instâncias por horizonte (40 pares horizonte-instância), e não apenas as 8 de validação. As Tabelas~\ref{tab:exato-mipgap} e~\ref{tab:exato-runtime} concentram os indicadores agregados que sustentam a leitura comparativa da Seção~\ref{cap:resultados:comparacao}; o detalhamento por instância e por ponto $\varepsilon$, incluindo fronteiras de Pareto sobrepostas, mapas de calor do gap e diagramas de caixa de tempo, é apresentado no Apêndice~\ref{apendice-cap:exato-detalhe}.

\subsection{Convergência e Gap MIP}
\label{cap:resultados:exato-gap}

A Tabela~\ref{tab:exato-mipgap} reporta o \textit{MIP gap} médio ao final do limite de tempo de 3.600~s por ponto $\varepsilon$, definido como $\text{MIP	gap} = (z_{\text{best}} - z_{\text{lb}})/|z_{\text{best}}|$, e o número médio de pontos da grade resolvidos até a otimalidade comprovada por instância.

\begin{table}[h]
	\centering
	\caption{MIP gap médio do método exato por ponto $\varepsilon$.}
	\begin{tabular}{M{3cm}M{3cm}M{3cm}}
		\toprule
		\textbf{Horizonte} & \textbf{MIP gap médio (\%)} & \textbf{Pontos $\varepsilon$ ótimos} \\
		\midrule
		15 dias            & 4{,}10 $\pm$ 8{,}32         & 3 / 7                                \\
		30 dias            & 42{,}66 $\pm$ 13{,}69       & 0 / 7                                \\
		60 dias            & 78{,}11 $\pm$ 14{,}73       & 0 / 7                                \\
		90 dias            & 96{,}80 $\pm$ 6{,}22        & 0 / 7                                \\
		\bottomrule
	\end{tabular}
	\label{tab:exato-mipgap}
	\fautor
\end{table}

Apenas o horizonte quinzenal atinge otimalidade comprovada em uma fração relevante dos pontos $\varepsilon$ (3 dos 7, em média); a partir de 30 dias o crescimento combinatório das variáveis $x_{iudt}$ (ordem de $10^6$ em 90 dias) impede a convergência em qualquer ponto da grade. A elevada dispersão observada em 15 dias ($\pm 8{,}32$~pp) é dominada por duas instâncias outlier discutidas no Apêndice~\ref{apendice-cap:exato-detalhe}.

\subsection{Tempo de Execução}
\label{cap:resultados:exato-runtime}

A Tabela~\ref{tab:exato-runtime} consolida o tempo médio por ponto $\varepsilon$, o número médio de pontos no limite de tempo e o tempo total da fronteira completa.

\begin{table}[h]
	\centering
	\caption{Tempo de execução do método exato por ponto $\varepsilon$ e tempo total da fronteira.}
	\begin{tabular}{M{2.5cm}M{3.5cm}M{3cm}M{3.5cm}}
		\toprule
		\textbf{Horizonte} & \textbf{Tempo médio por ponto $\varepsilon$ (s)} & \textbf{Pontos no limite de tempo} & \textbf{Tempo total da fronteira (s)} \\
		\midrule
		15 dias            & 2.702 $\pm$ 1.192                                & 3{,}9 / 7                          & 18.911 $\pm$ 4.788                    \\
		30 dias            & 3.600 $\pm$ 0                                    & 7{,}0 / 7                          & 25.201 $\pm$ 0                        \\
		60 dias            & 3.603 $\pm$ 10                                   & 7{,}0 / 7                          & 25.219 $\pm$ 27                       \\
		90 dias            & 3.605 $\pm$ 14                                   & 7{,}0 / 7                          & 25.233 $\pm$ 41                       \\
		\bottomrule
	\end{tabular}
	\label{tab:exato-runtime}
	\fautor
\end{table}

O tempo total se aproxima do produto $7 \times 3.600 = 25.200$~s a partir de 30 dias, confirmando que todos os pontos $\varepsilon$ saturam o limite de tempo sem convergência. Em 15 dias, o tempo total inferior ($\sim 18.911$~s) reflete a finalização antecipada dos pontos resolvidos à otimalidade.


\subsection{Implicações para a Análise Comparativa}
\label{cap:resultados:exato-implicacoes}

A caracterização do método exato motiva duas leituras distintas das métricas comparativas apresentadas na Seção~\ref{cap:resultados:comparacao}. No horizonte de 15 dias, o exato fornece uma fronteira de Pareto verdadeira (ou quase verdadeira), de modo que o hipervolume e o \textit{gap} de referência das demais famílias podem ser interpretados como medidas absolutas de qualidade. Nos horizontes de 30, 60 e 90 dias, a fronteira do exato é, ela própria, sub-ótima, e o hipervolume relativo passa a representar um indicador de proximidade à melhor fronteira conhecida no mesmo limite de tempo, e não à fronteira ótima global.

Adicionalmente, a transição de regime entre os horizontes de 15 e 30 dias estabelece, para o problema em estudo, o ponto a partir do qual o uso de heurísticas torna-se necessário em ambientes operacionais com restrição de tempo. Para o tomador de decisão da rede de clínicas, isso significa que o ciclo quinzenal de reprogramação pode ser tratado por método exato, enquanto os horizontes maiores (mensal, bimestral e trimestral) requerem necessariamente abordagens aproximadas, justificando a investigação das MIP-heurísticas e das meta-heurísticas conduzida nas seções subsequentes.

\section{Parametrização}
\label{cap:resultados:parametrizacao}

Antes da execução dos experimentos de validação, foi conduzida uma fase de parametrização para definir as configurações utilizadas em cada família de método. As três famílias avaliadas neste trabalho recebem tratamentos distintos nesta fase: o método exato por $\varepsilon$-restrito não possui hiperparâmetros próprios sujeitos a sintonia, atuando como referência para os demais; as MIP-heurísticas (Relax-and-Fix e Híbrido) requerem a seleção do tamanho da janela temporal e da política de sobreposição; e as meta-heurísticas NSGA-II (Estendido e Compacto) requerem a definição de tamanho de população, probabilidade de mutação e modo de extensão. Conforme descrito na Seção~\ref{cap:resultados:instancias-horizontes}, a fase utilizou as 2 instâncias de parametrização por horizonte (8 instâncias ao todo), representativas da variabilidade sazonal e de demanda dos dados reais. O objetivo é identificar a configuração que produz a melhor qualidade de fronteira de Pareto dentro do limite de tempo disponível, para posterior comparação com o método exato na fase de validação (Seção~\ref{cap:resultados:comparacao}).

\subsection{Método Exato (referência)}
\label{cap:resultados:param-exato}

O método exato por $\varepsilon$-restrito é a única família para a qual não há configuração de hiperparâmetros sujeita a investigação experimental: o seu funcionamento depende apenas das definições padrão do \textit{solver} Gurobi 12.0.3 e do limite de tempo de 3.600 segundos por ponto da fronteira, descritos na Seção~\ref{cap:resultados:config-experimental}. A grade de valores de $\varepsilon$ utilizada para enumerar os pontos da fronteira é composta por sete valores fixos por instância, obtidos por otimização lexicográfica que determina os limitantes mínimo e máximo do número de trocas $Z_2$ na instância considerada. Essa mesma grade é compartilhada com todas as MIP-heurísticas, garantindo comparabilidade ponto a ponto entre as fronteiras geradas pelos diferentes métodos sobre uma mesma instância.

Por consequência, os resultados do método exato cumprem dois papéis distintos na análise comparativa apresentada na
Seção~\ref{cap:resultados:comparacao}. Em horizontes nos quais o Gurobi atinge otimalidade, o exato fornece a fronteira de Pareto verdadeira, contra a qual a qualidade dos métodos heurísticos pode ser medida em termos absolutos. Em horizontes nos quais o limite de tempo é insuficiente para convergência, o exato fornece uma fronteira sub-ótima cuja qualidade decresce com o horizonte, fornecendo um limitante superior para a qualidade alcançável pelos demais métodos no mesmo limite de tempo. Esta distinção é central para a interpretação das métricas de hipervolume e gap de referência apresentadas adiante.

\subsection{MIP-Heurísticas}
\label{cap:resultados:param-mip}

Para as MIP-heurísticas, o parâmetro principal é o tamanho da janela temporal de decomposição, expresso como fração do horizonte total. Foram avaliadas seis variantes do Relax-and-Fix (RF) e seis do método híbrido RF+Fix-and-Optimize (FO), combinando tamanhos de janela de 25\%, 50\% e 75\% do horizonte, com e sem sobreposição entre janelas consecutivas (0 ou 2 períodos). A Figura \ref{fig:param-mip-barchart} apresenta o gap de referência médio em relação ao método exato para cada configuração, agregando as instâncias de parametrização no horizonte de 15 dias, em que a comparação com o método exato é mais confiável (no horizonte de 15 dias o Gurobi atinge otimalidade ou \textit{gaps} inferiores a 1\%).

\begin{figure}[h]
	\centering
	\includegraphics[width=\textwidth]{images/param_mip_barchart.png}
	\caption{Gap de referência médio por configuração de janela, horizonte de 15 dias.}
	\label{fig:param-mip-barchart}
	\fautor
\end{figure}

Com base nos resultados de parametrização, a configuração \textbf{RF com janela de 75\% e sobreposição de 2 períodos} (\texttt{rfw75\_rfo2}) foi selecionada para o Relax-and-Fix (gap médio de 3{,}0\%) e a configuração \textbf{Híbrido com janela de 50\% e sobreposição de 2 períodos} (\texttt{rfw50\_rfo2\_fow50}) foi selecionada para o método híbrido (gap médio de 2{,}9\%), equilibrando qualidade de solução e tempo de execução. Janelas maiores reduzem a fragmentação do problema, permitindo ao solver capturar dependências temporais mais longas em cada subproblema; a sobreposição de dois períodos entre janelas consecutivas suaviza as transições, evitando descontinuidades na fixação das variáveis $w_{iud}$.
\subsection{Meta-Heurísticas (NSGA-II)}
\label{cap:resultados:param-nsga2}

Para o NSGA-II, foram avaliadas combinações de tamanho de população ($N \in \{50, 100, 200\}$), probabilidade de mutação ($p_m \in \{0{,}01, 0{,}03,0{,}05\}$) e modo de melhoria evolutiva: com todas as extensões habilitadas (\textit{allon}), sem busca local (\textit{nols}) e sem nenhuma extensão (\textit{noenhance}). As avaliações foram conduzidas separadamente para o NSGA-II Estendido (representação binária) e para o NSGA-II Compacto
(representação inteira).

A métrica primária de avaliação é o hipervolume (HV) da fronteira aproximada, normalizado pelo HV do método exato na mesma instância. A Figura
\ref{fig:param-nsga-heatmap} apresenta os mapas de calor do HV relativo por tamanho de população e probabilidade de mutação, separados por modo de
extensão, para cada variante do NSGA-II.

\begin{figure}[h]
	\centering
	\includegraphics[width=\textwidth]{images/param_nsga2_heatmap.png}
	\caption{Mapa de calor do hipervolume relativo ao método exato, por configuração de população e mutação (NSGA-II Compacto na linha superior, Estendido na linha inferior; colunas: \textit{allon}, \textit{nols}, \textit{noenhance}).}
	\label{fig:param-nsga-heatmap}
	\fautor
\end{figure}

A análise dos mapas evidencia que o modo \textit{allon} supera consistentemente os modos \textit{nols} e \textit{noenhance}, em ambas as representações. A representação compacta apresenta HV relativo entre 0{,}67 e 0{,}72 no modo \textit{allon}, enquanto a representação estendida atinge valores de 0{,}47 a 0{,}51, refletindo o maior número de gerações viáveis dentro do limite de tempo quando a dimensionalidade do cromossomo é reduzida. A variação em relação ao tamanho da população e probabilidade de mutação é relativamente pequena (inferior a 5\% em termos absolutos).

Com base nos resultados da parametrização, a configuração $N=50$, $p_m=0{,}03$, modo \textit{allon} foi selecionada para o \textbf{NSGA-II Estendido} (HV relativo médio 0{,}510) e a configuração $N=50$, $p_m=0{,}05$, modo \textit{allon} foi selecionada para o \textbf{NSGA-II Compacto} (HV relativo médio 0{,}715).

\section{Análise Comparativa}
\label{cap:resultados:comparacao}

Nesta seção, são apresentados os resultados dos experimentos de validação entre os cinco métodos avaliados: método exato por $\varepsilon$-restrito (\textbf{Exato}), Relax-and-Fix (\textbf{RF}), método híbrido RF+FO (\textbf{Híbrido}), NSGA-II Estendido (\textbf{NSGA-IIe}) e NSGA-II Compacto (\textbf{NSGA-IIc}). Os experimentos foram conduzidos nas 8 instâncias de validação por horizonte de planejamento (15, 30, 60, 90 dias) catalogadas no Quadro~\ref{quad:instancias-datas}, totalizando 32 pares instância-método por método. As configurações de hiperparâmetros das MIP-heurísticas e do NSGA-II são aquelas selecionadas na fase de parametrização (Seção~\ref{cap:resultados:parametrizacao}); o método exato é executado em sua configuração padrão descrita na Seção~\ref{cap:resultados:param-exato}, atuando como referência para a avaliação dos demais métodos. Todas as figuras e tabelas comparativas a seguir incluem o método exato como linha de base, exceto pelas figuras de parametrização (Figuras~\ref{fig:param-mip-barchart}~e~\ref{fig:param-nsga-heatmap}), nas quais ele aparece implicitamente como referência para o cálculo dos \textit{gaps} e do hipervolume relativo.

\subsection{Fronteiras de Pareto}
\label{cap:resultados:pareto}

A Figura \ref{fig:pareto-todos} apresenta uma visão comparativa das fronteiras de Pareto aproximadas obtidas pelos cinco métodos em uma instância representativa de cada horizonte de planejamento. O eixo horizontal representa o número total de exames não atendidos ($Z_1$) e o eixo vertical o número total de trocas de clínica ($Z_2$), ambos com orientação de minimização.

\begin{figure}[h]
	\centering
	\includegraphics[width=0.9\textwidth]{images/pareto_comparativo_todos_horizontes.png}
	\caption{Fronteiras de Pareto para instâncias representativas em cada horizonte.}
	\label{fig:pareto-todos}
	\fautor
\end{figure}

Observa-se que, no horizonte de 15 dias, as cinco curvas se aproximam da fronteira ótima determinada pelo método exato, com maior dispersão para as variantes do NSGA-II no trecho de $Z_1$ elevado. Em horizontes maiores (60 e 90 dias), o método exato não atinge otimalidade no tempo disponível, o que resulta em uma fronteira sub-ótima e, em consequência, em pontos heurísticos que ocasionalmente dominam pontos do exato.

\subsection{Qualidade das Soluções}
\label{cap:resultados:qualidade}

A Tabela \ref{tab:hipervolume} apresenta o hipervolume médio das fronteiras de Pareto obtidas pelos cinco métodos, para cada horizonte de planejamento. O ponto de referência utilizado no cálculo do hipervolume é o vetor de piores valores observados na fronteira do método exato para cada instância, garantindo comparabilidade entre métodos. Os valores foram normalizados pelo hipervolume do método exato (referência $= 1{,}000$) para facilitar a interpretação. Valores superiores a 1{,}000 ocorrem em horizontes em que o Gurobi não atinge otimalidade no tempo limite, de modo que a fronteira de referência fica sub-ótima e pode ser dominada pontualmente pelos métodos heurísticos.


\begin{table}[h]
	\centering
	\caption{Hipervolume relativo ao método exato.}
	\resizebox{0.9\textwidth}{!}{\begin{minipage}{\textwidth}
			\centering
			\begin{tabular}{ccccc}
				\toprule
				\textbf{Método} & \textbf{15 dias}             & \textbf{30 dias}             & \textbf{60 dias}             & \textbf{90 dias}                 \\
				\midrule
				Exato           & 1{,}000\newline$\pm$ 0{,}000 & 1{,}000\newline$\pm$ 0{,}000 & 1{,}000\newline$\pm$ 0{,}000 & 1{,}000\newline$\pm$ 0{,}000     \\
				RF              & 0{,}984\newline$\pm$ 0{,}017 & 0{,}852\newline$\pm$ 0{,}028 & 1{,}158\newline$\pm$ 0{,}093 & 4{,}213\newline$\pm$ 5{,}158     \\
				Híbrido         & 0{,}996\newline$\pm$ 0{,}010 & 1{,}008\newline$\pm$ 0{,}039 & 0{,}938\newline$\pm$ 0{,}139 & 1{,}941\newline$\pm$ 3{,}233     \\
				NSGA-IIe        & 0{,}279\newline$\pm$ 0{,}121 & 0{,}023\newline$\pm$ 0{,}050 & 0{,}736\newline$\pm$ 0{,}172 & 10{,}190\newline$\pm$ 21{,}618   \\
				NSGA-IIc        & 0{,}382\newline$\pm$ 0{,}118 & 0{,}217\newline$\pm$ 0{,}156 & 1{,}039\newline$\pm$ 0{,}227 & 12{,}888 \newline $\pm$ 26{,}976 \\
				\bottomrule
			\end{tabular}
		\end{minipage}}
	\label{tab:hipervolume}
	\fautor
\end{table}

Os resultados indicam que o método Híbrido mantém HV próximo ao do exato nos horizontes de 15 e 30 dias, enquanto o NSGA-II Compacto supera o Estendido em todos os horizontes em razão de sua representação reduzida, que permite maior número de gerações dentro do limite de tempo. Para horizontes longos, a alta variância dos valores (especialmente em 90 dias) reflete a instabilidade da referência construída a partir do exato não convergido.

A Tabela \ref{tab:gap-referencia} apresenta o gap de referência médio dos métodos MIP-heurísticos em relação ao melhor limitante inferior reportado pelo método exato para o mesmo ponto $\varepsilon$. O gap de referência $\Delta_{\text{ref}}$ é calculado como $\Delta_{\text{ref}} = |f(x_h) - \underline{f}^\ast| / |f(x_h)| \cdot 100$, onde $f(x_h)$ é o valor objetivo da solução heurística e $\underline{f}^\ast$ o melhor limitante reportado pelo Gurobi.

\begin{table}[h]
	\centering
	\caption{Gap de referência médio dos métodos MIP-heurísticos em relação ao exato.}
	\resizebox{0.9\textwidth}{!}{\begin{minipage}{\textwidth}
			\centering
			\begin{tabular}{M{2.8cm}M{2.9cm}M{2.9cm}M{2.9cm}M{3cm}}
				\toprule
				\textbf{Método} & \textbf{15 dias}  & \textbf{30 dias}   & \textbf{60 dias}   & \textbf{90 dias}   \\
				\midrule
				RF              & 6{,}4 $\pm$ 7{,}3 & 52{,}2 $\pm$ 7{,}8 & 80{,}2 $\pm$ 8{,}3 & 97{,}1 $\pm$ 2{,}6 \\
				Híbrido         & 6{,}7 $\pm$ 8{,}6 & 44{,}2 $\pm$ 7{,}6 & 81{,}0 $\pm$ 5{,}9 & 97{,}0 $\pm$ 3{,}3 \\
				\bottomrule
			\end{tabular}
		\end{minipage}}
	\label{tab:gap-referencia}
	\fautor
\end{table}


Os valores elevados de gap para horizontes $\geq 30$ dias refletem a qualidade do limitante inferior produzido pelo Gurobi na janela de 1~hora por ponto $\varepsilon$, e não necessariamente a qualidade intrínseca das soluções heurísticas. Como o método exato também atinge o limite de tempo sem convergir (Tabela~\ref{tab:runtime}), o gap reportado é um limitante superior da distância ao ótimo global. Para o horizonte de 15 dias, em que o exato encontra soluções ótimas ou com gap $< 1\%$, os gaps médios de 6--7\% indicam boa qualidade relativa das heurísticas.

A aparente superioridade do Relax-and-Fix sobre o método Híbrido no horizonte de 60 dias (HV relativo de 1{,}158 vs.\ 0{,}938, Tabela~\ref{tab:hipervolume}) decorre da divisão de tempo adotada no modo híbrido: a fase construtiva RF ocupa 25\% do limite de tempo por ponto $\varepsilon$ ($\approx 900$~s), e a fase de refinamento Fix-and-Optimize (FO) ocupa os 75\% restantes ($\approx 2.700$~s). No horizonte de 60 dias, os 900~s disponíveis para a fase RF frequentemente são insuficientes para construir uma solução incumbente de qualidade, tornando o ponto de partida da fase FO fraco ou inviável. Em contraste, o RF puro dispõe de 3.600~s para a decomposição completa, o que explica a inversão de desempenho neste horizonte.

As soluções geradas pelo NSGA-II foram verificadas quanto à factibilidade mediante resolução do modelo Gurobi com as variáveis fixadas nos valores propostos pela meta-heurística. A Tabela \ref{tab:factibilidade-nsga2} apresenta a taxa de factibilidade global (satisfação de todas as restrições \ref{eq:r1}--\ref{eq:r6}) e a taxa de violação da restrição de cobertura mínima (restrições \ref{eq:r5}) para cada horizonte.

\begin{table}[h]
	\centering
	\caption{Factibilidade das soluções NSGA-II: global $vs.$ violação da restrição de cobertura mínima.}
	\resizebox{0.9\textwidth}{!}{\begin{minipage}{\textwidth}
			\centering
			\begin{tabular}{M{2.8cm}M{2.9cm}M{2.9cm}M{2.9cm}M{3cm}}
				\toprule
				\textbf{Método} & \textbf{15 dias}    & \textbf{30 dias}    & \textbf{60 dias}    & \textbf{90 dias}    \\
				\midrule
				NSGA-IIe        & 0{,}00 $\pm$ 0{,}00 & 0{,}00 $\pm$ 0{,}00 & 0{,}00 $\pm$ 0{,}00 & 0{,}00 $\pm$ 0{,}00 \\
				NSGA-IIc        & 0{,}00 $\pm$ 0{,}00 & 0{,}00 $\pm$ 0{,}00 & 0{,}00 $\pm$ 0{,}00 & 0{,}00 $\pm$ 0{,}00 \\
				\bottomrule
			\end{tabular}
		\end{minipage}}
	\label{tab:factibilidade-nsga2}
	\fautor
\end{table}

A taxa de factibilidade global nula indica que todas as soluções NSGA-II violam ao menos a restrição \ref{eq:r5} (cobertura mínima). A análise por variante revela, porém, distinções relevantes: no NSGA-II Estendido, as restrições estruturais \ref{eq:r1}--\ref{eq:r4} são satisfeitas na totalidade das soluções, e os valores de $Z_1$ e $Z_2$ foram verificados por resolução Gurobi com variáveis fixadas, sem discrepância; no NSGA-II Compacto, 89{,}8\% das soluções satisfazem \ref{eq:r1}--\ref{eq:r4}, enquanto 10{,}2\% apresentam violação da restrição de capacidade de sala (\ref{eq:r4}), concentradas na clínica U06, turno vespertino, indicando uma limitação do decodificador de cromossomo compacto nesta configuração específica. Em ambas as variantes, a restrição \ref{eq:r5} é violada em 4--6\% das combinações clínica--período, com concentração no turno vespertino.

\noindent\textbf{Nota metodológica.} Os resultados de hipervolume e fronteira de Pareto do NSGA-II (Tabelas~\ref{tab:hipervolume} e~\ref{tab:factibilidade-nsga2}) são apresentados em paralelo aos métodos MIP exclusivamente para contextualização. Dado que nenhuma solução NSGA-II satisfaz \ref{eq:r5}, e uma fração das soluções compactas também viola \ref{eq:r4}, esses resultados não constituem comparação direta com métodos que garantem factibilidade. A análise evolutiva deve ser lida isoladamente, com ênfase na qualidade das soluções dentro da família de meta-heurísticas.

\subsection{Desempenho Computacional}
\label{cap:resultados:desempenho}

A Tabela \ref{tab:runtime} apresenta o tempo médio total de execução de cada método por horizonte. Para os métodos MIP, o tempo inclui a resolução de todos os 7 pontos $\varepsilon$ da fronteira, cada um limitado a 3.600 segundos. Para o NSGA-II, o tempo total corresponde ao limite de 3.600 segundos por instância, já que uma única execução produz a fronteira aproximada completa.

\begin{table}[!htb]
	\centering
	\caption{Tempo total médio de execução por método e horizonte.}
	\resizebox{0.95\textwidth}{!}{\begin{minipage}{\textwidth}
			\centering
			\begin{tabular}{M{2.5cm}M{2.75cm}M{2.75cm}M{2.75cm}M{2.75cm}}
				\toprule
				\textbf{Método} & \textbf{15 dias}   & \textbf{30 dias} & \textbf{60 dias}   & \textbf{90 dias} \\
				\midrule
				Exato           & 20.124 $\pm$ 3.862 & 25.201 $\pm$ 0   & 25.222 $\pm$ 29    & 25.233 $\pm$ 44  \\
				RF              & 15.667 $\pm$ 1.905 & 14.879 $\pm$ 999 & 13.112 $\pm$ 1.367 & 11.554 $\pm$ 36  \\
				Híbrido         & 20.509 $\pm$ 3.281 & 25.201 $\pm$ 0   & 25.227 $\pm$ 34    & 25.231 $\pm$ 34  \\
				NSGA-IIe        & 3.389 $\pm$ 121    & 3.230 $\pm$ 216  & 2.861 $\pm$ 399    & 2.109 $\pm$ 200  \\
				NSGA-IIc        & 3.569 $\pm$ 24     & 3.537 $\pm$ 24   & 3.466 $\pm$ 61     & 3.407 $\pm$ 70   \\
				\bottomrule
			\end{tabular}
		\end{minipage}}
	\label{tab:runtime}
	\fautor
\end{table}

Os tempos dos métodos MIP (Exato, RF, Híbrido) acompanham a soma dos 7 limites de 3.600 segundos (25.200 s) nos horizontes $\geq 30$ dias, indicando que o \textit{solver} atinge o limite de tempo em todos os pontos $\varepsilon$ sem convergência. O Relax-and-Fix apresenta o menor tempo total por reduzir a dimensão do subproblema em cada janela. As variantes do NSGA-II, por construção, são limitadas a 3.600 segundos por instância, representando uma redução de aproximadamente 7 vezes no tempo total em relação aos métodos MIP.

A Figura \ref{fig:escalabilidade} ilustra a escalabilidade dos métodos em função do horizonte de planejamento. Observa-se que os métodos MIP permanecem em regime de tempo limite, enquanto o NSGA-II mantém tempo aproximadamente constante pelo critério de parada por tempo.

\begin{figure}[h]
	\centering
	\includegraphics[width=12cm]{images/escalabilidade_metodos.png}
	\caption{Escalabilidade dos métodos em função do horizonte de planejamento.}
	\label{fig:escalabilidade}
	\fautor
\end{figure}

\subsection{Discussão dos Resultados}
\label{cap:resultados:discussao}

A análise conjunta dos resultados permite extrair quatro observações principais sobre o problema e os métodos avaliados.

Em primeiro lugar, a \textbf{dificuldade computacional do modelo exato} cresce acentuadamente com o horizonte de planejamento. Para 15 dias, o Gurobi atinge otimalidade ou \textit{gap} inferior a 1\% em todos os pontos $\varepsilon$; a partir de 30 dias, o \textit{solver} atinge o limite de tempo em cada ponto sem convergência, como evidenciado pelos tempos totais próximos a 25.200 s (Tabela~\ref{tab:runtime}) e pelos \textit{gaps} de referência acima de 40\% (Tabela~\ref{tab:gap-referencia}). Isso motiva diretamente o uso de heurísticas em contextos operacionais reais, nos quais o horizonte típico é de 60 ou 90 dias.

Em segundo lugar, entre as \textbf{MIP-heurísticas}, o método Híbrido (Relax-and-Fix seguido de Fix-and-Optimize) supera consistentemente o Relax-and-Fix isolado em qualidade de hipervolume nos horizontes curtos, embora ao custo de tempo total similar ao do exato. O Relax-and-Fix puro, por sua vez, é aproximadamente $30$--$55\%$ mais rápido e mantém HV relativo de 0{,}85--0{,}98 nos horizontes de 15 e 30 dias, constituindo uma alternativa viável quando o tempo é restritivo.

Em terceiro lugar, as variantes do \textbf{NSGA-II} apresentam um \textit{trade-off} claro entre representação e número de gerações efetivas. A representação compacta ($w_{id}$) permite ao NSGA-IIc executar centenas de gerações dentro dos 3.600 segundos, produzindo HV relativo superior ao da representação binária ($x_{iudt}$) em todos os horizontes. No horizonte de 60 dias, o NSGA-IIc apresenta HV relativo de 1{,}039 --- valor superior à referência exata. Essa aparente vantagem, contudo, não deve ser interpretada como superioridade da meta-heurística: o método exato neste horizonte apresenta MIP gap médio de 78{,}1\% (Tabela~\ref{tab:exato-mipgap}), de modo que a fronteira de referência é fortemente sub-ótima. O NSGA-IIc supera uma referência degradada, não a fronteira ótima global.

Em quarto lugar, a verificação estrita de factibilidade indica que todas as soluções NSGA-II violam ao menos a restrição \ref{eq:r5} (cobertura mínima). No NSGA-II Estendido, as restrições estruturais \ref{eq:r1}--\ref{eq:r4} são integralmente satisfeitas; no NSGA-II Compacto, 10{,}2\% das soluções apresentam adicionalmente violação de \ref{eq:r4}, associada a uma limitação do decodificador de cromossomo compacto para a clínica U06. Essas limitações do mecanismo de reparo são relevantes para uso prático da meta-heurística e apontam para trabalhos futuros focados em operadores de reparo de cobertura e na correção de \ref{eq:r4} no cromossomo compacto.

Por fim, quando considerado o aspecto operacional, a redução da dimensionalidade do cromossomo e o uso de operadores específicos para a estrutura temporal do problema (cruzamento por blocos diários, realocação de \textit{slot}) mostraram-se decisivos para o desempenho do NSGA-II. A versão Compacta, combinada com os mecanismos da variante Estendida (diversidade adaptativa, evolução em duas fases e inicialização enviesada para extremos), oferece o melhor compromisso entre qualidade e tempo dentre as meta-heurísticas avaliadas.

\chapter{Conclusões e Perspectivas} % 6
\label{cap:fechamento}

Este trabalho investigou o problema de alocação de médicos em uma rede de clínicas de exames com múltiplas unidades autônomas, com base em dados reais de 250 profissionais e 37 unidades coletados ao longo de 2022 e 2023. O problema foi formulado como otimização inteira mista bi-objetiva, com minimização simultânea de exames não atendidos e trocas de clínica em horizontes acíclicos de planejamento, e resolvido por três famílias metodológicas distintas: o método exato por $\varepsilon$-restrito, MIP-heurísticas baseadas em \textit{Relax-and-Fix} e \textit{Fix-and-Optimize}, e meta-heurísticas evolutivas multiobjetivo nas variantes NSGA-II Estendido e Compacto.

As contribuições centrais do estudo podem ser sintetizadas em cinco pontos. Primeiro, a modelagem matemática inteira mista, de caráter bi-objetivo, para escalonamento médico em múltiplas clínicas autônomas, com tratamento conjunto da cobertura operacional e da estabilidade das alocações ao longo do horizonte. Segundo, a adaptação da variável de estado $\gamma_{iud}$ a partir do \textit{General Lot-sizing and Scheduling Problem} \cite{Fleischmann_Meyr_1997}, que viabiliza a detecção correta de trocas mesmo em agendas com lacunas. Terceiro, o acoplamento entre o método $\varepsilon$-restrito e mat-heurísticas para geração estruturada da fronteira de Pareto sob orçamento computacional comum. Quarto, a proposição das variantes NSGA-II Estendido e Compacto, equipadas com operadores customizados de amostragem com consciência de restrições, cruzamento por blocos diários e mutação por realocação de \textit{slot}. Por fim, a avaliação comparativa sistemática das três famílias metodológicas sobre dados reais, organizada em protocolo de dois estágios (parametrização e validação) ao longo de quatro horizontes de planejamento operacionalmente significativos.

\section{Principais Resultados}

A caracterização do método exato, apresentada no Capítulo~\ref{cap:resultados}, evidenciou a transição de regime entre os horizontes de 15 e 30 dias. No ciclo quinzenal, o Gurobi atinge otimalidade comprovada em parte significativa dos pontos $\varepsilon$, com gap MIP médio de $4{,}1\%$, e tempo total inferior ao limite de tempo alocado. A partir de 30 dias, o crescimento combinatório das variáveis $x_{iudt}$ inviabiliza a convergência: o gap médio sobe para $42{,}7\%$ aos 30 dias, $78{,}1\%$ aos 60 dias e $96{,}8\%$ aos 90 dias, com saturação completa do limite de 3.600 segundos por ponto $\varepsilon$. Esse padrão fundamenta a necessidade de heurísticas em horizontes operacionais maiores.

Entre as MIP-heurísticas, o método híbrido (Relax-and-Fix seguido de Fix-and-Optimize) apresentou desempenho próximo ao do método exato nos horizontes em que este converge, com hipervolume relativo de $0{,}996$ aos 15 dias e $1{,}008$ aos 30 dias. O Relax-and-Fix puro, embora ligeiramente inferior em qualidade ($0{,}984$ e $0{,}852$ nos mesmos horizontes), reduziu o tempo total de execução em até $30\%$ em relação ao exato, consolidando-se como alternativa viável em cenários com restrição temporal acentuada.

Nas meta-heurísticas, a representação compacta mostrou-se decisiva. O NSGA-II Compacto, ao codificar diretamente a associação diária médico-clínica, executa um número substancialmente maior de gerações dentro do limite de tempo, atingindo hipervolume relativo de $0{,}715$ contra $0{,}510$ da variante Estendida na fase de parametrização; no horizonte de 60 dias, o NSGA-IIc obteve HV relativo de $1{,}039$ em relação à referência exata --- que, neste horizonte, opera com MIP gap médio de $78{,}1\%$, tornando a comparação direta não conclusiva quanto à qualidade absoluta. A combinação da representação reduzida com os mecanismos de diversidade adaptativa, evolução em duas fases e inicialização enviesada para extremos, herdados da variante Estendida, configurou a abordagem evolutiva mais robusta entre as avaliadas.

Em conjunto, os resultados sustentam uma recomendação operacional clara para a rede de clínicas: o ciclo quinzenal de reprogramação pode ser tratado pelo método exato; horizontes mensal, bimestral e trimestral exigem necessariamente abordagens aproximadas, com o método híbrido como melhor escolha quando há tempo computacional comparável ao do \textit{solver} exato, e o NSGA-II Compacto como melhor escolha quando o limite de tempo é estritamente limitado.

\section{Limitações}

A principal limitação identificada reside no mecanismo de reparo de restrições do NSGA-II. A verificação estrita das soluções revelou que todas violam a restrição de cobertura mínima (\ref{eq:r5}) em 4 a 6\% das combinações clínica-período, inviabilizando o uso operacional direto sem inspeção e reparo posterior. A análise por variante revela distinções relevantes: no NSGA-II Estendido, as restrições estruturais \ref{eq:r1}--\ref{eq:r4} são integralmente satisfeitas e os valores de $Z_1$ e $Z_2$ foram verificados por Gurobi com variáveis fixadas sem discrepância; no NSGA-II Compacto, 10{,}2\% das soluções apresentam adicionalmente violação da restrição de capacidade de sala (\ref{eq:r4}), concentradas na clínica U06 no turno vespertino, e os valores de $Z_1$ reportados subestimam os exames não atendidos em até 23\% nas soluções no extremo de mínima cobertura --- comprometendo a validade dos objetivos para essa fração das soluções. Em consequência, nenhuma das variantes é diretamente comparável aos métodos MIP que garantem factibilidade, e os resultados do NSGA-II devem ser lidos isoladamente, dentro da família de meta-heurísticas.

A modelagem assume disponibilidade médica e demanda por exames como parâmetros determinísticos, obtidos a partir de registros históricos da rede de clínicas. Na operação real, ambas as grandezas estão sujeitas a variabilidade originada por absenteísmo, picos sazonais de demanda e reagendamentos, que o modelo atual absorve apenas indiretamente mediante otimização periódica em horizontes mais curtos. A variável de estado $\gamma_{iud}$ também exige uma condição inicial não nula para todos os médicos, o que, na ausência de histórico imediato, depende de inicialização aleatória ou de estimativa fornecida pelo gestor.

Permanecem fora da formulação algumas restrições ergonômicas presentes na operação cotidiana da rede: limites superiores ao número de finais de semana consecutivos trabalhados, número mínimo de domingos atribuídos por médico e descansos mandatórios entre turnos. Adicionalmente, a possibilidade de exceções negociadas às restrições~\eqref{eq:r2}-\eqref{eq:r3}, que atualmente impedem o profissional de atender múltiplas clínicas no mesmo dia, não está representada. Essas regras refletem acordos contratuais e operacionais que não foram incorporados para preservar a generalidade do modelo, mas constituem extensões naturais para implantação real.

\section{Trabalhos Futuros}

Dado o escopo abrangido neste trabalho, aliada a diversidade das possibilidades de extensão em pesquisas de Otimização e a riqueza nos resultados obtidos, torna-se natural que certas oportunidades não sejam abordadas. Nesse sentido, existe uma gama de direções investigativas em aberto que podem ser exploradas para aprofundar e ampliar os resultados aqui apresentados, tanto em termos de modelagem quanto de solução.

A princípio, pensa-se no desenvolvimento de operadores de reparo específicos para as limitações identificadas nas variantes do NSGA-II. Para a restrição de cobertura mínima (\ref{eq:r5}), violada em ambas as variantes, estratégias plausíveis incluem reparo \textit{post hoc} via subproblema MIP local, busca local guiada por unidades subcobertas e reformulação dos operadores de cruzamento e mutação para preservar a satisfação de \ref{eq:r5} por construção. Para a violação de capacidade de sala (\ref{eq:r4}) no NSGA-II Compacto, concentrada na clínica U06 no turno vespertino, a correção consiste em revisar o decodificador de cromossomo compacto para garantir o respeito aos limites de salas em todas as configurações de clínica e turno. A literatura recente em escalonamento na saúde \cite{Turhan_Bilgen_2020,Wickert_Kummer_Neto_Boniatti_Buriol_2021} indica que reparos baseados em MIPs locais combinam-se naturalmente com meta-heurísticas evolutivas, padrão coerente com a arquitetura adotada neste estudo.

Outra linha de investigação refere-se ao tratamento estocástico da demanda e da disponibilidade. Modelos robustos ou de programação estocástica em dois estágios podem absorver explicitamente a variabilidade observada na operação, complementados por mecanismos de reescalonamento \textit{online} que ajustem a escala diante de faltas de profissionais, picos de demanda e reagendamentos. A aplicação de técnicas de Aprendizado de Máquina apresenta-se como outra direção promissora em duas frentes complementares: a previsão da demanda por unidade-período, com vistas à substituição das estimativas históricas por projeções calibradas; e a geração de soluções iniciais informadas por padrões aprendidos, fornecendo \textit{warm-starts} de qualidade tanto para o método exato quanto para as mat-heurísticas e variantes do NSGA-II.

Por fim, a incorporação de uma metologia hibrida, aplicando conceitos de solução exata aliados a meta-heurísticas evolutivas em uma única \textit{pipeline} geracional pode explorar os limites computacionais de ambas as abordagens, reduzindo assim o \textit{gap} observado ao aplicar apenas uma metodologia individualizada. A combinação de métodos exatos e meta-heurísticas pode superar as limitações de cada abordagem isoladamente, especialmente em problemas de grande escala e alta complexidade como o escalonamento médico em múltiplas clínicas autônomas.


% ---
% Finaliza a parte no bookmark do PDF, para que se inicie o bookmark na raiz
% ---
\bookmarksetup{startatroot}% 
% ---

% ----------------------------------------------------------
% ELEMENTOS PÓS-TEXTUAIS
% ----------------------------------------------------------
\postextual

% ----------------------------------------------------------
% Referências bibliográficas
% ----------------------------------------------------------
\bibliography{bib/references}

% ---------------------------------------------------------------------
% GLOSSÁRIO
% ---------------------------------------------------------------------

% % Arquivo que contém as definições que vão aparecer no glossário
% \input{tex/3-pos-textual/glossario.tex}
% % Comando para incluir todas as definições do arquivo glossario.tex
% \glsaddall
% % Impressão do glossário
% \printglossaries

% ----------------------------------------------------------
% Apêndices
% ----------------------------------------------------------

% ---
% Inicia os apêndices
% ---
\begin{apendicesenv}

    \chapter{Análise da Base de Dados por Instância}
    \label{apendice-cap:analise-base-dados}
    Este apêndice complementa a Seção~\ref{cap:resultados:sec-analise-base-dados} com o detalhamento numérico do conjunto histórico e das 40 instâncias (10 por horizonte) utilizadas nos experimentos. A organização espelha a do capítulo: a Seção~\ref{apendice:base-geral-detalhe} aprofunda a análise descritiva do conjunto completo (02/01/2022 a 31/07/2023), e as Seções~\ref{apendice:base-h15}--\ref{apendice:base-h90} apresentam, para cada horizonte de planejamento, a tabela estendida das 10 instâncias, com colunas adicionais (médias diárias, número de unidades ativas) que foram omitidas do capítulo por concisão.

\noindent\textbf{Estimativa.}\, Como na Seção~\ref{cap:resultados:base-geral}, todas as razões \emph{Demanda / (Salas\,$\times$\,Turnos)} e \emph{Disponibilidade / Demanda} reportadas neste apêndice são \emph{limitantes ingênuos}. Elas \textbf{ignoram} as restrições \ref{eq:r1}--\ref{eq:r5} do modelo (qualificação médico--unidade, unicidade diária, capacidade por turno, cobertura mínima) e \textbf{não correspondem a uma alocação factível}; servem apenas como referência rápida para o regime oferta\,$\times$\,demanda esperado em cada recorte.

\section{Conjunto Histórico Completo}
\label{apendice:base-geral-detalhe}

A Figura~\ref{fig:apendice-scatter-salas-outros} examina a relação entre a soma mensal de salas disponíveis e três indicadores agregados na mesma granularidade: a demanda mensal de exames, a disponibilidade mensal de médicos e a quantidade de médicos distintos com pelo menos um turno ativo no mês. Demanda e disponibilidade exibem associação aproximadamente linear com a oferta de salas, enquanto a contagem de médicos distintos cresce de forma menos regular: o número de profissionais cadastrados aumenta ao longo do período sem que isso se traduza em proporção direta de salas adicionais, indicando que parte do crescimento da força de trabalho é absorvida pelo aumento da utilização das salas existentes.

\newpage

\begin{figure}[h]
	\centering
	\begin{subfigure}{0.333\textwidth}
		\includegraphics[width=\textwidth]{images/scatter-sala-demanda_2.png}
		\caption{Salas $\times$ Demanda}
		\label{fig:apendice-scatter-sala-dem}
	\end{subfigure}
	\begin{subfigure}{0.333\textwidth}
		\includegraphics[width=\textwidth]{images/scatter-sala-disponibilidade.png}
		\caption{Salas $\times$ Disponibilidade}
		\label{fig:apendice-scatter-sala-disp}
	\end{subfigure}
	\begin{subfigure}{0.333\textwidth}
		\includegraphics[width=\textwidth]{images/scatter-sala-coutmed.png}
		\caption{Salas $\times$ Médicos distintos}
		\label{fig:apendice-scatter-sala-med}
	\end{subfigure}
	\caption{Figuras de dispersão entre indicadores agregados por mês.}
	\label{fig:apendice-scatter-salas-outros}
	\fautor
\end{figure}

A distribuição de salas entre as 37 unidades clínicas, ilustrada na Figura~\ref{fig:apendice-bar-sala-unidade}, é fortemente assimétrica: poucas unidades concentram a maior parte da capacidade agregada de salas-turno, enquanto a cauda longa contém clínicas com participação marginal. Essa heterogeneidade é uma das fontes de dificuldade do modelo MIP: clínicas de grande porte impõem restrições de cobertura mínima (C5) sobre um volume substancial de exames, e a mobilidade dos médicos entre essas unidades (qualificação dada pela matriz $I \times U$) é frequentemente limitada.

\newpage

\begin{figure}[h]
	\centering
	\includegraphics[width=0.8\textwidth]{images/bar-sala-unidade.png}
	\caption{Soma de salas $\times$ turnos por unidade clínica no conjunto histórico (2022--2023).}
	\label{fig:apendice-bar-sala-unidade}
	\fautor
\end{figure}

A Figura~\ref{fig:apendice-histogramas} apresenta as distribuições marginais da demanda e da quantidade de salas disponíveis ao nível mais granular do conjunto (uma observação por tripla data\,$\times$\,turno\,$\times$\,unidade). A demanda exibe distribuição fortemente assimétrica à direita, com massa concentrada em valores baixos a moderados e cauda longa para unidades de grande porte; a quantidade de salas disponíveis concentra-se em 1, 2 e 3 salas por turno, refletindo a estrutura modular das unidades.

\begin{figure}[h]
	\centering
	\begin{subfigure}{0.45\textwidth}
		\includegraphics[width=\textwidth]{images/hist-demanda.png}
		\caption{Demanda por (data, turno, unidade)}
		\label{fig:apendice-hist-demanda}
	\end{subfigure}
	\begin{subfigure}{0.45\textwidth}
		\includegraphics[width=\textwidth]{images/hist-salas.png}
		\caption{Salas disponíveis por (data, turno, unidade)}
		\label{fig:apendice-hist-salas}
	\end{subfigure}
	\caption{Distribuições marginais de demanda e capacidade de salas no conjunto histórico (2022--2023).}
	\label{fig:apendice-histogramas}
	\fautor
\end{figure}

\clearpage

\section{Horizonte de 15 dias}
\label{apendice:base-h15}

O Quadro~\ref{quad:apendice-base-h15}, ao final desta seção, apresenta a caracterização estendida das 10 instâncias quinzenais, ordenadas cronologicamente. Em relação ao Quadro~\ref{quad:base-h15} do capítulo, são acrescentadas as colunas de médias diárias (demanda, disponibilidade e salas\,$\times$\,turnos) e o número de unidades clínicas ativas no recorte.

\noindent No horizonte quinzenal, a soma de exames demandados oscila entre $\sim$28k (recorte de janeiro/2022) e $\sim$44k (recortes de 2023), refletindo o crescimento gradual do volume operacional ao longo do período. A razão \emph{Demanda / (Salas\,$\times$\,Turnos)} permanece em uma faixa estreita (20{,}8--23{,}5), sem ruptura clara entre instâncias de parametrização e de validação. A razão \emph{Disponibilidade / Demanda} inicia em $\sim$1{,}08 nos primeiros recortes de 2022 (excesso de mão-de-obra disponível) e cai para $\sim$0{,}98 nos recortes de meados de 2023 (déficit incipiente), evidenciando que o crescimento da demanda nos últimos meses do conjunto histórico supera o crescimento da disponibilidade médica registrada na base.

\begingroup\let\cleardoublepage\clearpage
\begin{sidewaysquadro}[p]
	\centering
	\caption{Caracterização estendida das 10 instâncias do horizonte de 15 dias.}
	\begin{tabular}{|M{2.6cm}|c|c|c|c|c|c|c|c|c|c|}
\hline
\textbf{Inst\^ancia} & \textbf{M\'ed.} & \textbf{Unid.} & \textbf{$\Sigma$ Dem.} & \textbf{Dem./dia} & \textbf{$\Sigma$ Disp.} & \textbf{Disp./dia} & \textbf{$\Sigma$ Salas$\times$T} & \textbf{Salas/dia} & \textbf{Dem./Oferta} & \textbf{Disp./Dem.} \\
\hline\hline
02/01/22 a 16/01/22 & 127 & 35 & 27.855 & 1857 & 29.950 & 1997 & 1.338 & 89{,}2 & 20{,}82 & 1{,}08 \\
\hline
01/03/22 a 15/03/22 & 136 & 36 & 36.397 & 2426 & 37.374 & 2492 & 1.575 & 105{,}0 & 23{,}11 & 1{,}03 \\
\hline
01/05/22 a 15/05/22 & 145 & 36 & 36.358 & 2424 & 36.726 & 2448 & 1.548 & 103{,}2 & 23{,}49 & 1{,}01 \\
\hline
01/07/22 a 15/07/22 & 143 & 36 & 35.136 & 2342 & 37.459 & 2497 & 1.648 & 109{,}9 & 21{,}32 & 1{,}07 \\
\hline
01/09/22 a 15/09/22 & 139 & 37 & 36.467 & 2431 & 38.108 & 2541 & 1.691 & 112{,}7 & 21{,}57 & 1{,}04 \\
\hline
01/11/22 a 15/11/22 & 147 & 37 & 36.316 & 2421 & 37.454 & 2497 & 1.609 & 107{,}3 & 22{,}57 & 1{,}03 \\
\hline
01/01/23 a 15/01/23 & 150 & 37 & 31.350 & 2090 & 32.371 & 2158 & 1.416 & 94{,}4 & 22{,}14 & 1{,}03 \\
\hline
01/03/23 a 15/03/23 & 162 & 37 & 44.210 & 2947 & 43.319 & 2888 & 1.920 & 128{,}0 & 23{,}03 & 0{,}98 \\
\hline
01/05/23 a 15/05/23 & 171 & 37 & 43.904 & 2927 & 43.193 & 2880 & 1.895 & 126{,}3 & 23{,}17 & 0{,}98 \\
\hline
01/07/23 a 15/07/23 & 171 & 37 & 44.332 & 2955 & 44.019 & 2935 & 2.046 & 136{,}4 & 21{,}67 & 0{,}99 \\
\hline
\end{tabular}

	\label{quad:apendice-base-h15}
	\fautor
\end{sidewaysquadro}
\endgroup

\clearpage

\section{Horizonte de 30 dias}
\label{apendice:base-h30}

O Quadro~\ref{quad:apendice-base-h30}, ao final desta seção, apresenta a caracterização estendida das 10 instâncias mensais. Em relação ao Quadro~\ref{quad:base-h30} do capítulo, são acrescentadas as colunas de médias diárias (demanda, disponibilidade e salas\,$\times$\,turnos) e o número de unidades clínicas ativas no recorte.

\noindent No horizonte mensal, a soma de exames demandados varia entre $\sim$57k (janeiro/2022, instância parcial por iniciar em 02/01) e $\sim$91k (março/2023). A média diária de demanda evolui de $\sim$1.918 para $\sim$3.031 exames, com aumento simultâneo no número de médicos distintos ativos (141 para 182). A razão \emph{Demanda / (Salas\,$\times$\,Turnos)} oscila entre 21{,}5 e 23{,}6, e a razão \emph{Disponibilidade / Demanda} segue o mesmo padrão observado no horizonte quinzenal: começa próxima de 1{,}07 em 2022 e converge para valores ligeiramente inferiores a 1 nos recortes de 2023. A janela mensal já é suficiente para suavizar parte da variabilidade semanal observada no horizonte de 15 dias, mas preserva integralmente o efeito sazonal anual.

\begingroup\let\cleardoublepage\clearpage
\begin{sidewaysquadro}[p]
	\centering
	\caption{Caracterização estendida das 10 instâncias do horizonte de 30 dias.}
	\begin{tabular}{|M{2.6cm}|c|c|c|c|c|c|c|c|c|c|}
\hline
\textbf{Inst\^ancia} & \textbf{M\'ed.} & \textbf{Unid.} & \textbf{$\Sigma$ Dem.} & \textbf{Dem./dia} & \textbf{$\Sigma$ Disp.} & \textbf{Disp./dia} & \textbf{$\Sigma$ Salas$\times$T} & \textbf{Salas/dia} & \textbf{Dem./Oferta} & \textbf{Disp./Dem.} \\
\hline\hline
02/01/22 a 31/01/22 & 141 & 35 & 57.553 & 1918 & 61.723 & 2057 & 2.676 & 89{,}2 & 21{,}51 & 1{,}07 \\
\hline
01/03/22 a 30/03/22 & 144 & 36 & 74.772 & 2492 & 76.578 & 2553 & 3.197 & 106{,}6 & 23{,}39 & 1{,}02 \\
\hline
01/05/22 a 30/05/22 & 154 & 36 & 74.067 & 2469 & 75.659 & 2522 & 3.218 & 107{,}3 & 23{,}02 & 1{,}02 \\
\hline
01/07/22 a 30/07/22 & 157 & 36 & 70.942 & 2365 & 74.545 & 2485 & 3.279 & 109{,}3 & 21{,}64 & 1{,}05 \\
\hline
01/09/22 a 30/09/22 & 146 & 37 & 73.809 & 2460 & 76.854 & 2562 & 3.367 & 112{,}2 & 21{,}92 & 1{,}04 \\
\hline
01/11/22 a 30/11/22 & 154 & 37 & 75.006 & 2500 & 76.229 & 2541 & 3.266 & 108{,}9 & 22{,}97 & 1{,}02 \\
\hline
01/01/23 a 30/01/23 & 163 & 37 & 70.932 & 2364 & 72.476 & 2416 & 3.146 & 104{,}9 & 22{,}55 & 1{,}02 \\
\hline
01/03/23 a 30/03/23 & 175 & 37 & 90.943 & 3031 & 88.115 & 2937 & 3.862 & 128{,}7 & 23{,}55 & 0{,}97 \\
\hline
01/05/23 a 30/05/23 & 178 & 37 & 90.052 & 3002 & 87.432 & 2914 & 3.830 & 127{,}7 & 23{,}51 & 0{,}97 \\
\hline
01/07/23 a 30/07/23 & 182 & 37 & 88.366 & 2946 & 87.615 & 2920 & 3.979 & 132{,}6 & 22{,}21 & 0{,}99 \\
\hline
\end{tabular}

	\label{quad:apendice-base-h30}
	\fautor
\end{sidewaysquadro}
\endgroup

\clearpage

\section{Horizonte de 60 dias}
\label{apendice:base-h60}

O Quadro~\ref{quad:apendice-base-h60}, ao final desta seção, apresenta a caracterização estendida das 10 instâncias bimestrais. Em relação ao Quadro~\ref{quad:base-h60} do capítulo, são acrescentadas as colunas de médias diárias (demanda, disponibilidade e salas\,$\times$\,turnos) e o número de unidades clínicas ativas no recorte.

\noindent No horizonte bimestral, a soma de exames demandados ultrapassa a casa das centenas de milhares (de $\sim$130k a $\sim$180k por instância), com média diária estável na faixa $\sim$2.160--2.993 exames. O número de médicos distintos cresce de 146 (recorte iniciado em 02/01/2022) para 193 (recortes de meados de 2023), confirmando a tendência de expansão da força de trabalho registrada na base. A amplitude das duas razões \emph{oferta\,$\times$\,demanda} se reduz em relação aos horizontes mais curtos --- \emph{Demanda / (Salas\,$\times$\,Turnos)} concentra-se entre 21{,}6 e 23{,}7, e \emph{Disponibilidade / Demanda} entre 0{,}96 e 1{,}06 ---, indicando que janelas mais longas amortecem flutuações pontuais de demanda e disponibilidade.

\begingroup\let\cleardoublepage\clearpage
\begin{sidewaysquadro}[p]
	\centering
	\caption{Caracterização estendida das 10 instâncias do horizonte de 60 dias.}
	\begin{tabular}{|M{2.6cm}|c|c|c|c|c|c|c|c|c|c|}
\hline
\textbf{Inst\^ancia} & \textbf{M\'ed.} & \textbf{Unid.} & \textbf{$\Sigma$ Dem.} & \textbf{Dem./dia} & \textbf{$\Sigma$ Disp.} & \textbf{Disp./dia} & \textbf{$\Sigma$ Salas$\times$T} & \textbf{Salas/dia} & \textbf{Dem./Oferta} & \textbf{Disp./Dem.} \\
\hline\hline
02/01/22 a 02/03/22 & 146 & 35 & 129.617 & 2160 & 136.118 & 2269 & 5.784 & 96{,}4 & 22{,}41 & 1{,}05 \\
\hline
01/02/22 a 01/04/22 & 149 & 36 & 147.997 & 2467 & 151.902 & 2532 & 6.333 & 105{,}5 & 23{,}37 & 1{,}03 \\
\hline
01/04/22 a 30/05/22 & 162 & 36 & 147.217 & 2454 & 149.008 & 2483 & 6.421 & 107{,}0 & 22{,}93 & 1{,}01 \\
\hline
01/06/22 a 30/07/22 & 169 & 36 & 141.915 & 2365 & 149.923 & 2499 & 6.571 & 109{,}5 & 21{,}60 & 1{,}06 \\
\hline
01/08/22 a 29/09/22 & 160 & 37 & 152.426 & 2540 & 158.342 & 2639 & 6.843 & 114{,}0 & 22{,}27 & 1{,}04 \\
\hline
01/10/22 a 29/11/22 & 158 & 37 & 148.527 & 2475 & 152.139 & 2536 & 6.542 & 109{,}0 & 22{,}70 & 1{,}02 \\
\hline
01/12/22 a 29/01/23 & 175 & 37 & 133.009 & 2217 & 136.052 & 2268 & 5.931 & 98{,}8 & 22{,}43 & 1{,}02 \\
\hline
01/02/23 a 01/04/23 & 184 & 37 & 174.722 & 2912 & 171.118 & 2852 & 7.552 & 125{,}9 & 23{,}14 & 0{,}98 \\
\hline
01/04/23 a 30/05/23 & 193 & 37 & 177.365 & 2956 & 170.870 & 2848 & 7.489 & 124{,}8 & 23{,}68 & 0{,}96 \\
\hline
01/06/23 a 30/07/23 & 193 & 37 & 179.589 & 2993 & 177.056 & 2951 & 7.905 & 131{,}8 & 22{,}72 & 0{,}99 \\
\hline
\end{tabular}

	\label{quad:apendice-base-h60}
	\fautor
\end{sidewaysquadro}
\endgroup

\clearpage

\section{Horizonte de 90 dias}
\label{apendice:base-h90}

O Quadro~\ref{quad:apendice-base-h90}, ao final desta seção, apresenta a caracterização estendida das 10 instâncias trimestrais. Em relação ao Quadro~\ref{quad:base-h90} do capítulo, são acrescentadas as colunas de médias diárias (demanda, disponibilidade e salas\,$\times$\,turnos) e o número de unidades clínicas ativas no recorte.

\noindent No horizonte trimestral, cada instância corresponde a uma janela operacional completa de planejamento primário, com soma de exames demandados entre $\sim$205k e $\sim$273k. A média diária de demanda evolui de $\sim$2.284 (primeiro trimestre de 2022) para $\sim$3.030 (recortes de meados de 2023), acompanhada pela ampliação do quadro médico ativo (de 154 para 198 profissionais). As razões \emph{oferta\,$\times$\,demanda} apresentam a menor amplitude entre os quatro horizontes (\emph{Demanda / (Salas\,$\times$\,Turnos)} entre 21{,}9 e 23{,}6, \emph{Disponibilidade / Demanda} entre 0{,}97 e 1{,}05), refletindo o efeito de média do horizonte mais longo. Os $\sim$10\textsuperscript{6} variáveis binárias geradas por uma instância deste horizonte (cf.~Seção~\ref{cap:resultados:exato-gap}) impedem a convergência do método exato no limite de tempo adotado, motivando a comparação com as MIP-heurísticas e meta-heurísticas conduzida no Capítulo~\ref{cap:resultados}.

\begingroup\let\cleardoublepage\clearpage
\begin{sidewaysquadro}[p]
	\centering
	\caption{Caracterização estendida das 10 instâncias do horizonte de 90 dias.}
	\begin{tabular}{|M{2.6cm}|c|c|c|c|c|c|c|c|c|c|}
\hline
\textbf{Inst\^ancia} & \textbf{M\'ed.} & \textbf{Unid.} & \textbf{$\Sigma$ Dem.} & \textbf{Dem./dia} & \textbf{$\Sigma$ Disp.} & \textbf{Disp./dia} & \textbf{$\Sigma$ Salas$\times$T} & \textbf{Salas/dia} & \textbf{Dem./Oferta} & \textbf{Disp./Dem.} \\
\hline\hline
02/01/22 a 01/04/22 & 154 & 36 & 205.550 & 2284 & 213.625 & 2374 & 9.009 & 100{,}1 & 22{,}82 & 1{,}04 \\
\hline
01/02/22 a 01/05/22 & 162 & 36 & 218.818 & 2431 & 222.837 & 2476 & 9.439 & 104{,}9 & 23{,}18 & 1{,}02 \\
\hline
01/04/22 a 29/06/22 & 175 & 36 & 218.683 & 2430 & 224.851 & 2498 & 9.728 & 108{,}1 & 22{,}48 & 1{,}03 \\
\hline
01/06/22 a 29/08/22 & 174 & 37 & 217.185 & 2413 & 227.748 & 2531 & 9.910 & 110{,}1 & 21{,}92 & 1{,}05 \\
\hline
01/08/22 a 29/10/22 & 162 & 37 & 228.215 & 2536 & 236.405 & 2627 & 10.217 & 113{,}5 & 22{,}34 & 1{,}04 \\
\hline
01/09/22 a 29/11/22 & 162 & 37 & 222.336 & 2470 & 228.993 & 2544 & 9.909 & 110{,}1 & 22{,}44 & 1{,}03 \\
\hline
01/11/22 a 29/01/23 & 180 & 37 & 208.015 & 2311 & 212.281 & 2359 & 9.197 & 102{,}2 & 22{,}62 & 1{,}02 \\
\hline
01/01/23 a 31/03/23 & 189 & 37 & 246.456 & 2738 & 244.418 & 2716 & 10.724 & 119{,}2 & 22{,}98 & 0{,}99 \\
\hline
01/03/23 a 29/05/23 & 198 & 37 & 267.207 & 2969 & 257.980 & 2866 & 11.323 & 125{,}8 & 23{,}60 & 0{,}97 \\
\hline
01/05/23 a 29/07/23 & 195 & 37 & 272.689 & 3030 & 267.385 & 2971 & 11.851 & 131{,}7 & 23{,}01 & 0{,}98 \\
\hline
\end{tabular}

	\label{quad:apendice-base-h90}
	\fautor
\end{sidewaysquadro}
\endgroup

\clearpage

\section{Síntese}
\label{apendice:base-sintese}

A leitura conjunta dos quatro horizontes evidencia três padrões consistentes ao longo do conjunto histórico. Primeiro, o crescimento gradual do volume operacional --- demanda, disponibilidade e quantidade de médicos --- progride de forma aproximadamente linear entre 02/01/2022 e 31/07/2023, sem rupturas estruturais identificáveis nos quadros. Segundo, a razão \emph{Disponibilidade / Demanda} migra de um regime de leve excesso ($>1$) em 2022 para um regime de leve déficit ($<1$) em 2023, indicando que a expansão da força de trabalho registrada na base não acompanhou integralmente a ampliação da demanda no último semestre coberto. Terceiro, a amplitude das duas razões \emph{oferta\,$\times$\,demanda} diminui monotonicamente com o aumento do horizonte (de [20{,}8;\,23{,}5] em 15 dias para [21{,}9;\,23{,}6] em 90 dias na razão \emph{Demanda / (Salas\,$\times$\,Turnos)}), confirmando o efeito esperado de amortecimento estatístico em janelas mais longas. \textbf{Reitera-se} que estas razões mantêm o caráter de \emph{limitantes ingênuos} discutido na Seção~\ref{cap:resultados:base-geral}: oferecem um diagnóstico de regime, não uma medida de factibilidade.


    \chapter{Análise Detalhada dos Resultados Exatos}
    \label{apendice-cap:exato-detalhe}
    Este apêndice complementa a Seção~\ref{cap:resultados:exato} com o detalhamento, por instância e por ponto $\varepsilon$, do método exato. Para cada um dos quatro horizontes, é apresentada uma tabela com os 7 pontos $\varepsilon$ nas colunas e as 10 instâncias nas linhas; em cada célula, reporta-se o terno $Z_1$ / gap (\%) / tempo (s).

Valores em negrito indicam pontos resolvidos até a otimalidade comprovada; o símbolo $^{\dagger}$ assinala pontos cuja execução atingiu o limite de tempo de 3.600~s; e a notação \emph{n/f} (não factível) é utilizada quando o \textit{solver} encerra o orçamento sem encontrar nenhuma solução factível para o ponto $\varepsilon$. Cada tabela é acompanhada de três figuras: (i) a sobreposição das fronteiras de Pareto das 10 instâncias, (ii) um mapa de calor do MIP gap (instâncias~$\times$~pontos $\varepsilon$) e (iii) um diagrama de caixa do tempo de execução por ponto $\varepsilon$.

\section{Horizonte de 15 dias}
\label{apendice:exato-h15}

No horizonte de 15 dias, observam-se dois regimes distintos de comportamento entre as instâncias. Sete das 10 instâncias atingem otimalidade comprovada nos pontos $\varepsilon$ extremos (próximos ao mínimo de $Z_1$ e ao mínimo de $Z_2$), com tempos de execução tipicamente abaixo de 2.000~s. Os pontos intermediários da grade ($\varepsilon_2$ a $\varepsilon_4$) concentram a maior parte dos casos não convergentes, com gaps tipicamente entre 1\% e 7\%. Duas instâncias \textit{outlier} (\textit{2023mar01} e \textit{2023jul01}) destoam do padrão: nelas, o \textit{solver} sequer encontra solução factível em $\varepsilon_1$ ou apresenta gap acima de 30\% nos pontos centrais, sugerindo perfis de demanda particularmente apertados em relação à oferta de salas. Essas instâncias deslocam para cima a média e o desvio padrão reportados na Tabela~\ref{tab:exato-mipgap} e respondem pela maior parte da dispersão observada.

A figura \ref{fig:exato-pareto-15dias} apresenta a sobreposição das fronteiras de Pareto das 10 instâncias para o horizonte de 15 dias. Para essa janela de tempo, observa-se o padrão esperado de fronteiras decrescentes, com $Z_2$ diminuindo à medida que $Z_1$ aumenta. A forma das fronteiras é relativamente suave, sem grandes saltos ou irregularidades, o que sugere que o método exato consegue explorar de forma eficaz a região próxima à fronteira ótima. A dispersão entre as fronteiras das diferentes instâncias é moderada, indicando que os perfis sazonais de demanda em 2022 e 2023 não geram variações extremas na relação entre $Z_1$ e $Z_2$ para esse horizonte.

\begin{figure}[h]
	\centering
	\includegraphics[width=0.69\textwidth]{images/exato_pareto_overlay_15days.png}
	\caption{Fronteiras de Pareto sobrepostas no horizonte de 15 dias.}
	\label{fig:exato-pareto-15dias}
	\fautor
\end{figure}

O mapa de calor, presente em Figura~\ref{fig:exato-heatmap-15dias}, sintetiza a distribuição dos gaps MIP por instância e ponto $\varepsilon$. Os pontos $\varepsilon_1$ e $\varepsilon_7$ apresentam uma concentração de gaps próximos a zero, refletindo os casos resolvidos até a otimalidade comprovada. Já os pontos intermediários exibem uma maior variabilidade, com alguns casos apresentando gaps acima de 5\%, o que evidencia a dificuldade de convergência para essas configurações.

\begin{figure}[h]
	\centering
	\includegraphics[width=0.64\textwidth]{images/exato_mipgap_heatmap_15days.png}
	\caption{Mapa de calor do MIP gap no horizonte de 15 dias.}
	\label{fig:exato-heatmap-15dias}
	\fautor
\end{figure}

O diagrama de caixa do tempo de execução por ponto $\varepsilon$, apresentado na Figura~\ref{fig:exato-runtime-15dias}, mostra que os tempos de execução variam significativamente entre os pontos da grade. Os pontos extremos ($\varepsilon_1$ e $\varepsilon_7$) apresentam tempos menores, enquanto os pontos intermediários valores mais elevados, com alguns casos atingindo o limite de tempo. A presença de outliers nos pontos intermediários sugere que, para certas instâncias, o método exato enfrenta dificuldades significativas para encontrar soluções factíveis ou para comprovar a otimalidade dentro do limite de tempo estabelecido. Apesar disso, observa-se alta variabilidade nos pontos mais extremos, indicativo de que o padrão de sazonalidade de demanda pode influenciar fortemente a dificuldade de resolução mesmo para configurações mais restritivas ou mais relaxadas.

\begin{figure}[h]
	\centering
	\includegraphics[width=0.65\textwidth]{images/exato_runtime_boxplot_15days.png}
	\caption{Tempo de execução por ponto $\varepsilon$ no horizonte de 15 dias.}
	\label{fig:exato-runtime-15dias}
	\fautor
\end{figure}

A Tabela~\ref{tab:exato-triplet-15dias} apresenta os resultados completos por instância e ponto $\varepsilon$ no horizonte de 15 dias.

\clearpage

% Tabela
\begin{sidewaystable}[p]
	\centering
	\caption{Resultados por ponto $\varepsilon$ e instância, horizonte de 15 dias.}
	\label{tab:exato-triplet-15dias}
	\resizebox{\textheight}{!}{{\footnotesize
\begin{tabular}{@{}M{2.5cm}ccccccccccccccccccccc@{}}
    \toprule
    \multirow{2}{*}{\textbf{Inst\^ancia}} & \multicolumn{3}{c}{$\varepsilon_{1}$} & \multicolumn{3}{c}{$\varepsilon_{2}$} & \multicolumn{3}{c}{$\varepsilon_{3}$} & \multicolumn{3}{c}{$\varepsilon_{4}$} & \multicolumn{3}{c}{$\varepsilon_{5}$} & \multicolumn{3}{c}{$\varepsilon_{6}$} & \multicolumn{3}{c}{$\varepsilon_{7}$}                                                                                                                                                                                                                                                                                                                                                                                      \\

    \cmidrule(lr){2-4}\cmidrule(lr){5-7}\cmidrule(lr){8-10}\cmidrule(lr){11-13}\cmidrule(lr){14-16}\cmidrule(lr){17-19}\cmidrule(lr){20-22}
                                          & $Z_1$                                 & gap\%                                   & t(s)                                    & $Z_1$                                 & gap\%                                   & t(s)                                    & $Z_1$                                 & gap\%          & t(s)                              & $Z_1$                   & gap\%         & t(s)                              & $Z_1$                   & gap\%         & t(s)                              & $Z_1$                   & gap\%         & t(s)                              & $Z_1$                   & gap\%         & t(s)                              \\
    \toprule
    2022jan02\newline a\newline 2022jan16 & \textbf{144}                            & 0.0                                     & 566                                     & 390                                     & 1.0                                     & 3600\textsuperscript{\textdagger}       & \textbf{631}                            & 0.0            & 913                               & \textbf{855}              & 0.0           & 2152                              & \textbf{1107}             & 0.0           & 721                               & \textbf{1328}             & 0.0           & 953                               & \textbf{1535}             & 0.0           & 515                               \\
    \midrule

    2022mar01\newline a\newline 2022mar15 & \textbf{191}                            & 0.0                                     & 1610                                    & 508                                     & 3.6                                     & 3600\textsuperscript{\textdagger}       & 813                                     & 2.9            & 3600\textsuperscript{\textdagger} & 1123                      & 2.1           & 3600\textsuperscript{\textdagger} & \textbf{1441}             & 0.0           & 1628                              & \textbf{1784}             & 0.0           & 1081                              & \textbf{1989}             & 0.0           & 3357                              \\
    \midrule

    2022may01\newline a\newline 2022may15 & \textbf{254}                            & 0.0                                     & 367                                     & 600                                     & 1.4                                     & 3600\textsuperscript{\textdagger}       & 947                                     & 1.1            & 3600\textsuperscript{\textdagger} & \textbf{1285}             & 0.0           & 1812                              & \textbf{1611}             & 0.0           & 2616                              & \textbf{1967}             & 0.0           & 753                               & \textbf{2322}             & 0.0           & 706                               \\
    \midrule

    2022jul01\newline a\newline 2022jul15 & \textbf{105}                            & 0.0                                     & 1256                                    & 394                                     & 7.1                                     & 3600\textsuperscript{\textdagger}       & 700                                     & 6.3            & 3600\textsuperscript{\textdagger} & 999                       & 2.5           & 3600\textsuperscript{\textdagger} & 1290                      & 0.9           & 3600\textsuperscript{\textdagger} & \textbf{1612}             & 0.0           & 1135                              & \textbf{1932}             & 0.0           & 1910                              \\
    \midrule

    2022sep01\newline a\newline 2022sep15 & 128                                     & 1.4                                     & 3600\textsuperscript{\textdagger}       & 411                                     & 9.1                                     & 3600\textsuperscript{\textdagger}       & 713                                     & 2.6            & 3600\textsuperscript{\textdagger} & 986                       & 3.5           & 3600\textsuperscript{\textdagger} & \textbf{1299}             & 0.0           & 1853                              & \textbf{1615}             & 0.0           & 1184                              & \textbf{1915}             & 0.0           & 652                               \\
    \midrule

    2022nov01\newline a\newline 2022nov15 & \textbf{150}                            & 0.0                                     & 2330                                    & 462                                     & 5.5                                     & 3600\textsuperscript{\textdagger}       & 781                                     & 1.8            & 3600\textsuperscript{\textdagger} & 1084                      & 2.6           & 3600\textsuperscript{\textdagger} & 1393                      & 2.4           & 3600\textsuperscript{\textdagger} & \textbf{1704}             & 0.0           & 3331                              & \textbf{2013}             & 0.0           & 896                               \\
    \midrule

    2023jan01\newline a\newline 2023jan15 & \textbf{158}                            & 0.0                                     & 487                                     & 447                                     & 2.7                                     & 3600\textsuperscript{\textdagger}       & 757                                     & 1.4            & 3600\textsuperscript{\textdagger} & 1034                      & 1.3           & 3600\textsuperscript{\textdagger} & \textbf{1334}             & 0.0           & 2027                              & 1583                      & 1.9           & 3600\textsuperscript{\textdagger} & \textbf{1955}             & 0.0           & 1963                              \\
    \midrule

    2023mar01\newline a\newline 2023mar15 & 400                                     & 3.2                                     & 3600\textsuperscript{\textdagger}       & 819                                     & 31.4                                    & 3600\textsuperscript{\textdagger}       & 1257                                    & 30.5           & 3600\textsuperscript{\textdagger} & 1677                      & 5.9           & 3600\textsuperscript{\textdagger} & 2078                      & 8.8           & 3600\textsuperscript{\textdagger} & 2540                      & 3.8           & 3600\textsuperscript{\textdagger} & 2865                      & 3.9           & 3600\textsuperscript{\textdagger} \\
    \midrule

    2023may01\newline a\newline 2023may15 & 446                                     & 2.1                                     & 3600\textsuperscript{\textdagger}       & 865                                     & 6.3                                     & 3600\textsuperscript{\textdagger}       & 1300                                    & 6.3            & 3600\textsuperscript{\textdagger} & 1717                      & 3.7           & 3600\textsuperscript{\textdagger} & 2137                      & 3.9           & 3600\textsuperscript{\textdagger} & \textbf{2592}             & 0.0           & 1970                              & 3032                      & 0.0           & 895                               \\
    \midrule

    2023jul01\newline a\newline 2023jul15 & ---                                     & ---                                     & 3600\textsuperscript{\textdagger}       & 1018                                    & 33.9                                    & 3600\textsuperscript{\textdagger}       & 1426                                    & 34.9           & 3600\textsuperscript{\textdagger} & 1827                      & 30.9          & 3600\textsuperscript{\textdagger} & 2251                      & 8.1           & 3600\textsuperscript{\textdagger} & \textbf{2670}             & 0.0           & 3552                              & \textbf{3070}             & 0.0           & 3519                              \\
    \midrule

    \midrule
    \textbf{M\'edia (10 inst.)}           & $\overline{Z_1}{=}220$                & 0.7$\pm$1.2                           & 2102                                    & $\overline{Z_1}{=}591$                & 10.2$\pm$12.1                         & 3600                                    & $\overline{Z_1}{=}932$                & 8.8$\pm$12.8 & 3331                              & $\overline{Z_1}{=}1259$ & 5.3$\pm$9.2 & 3276                              & $\overline{Z_1}{=}1594$ & 2.4$\pm$3.4 & 2685                              & $\overline{Z_1}{=}1940$ & 0.6$\pm$1.3 & 2116                              & $\overline{Z_1}{=}2263$ & 0.4$\pm$1.2 & 1801                              \\

    \bottomrule
\end{tabular}
}
}
	\fautor
\end{sidewaystable}

\clearpage

\section{Horizonte de 30 dias}
\label{apendice:exato-h30}

A passagem do horizonte quinzenal para o mensal marca uma ruptura qualitativa no comportamento do método exato. Nenhum dos 70 pontos $\varepsilon$ é resolvido até a otimalidade comprovada, e todos atingem o limite de tempo de 3.600~s. A particularidade deste horizonte está na frequência com que o \textit{solver} encerra o tempo sem encontrar nenhuma solução factível: 8 das 10 instâncias terminam em situação \emph{n/f} no ponto $\varepsilon_1$ (o mais restritivo em $Z_2$), refletindo a dificuldade do procedimento de \textit{branch-and-bound} em construir uma incumbente respeitando simultaneamente a restrição lexicográfica de $Z_2$ mínimo e a cobertura mínima C5. Para os pontos $\varepsilon_2$ a $\varepsilon_7$, o gap médio cresce de forma aproximadamente linear com a posição na grade, evidenciando que a relaxação $\varepsilon$ não é, por si só, suficiente para tornar o problema tratável neste horizonte.

Novamente a sobreposição das fronteiras de Pareto, apresentado na Figura~\ref{fig:exato-pareto-30dias}, revela o padrão decrescente esperado, mas a interpretação das fronteiras como aproximações da fronteira ótima global fica comprometida. A ausência de soluções factíveis em muitos pontos $\varepsilon$ e a elevada heterogeneidade entre as instâncias indicam que as fronteiras sobrepostas refletem mais a capacidade do método exato de encontrar soluções factíveis dentro do limite de tempo do que uma aproximação consistente da fronteira verdadeira. A dispersão entre as fronteiras é acentuada, com algumas instâncias apresentando formas mais irregulares e amplitudes maiores, o que sugere que os perfis sazonais de demanda em 2022 e 2023 geram variações significativas na dificuldade de resolução para esse horizonte.

\begin{figure}[h]
	\centering
	\includegraphics[width=0.69\textwidth]{images/exato_pareto_overlay_30days.png}
	\caption{Fronteiras de Pareto sobrepostas no horizonte de 30 dias.}
	\label{fig:exato-pareto-30dias}
	\vspace{-2mm}
	\fautor
\end{figure}

O mapa de calor do MIP gap, presente na Figura~\ref{fig:exato-heatmap-30dias}, sintetiza a heterogeneidade de dificuldade entre as instâncias e os pontos $\varepsilon$. Os pontos $\varepsilon_1$ e $\varepsilon_2$ apresenta uma grande concentração de instâncias cujos passos $\epsilon$ estão acima de 40\% do \textit(gap), com adendo para o ponto $\varepsilon_1$, onde 8 das 10 instâncias terminam sem solução factível. Nos pontos intermediários, observa-se uma tendência de aumento do gap médio, com alguns casos apresentando gaps acima de 80\%, evidenciando a dificuldade de convergência para essas configurações. As instâncias \textit{2023mar01} e \textit{2023jul01} continuam a destacar-se no mapa de calor por apresentarem gaps elevados em vários pontos, reforçando a percepção de que esses perfis de demanda são particularmente desafiadores para o método exato neste horizonte.

\begin{figure}[h]
	\centering
	\includegraphics[width=0.68\textwidth]{images/exato_mipgap_heatmap_30days.png}
	\caption{Mapa de calor do MIP gap no horizonte de 30 dias.}
	\label{fig:exato-heatmap-30dias}
	\vspace{-2mm}
	\fautor
\end{figure}

O diagrama de caixa do tempo de execução por ponto $\varepsilon$, apresentado na Figura~\ref{fig:exato-runtime-30dias}, mostra que os tempos de execução atingem o limite de 3.600~s para todas as instâncias em todos os pontos, o que reforça ainda mais a complexidade do problema, mesmo em um horizonte de planejamento baixo de 30 dias. A presença dos \textit{outliers} no gráfico dizem respeito ao \textit{overhead} computacional, que impedia o processo do \textit{solver} the terminar a tarefa exatamente em 3.600~s.

\begin{figure}[h]
	\centering
	\includegraphics[width=0.65\textwidth]{images/exato_runtime_boxplot_30days.png}
	\caption{Tempo de execução por ponto $\varepsilon$ no horizonte de 30 dias.}
	\label{fig:exato-runtime-30dias}
	\vspace{-2mm}
	\fautor
\end{figure}

A Tabela~\ref{tab:exato-triplet-30dias}, ao final desta seção, apresenta os resultados completos por instância e ponto $\varepsilon$ no horizonte de 30 dias.

% Tabela
\clearpage
\begin{sidewaystable}[p]
	\centering
	\caption{Resultados por ponto $\varepsilon$ e instância, horizonte de 30 dias.}
	\label{tab:exato-triplet-30dias}
	\resizebox{\textheight}{!}{{\footnotesize
\begin{tabular}{@{}M{2.5cm}ccccccccccccccccccccc@{}}
    \toprule
    \multirow{2}{*}{\textbf{Inst\^ancia}} & \multicolumn{3}{c}{$\varepsilon_{1}$} & \multicolumn{3}{c}{$\varepsilon_{2}$} & \multicolumn{3}{c}{$\varepsilon_{3}$} & \multicolumn{3}{c}{$\varepsilon_{4}$} & \multicolumn{3}{c}{$\varepsilon_{5}$} & \multicolumn{3}{c}{$\varepsilon_{6}$} & \multicolumn{3}{c}{$\varepsilon_{7}$}                                                                                                                                                                                                                                                                                                                                                                          \\
    \cmidrule(lr){2-4}\cmidrule(lr){5-7}\cmidrule(lr){8-10}\cmidrule(lr){11-13}\cmidrule(lr){14-16}\cmidrule(lr){17-19}\cmidrule(lr){20-22}
                                          & $Z_1$                                 & gap\%                                 & t(s)                                  & $Z_1$                                 & gap\%                                 & t(s)                                  & $Z_1$                                 & gap\%         & t(s)                              & $Z_1$                   & gap\%         & t(s)                              & $Z_1$                   & gap\%        & t(s)                              & $Z_1$                   & gap\%        & t(s)                              & $Z_1$                   & gap\%        & t(s)                              \\
    \toprule
    2022jan02\newline a\newline 2022jan31 & 294                                   & 5.8                                   & 3600\textsuperscript{\textdagger}     & 796                                   & 36.6                                  & 3600\textsuperscript{\textdagger}     & 1312                                  & 32.8          & 3600\textsuperscript{\textdagger} & 1822                    & 35.2          & 3600\textsuperscript{\textdagger} & 2305                    & 34.7         & 3600\textsuperscript{\textdagger} & 2803                    & 34.9         & 3600\textsuperscript{\textdagger} & 3156                    & 24.1         & 3600\textsuperscript{\textdagger} \\
    \midrule

    2022mar01\newline a\newline 2022mar30 & ---                                   & ---                                   & 3600\textsuperscript{\textdagger}     & 1136                                  & 41.6                                  & 3600\textsuperscript{\textdagger}     & 1795                                  & 45.9          & 3600\textsuperscript{\textdagger} & 2434                    & 38.1          & 3601\textsuperscript{\textdagger} & 3123                    & 28.5         & 3600\textsuperscript{\textdagger} & 3668                    & 30.4         & 3600\textsuperscript{\textdagger} & 4194                    & 28.7         & 3600\textsuperscript{\textdagger} \\
    \midrule

    2022may01\newline a\newline 2022may30 & ---                                   & ---                                   & 3600\textsuperscript{\textdagger}     & 1238                                  & 49.3                                  & 3600\textsuperscript{\textdagger}     & 2029                                  & 48.7          & 3600\textsuperscript{\textdagger} & 2808                    & 34.7          & 3600\textsuperscript{\textdagger} & 3569                    & 37.1         & 3600\textsuperscript{\textdagger} & 4338                    & 31.2         & 3600\textsuperscript{\textdagger} & 4949                    & 33.1         & 3600\textsuperscript{\textdagger} \\
    \midrule

    2022jul01\newline a\newline 2022jul30 & ---                                   & ---                                   & 3600\textsuperscript{\textdagger}     & 898                                   & 42.6                                  & 3600\textsuperscript{\textdagger}     & 1505                                  & 50.3          & 3600\textsuperscript{\textdagger} & 2114                    & 45.4          & 3600\textsuperscript{\textdagger} & 2704                    & 40.9         & 3600\textsuperscript{\textdagger} & 3296                    & 35.1         & 3600\textsuperscript{\textdagger} & 3845                    & 29.1         & 3600\textsuperscript{\textdagger} \\
    \midrule

    2022sep01\newline a\newline 2022sep30 & ---                                   & ---                                   & 3600\textsuperscript{\textdagger}     & 915                                   & 50.9                                  & 3600\textsuperscript{\textdagger}     & 1591                                  & 56.6          & 3600\textsuperscript{\textdagger} & 2275                    & 41.9          & 3600\textsuperscript{\textdagger} & 2965                    & 34.6         & 3600\textsuperscript{\textdagger} & 3625                    & 28.3         & 3600\textsuperscript{\textdagger} & 4320                    & 25.7         & 3600\textsuperscript{\textdagger} \\
    \midrule

    2022nov01\newline a\newline 2022nov30 & ---                                   & ---                                   & 3600\textsuperscript{\textdagger}     & 1099                                  & 43.4                                  & 3600\textsuperscript{\textdagger}     & 1791                                  & 47.0          & 3600\textsuperscript{\textdagger} & 2500                    & 38.8          & 3600\textsuperscript{\textdagger} & 3163                    & 35.6         & 3600\textsuperscript{\textdagger} & 3729                    & 31.9         & 3600\textsuperscript{\textdagger} & 4323                    & 31.4         & 3600\textsuperscript{\textdagger} \\
    \midrule

    2023jan01\newline a\newline 2023jan30 & ---                                   & ---                                   & 3600\textsuperscript{\textdagger}     & 1157                                  & 47.0                                  & 3600\textsuperscript{\textdagger}     & 1793                                  & 51.9          & 3600\textsuperscript{\textdagger} & 2463                    & 54.0          & 3600\textsuperscript{\textdagger} & 3084                    & 43.4         & 3600\textsuperscript{\textdagger} & 3767                    & 33.3         & 3600\textsuperscript{\textdagger} & 4374                    & 35.8         & 3600\textsuperscript{\textdagger} \\
    \midrule

    2023mar01\newline a\newline 2023mar30 & ---                                   & ---                                   & 3600\textsuperscript{\textdagger}     & 2091                                  & 61.6                                  & 3600\textsuperscript{\textdagger}     & 3143                                  & 68.5          & 3600\textsuperscript{\textdagger} & 4183                    & 58.3          & 3600\textsuperscript{\textdagger} & 5228                    & 45.8         & 3600\textsuperscript{\textdagger} & 6202                    & 34.8         & 3600\textsuperscript{\textdagger} & 7198                    & 28.9         & 3600\textsuperscript{\textdagger} \\
    \midrule

    2023may01\newline a\newline 2023may30 & 2089                                  & 71.7                                  & 3600\textsuperscript{\textdagger}     & 2958                                  & 71.0                                  & 3600\textsuperscript{\textdagger}     & 3830                                  & 61.6          & 3600\textsuperscript{\textdagger} & 4698                    & 52.8          & 3600\textsuperscript{\textdagger} & 5526                    & 48.1         & 3600\textsuperscript{\textdagger} & 6447                    & 35.2         & 3600\textsuperscript{\textdagger} & 7149                    & 29.6         & 3600\textsuperscript{\textdagger} \\
    \midrule

    2023jul01\newline a\newline 2023jul30 & ---                                   & ---                                   & 3600\textsuperscript{\textdagger}     & 2060                                  & 75.0                                  & 3600\textsuperscript{\textdagger}     & 2972                                  & 71.4          & 3600\textsuperscript{\textdagger} & 3881                    & 70.6          & 3600\textsuperscript{\textdagger} & 4764                    & 50.2         & 3600\textsuperscript{\textdagger} & 5657                    & 42.3         & 3600\textsuperscript{\textdagger} & 6345                    & 41.2         & 3600\textsuperscript{\textdagger} \\
    \midrule

    \midrule
    \textbf{M\'edia (10 inst.)}           & $\overline{Z_1}{=}1192$               & 38.7$\pm$46.6                         & 3600                                  & $\overline{Z_1}{=}1435$               & 51.9$\pm$13.0                         & 3600                                  & $\overline{Z_1}{=}2176$               & 53.5$\pm$11.5 & 3600                              & $\overline{Z_1}{=}2918$ & 47.0$\pm$11.7 & 3600                              & $\overline{Z_1}{=}3643$ & 39.9$\pm$6.9 & 3600                              & $\overline{Z_1}{=}4353$ & 33.7$\pm$3.8 & 3600                              & $\overline{Z_1}{=}4985$ & 30.7$\pm$5.0 & 3600                              \\
    \bottomrule
\end{tabular}
}
}
	\fautor
\end{sidewaystable}

\clearpage

\section{Horizonte de 60 dias}
\label{apendice:exato-h60}

No horizonte bimestral, o gap MIP médio dobra em relação ao horizonte mensal, alcançando $78{,}11\% \pm 14{,}73$ pontos percentuais. A elevada dispersão sinaliza que, em algumas instâncias, o limitante inferior reportado pelo Gurobi permanece próximo ao valor da relaxação linear, distante do incumbente factível encontrado. As fronteiras sobrepostas (Figura~\ref{fig:exato-pareto-60dias}) mantêm a forma esperada, decrescente de $Z_2$ em função de $Z_1$, mas a sua interpretação como aproximações da fronteira verdadeira fica comprometida: pontos que aparentam ser não-dominados podem, na realidade, ser dominados por soluções factíveis ainda não exploradas dentro do orçamento de 3.600~s. Esta limitação justifica, na Seção~\ref{cap:resultados:comparacao}, a leitura cautelosa das métricas comparativas a partir deste horizonte.

A partir deste horizonte de planejamento, já podemos observar instabilidades significativas no comportamento do método exato, evidenciados pelo comportamento anormal das fronteiras de Pareto na Figura \ref{fig:exato-pareto-60dias}. Embora instâncias com padrão de sazonalidade de demanda mais suave apresentem fronteiras relativamente regulares, algumas instâncias apresentaram comportamento errático, ou até mesmo degradação na fronteira em pontos intermediários oriundos da incapacidade do método de encontrar soluções com baixo \textit{gap} dentro do limite de tempo. A heterogeneidade entre as instâncias é acentuada, com algumas apresentando formas mais irregulares e amplitudes maiores, o que sugere que os perfis sazonais de demanda em 2022 e 2023 geram variações significativas na dificuldade de resolução para esse horizonte.

\begin{figure}[h]
	\centering
	\includegraphics[width=0.69\textwidth]{images/exato_pareto_overlay_60days.png}
	\caption{Fronteiras de Pareto sobrepostas no horizonte de 60 dias.}
	\label{fig:exato-pareto-60dias}
	\fautor
\end{figure}

Apesar disso, o mapa de calor do MIP gap, presente na Figura~\ref{fig:exato-heatmap-60dias}, revela que não houve instâncias com pontos infactíveis, sem solução encontrada dentro do limite de tempo, representando um avanço em relação ao horizonte mensal, provavelmente associado a uma maior flexibilidade na construção de incumbentes factíveis. O gap médio, embora elevado, apresenta uma amplitude menor do que no horizonte mensal, o que sugere que o método exato consegue encontrar soluções factíveis com qualidade relativamente consistente entre as instâncias, mesmo que ainda distantes do limitante inferior reportado pelo \textit{solver}.
\begin{figure}[h]
	\centering
	\includegraphics[width=0.68\textwidth]{images/exato_mipgap_heatmap_60days.png}
	\caption{Mapa de calor do MIP gap no horizonte de 60 dias.}
	\label{fig:exato-heatmap-60dias}
	\fautor
\end{figure}

Assim como observado nos tempos de execução do horizonte mensal, o diagrama de caixa do tempo de execução por ponto $\varepsilon$, apresentado na Figura~\ref{fig:exato-runtime-60dias}, mostra que os tempos de execução atingem o limite de 3.600~s para todas as instâncias em todos os pontos, o que reforça ainda mais a complexidade do problema, mesmo em um horizonte de planejamento de 60 dias. A presença dos \textit{outliers} no gráfico dizem respeito ao \textit{overhead} computacional, que impedia o processo do \textit{solver} de terminar a tarefa exatamente em 3.600~s.

\begin{figure}[h]
	\centering
	\includegraphics[width=0.65\textwidth]{images/exato_runtime_boxplot_60days.png}
	\caption{Tempo de execução por ponto $\varepsilon$ no horizonte de 60 dias.}
	\label{fig:exato-runtime-60dias}
	\fautor
\end{figure}


A Tabela~\ref{tab:exato-triplet-60dias}, ao final desta seção, apresenta os resultados completos por instância e ponto $\varepsilon$ no horizonte de 60 dias.

% Tabela
\clearpage

\begin{sidewaystable}[p]
	\centering
	\caption{Resultados por ponto $\varepsilon$ e instância, horizonte de 60 dias.}
	\label{tab:exato-triplet-60dias}
	\resizebox{\textheight}{!}{{\footnotesize
\begin{tabular}{@{}M{2.5cm}ccccccccccccccccccccc@{}}
\toprule
\multirow{2}{*}{\textbf{Inst\^ancia}} & \multicolumn{3}{c}{$\varepsilon_{1}$} & \multicolumn{3}{c}{$\varepsilon_{2}$} & \multicolumn{3}{c}{$\varepsilon_{3}$} & \multicolumn{3}{c}{$\varepsilon_{4}$} & \multicolumn{3}{c}{$\varepsilon_{5}$} & \multicolumn{3}{c}{$\varepsilon_{6}$} & \multicolumn{3}{c}{$\varepsilon_{7}$} \\
\cmidrule(lr){2-4}\cmidrule(lr){5-7}\cmidrule(lr){8-10}\cmidrule(lr){11-13}\cmidrule(lr){14-16}\cmidrule(lr){17-19}\cmidrule(lr){20-22}
 & $Z_1$ & gap\% & t(s) & $Z_1$ & gap\% & t(s) & $Z_1$ & gap\% & t(s) & $Z_1$ & gap\% & t(s) & $Z_1$ & gap\% & t(s) & $Z_1$ & gap\% & t(s) & $Z_1$ & gap\% & t(s) \\
\toprule
2022jan02\newline a\newline 2022mar02 & 1394 & 77.8 & 3600\textsuperscript{\textdagger}& 2722 & 79.9 & 3600\textsuperscript{\textdagger}& 4055 & 70.0 & 3600\textsuperscript{\textdagger}& 5358 & 69.0 & 3600\textsuperscript{\textdagger}& 6622 & 55.7 & 3600\textsuperscript{\textdagger}& 7990 & 54.8 & 3601\textsuperscript{\textdagger}& 8104 & 53.7 & 3600\textsuperscript{\textdagger} \\
\midrule

2022feb01\newline a\newline 2022apr01 & 1868 & 92.1 & 3601\textsuperscript{\textdagger}& 3324 & 91.7 & 3600\textsuperscript{\textdagger}& 4770 & 80.0 & 3600\textsuperscript{\textdagger}& 6227 & 69.8 & 3600\textsuperscript{\textdagger}& 7635 & 53.3 & 3600\textsuperscript{\textdagger}& 8933 & 44.2 & 3600\textsuperscript{\textdagger}& 10108 & 42.7 & 3600\textsuperscript{\textdagger} \\
\midrule

2022apr01\newline a\newline 2022may30 & 2570 & 91.2 & 3600\textsuperscript{\textdagger}& 3826 & 88.9 & 3600\textsuperscript{\textdagger}& 5064 & 82.5 & 3602\textsuperscript{\textdagger}& 6303 & 73.2 & 3600\textsuperscript{\textdagger}& 7505 & 69.9 & 3600\textsuperscript{\textdagger}& 8756 & 56.3 & 3600\textsuperscript{\textdagger}& 9933 & 52.8 & 3600\textsuperscript{\textdagger} \\
\midrule
2022jun01\newline a\newline 2022jul30 & 2070 & 93.0 & 3600\textsuperscript{\textdagger}& 3196 & 92.9 & 3600\textsuperscript{\textdagger}& 4344 & 86.3 & 3600\textsuperscript{\textdagger}& 5504 & 76.2 & 3600\textsuperscript{\textdagger}& 6649 & 73.0 & 3600\textsuperscript{\textdagger}& 7720 & 68.6 & 3600\textsuperscript{\textdagger}& 8668 & 60.7 & 3600\textsuperscript{\textdagger} \\
\midrule

2022aug01\newline a\newline 2022sep29 & 2305 & 94.7 & 3600\textsuperscript{\textdagger}& 3533 & 93.9 & 3603\textsuperscript{\textdagger}& 4720 & 95.5 & 3600\textsuperscript{\textdagger}& 5972 & 74.5 & 3601\textsuperscript{\textdagger}& 7192 & 72.7 & 3600\textsuperscript{\textdagger}& 8380 & 54.6 & 3600\textsuperscript{\textdagger}& 9512 & 50.4 & 3600\textsuperscript{\textdagger} \\
\midrule

2022oct01\newline a\newline 2022nov29 & 2263 & 93.6 & 3600\textsuperscript{\textdagger}& 3541 & 92.8 & 3600\textsuperscript{\textdagger}& 4785 & 88.1 & 3600\textsuperscript{\textdagger}& 6066 & 82.4 & 3600\textsuperscript{\textdagger}& 7331 & 72.5 & 3600\textsuperscript{\textdagger}& 8559 & 70.2 & 3600\textsuperscript{\textdagger}& 9419 & 62.9 & 3600\textsuperscript{\textdagger} \\
\midrule

2022dec01\newline a\newline 2023jan29 & 2141 & 91.0 & 3600\textsuperscript{\textdagger}& 3160 & 89.0 & 3600\textsuperscript{\textdagger}& 4173 & 86.5 & 3601\textsuperscript{\textdagger}& 5205 & 80.8 & 3600\textsuperscript{\textdagger}& 6221 & 68.9 & 3600\textsuperscript{\textdagger}& 7243 & 70.7 & 3600\textsuperscript{\textdagger}& 8053 & 69.5 & 3600\textsuperscript{\textdagger} \\
\midrule

2023feb01\newline a\newline 2023apr01 & 5894 & 96.2 & 3601\textsuperscript{\textdagger}& 7449 & 96.6 & 3600\textsuperscript{\textdagger}& 9044 & 97.0 & 3600\textsuperscript{\textdagger}& 10796 & 93.7 & 3655\textsuperscript{\textdagger}& 12403 & 81.0 & 3600\textsuperscript{\textdagger}& 14028 & 73.6 & 3600\textsuperscript{\textdagger}& 15363 & 72.6 & 3600\textsuperscript{\textdagger} \\
\midrule

2023apr01\newline a\newline 2023may30 & 7101 & 96.3 & 3602\textsuperscript{\textdagger}& 8439 & 94.1 & 3603\textsuperscript{\textdagger}& 8468 & 98.8 & 3645\textsuperscript{\textdagger}& 11160 & 92.2 & 3600\textsuperscript{\textdagger}& 12543 & 80.8 & 3600\textsuperscript{\textdagger}& 13891 & 77.4 & 3600\textsuperscript{\textdagger}& 15237 & 70.6 & 3600\textsuperscript{\textdagger} \\
\midrule

2023jun01\newline a\newline 2023jul30 & 6602 & 95.9 & 3600\textsuperscript{\textdagger}& 7970 & 96.1 & 3602\textsuperscript{\textdagger}& 9288 & 92.4 & 3601\textsuperscript{\textdagger}& 10682 & 86.5 & 3636\textsuperscript{\textdagger}& 12148 & 72.9 & 3600\textsuperscript{\textdagger}& 13553 & 75.7 & 3623\textsuperscript{\textdagger}& 14916 & 72.1 & 3600\textsuperscript{\textdagger} \\
\midrule

\midrule
\textbf{M\'edia (10 inst.)} & $\overline{Z_1}{=}3421$ & 92.2$\pm$5.4 & 3600& $\overline{Z_1}{=}4716$ & 91.6$\pm$4.9 & 3601& $\overline{Z_1}{=}5871$ & 87.7$\pm$8.8 & 3605& $\overline{Z_1}{=}7327$ & 79.8$\pm$8.8 & 3609& $\overline{Z_1}{=}8625$ & 70.1$\pm$9.2 & 3600& $\overline{Z_1}{=}9905$ & 64.6$\pm$11.2 & 3603& $\overline{Z_1}{=}10931$ & 60.8$\pm$10.5 & 3600 \\

\bottomrule
\end{tabular}
}
}
	\fautor
\end{sidewaystable}

\clearpage

\section{Horizonte de 90 dias}
\label{apendice:exato-h90}

No horizonte trimestral, o gap MIP médio aproxima-se do limite teórico, com média de $96{,}80\%$ e amplitude reduzida ($\pm 6{,}22$ pontos percentuais), indicando uniformidade entre instâncias. O limitante inferior reportado pelo \textit{solver} é, em essência, próximo de zero ou da própria relaxação linear, de modo que o gap deixa de funcionar como medida útil de qualidade da solução encontrada. Sob esse regime, o método exato preserva a função de fornecer soluções factíveis na grade de $\varepsilon$, mas a sua fronteira de Pareto não pode mais ser interpretada como aproximação confiável da fronteira ótima global. A análise por instância (Figura~\ref{fig:exato-pareto-90dias}) revela a coexistência de fronteiras com formas e amplitudes distintas, espelhando a heterogeneidade dos perfis sazonais de demanda de 2022 e 2023.

O horizonte de planejamento de 90 dias apresenta o comportamento mais errático, evidenciado na Figura \ref{fig:exato-pareto-90dias}. Observa-se que na maioria das instâncias, o comportamento é das fronteiras é de pontos relativamente estáveis nos extremos, especialmente nos pontos $\epsilon_{1}$, com uma grande degradação na fronteira a partir dos pontos intermediários e extremo $\epsilon_{7}$, demonstrando a dificuldade oriunda da complexidade combinatória do problema de redução de trocas. Em outros casos, no entanto, observa-se o comportamento oposto, com \textit{gaps} elevados para os pontos iniciais, e melhor desempenho nos extremos de troca. Algumas instâncias, como \textbf{fevereiro de 2022 a maio de 2022} apresentam o comportamento de decrescimento esperado entre os pontos a medida que os passos deslocam-sem, admitindo menos exames atendidos em detrimento a mais trocas, com o adendo de que os pontos intermediários apresentam degradação significativa.

\begin{figure}[h]
	\centering
	\includegraphics[width=0.72\textwidth]{images/exato_pareto_overlay_90days.png}
	\caption{Fronteiras de Pareto sobrepostas no horizonte de 90 dias.}
	\label{fig:exato-pareto-90dias}
	\fautor
\end{figure}

\newpage

O mapa de calor do MIP gap, presente na Figura~\ref{fig:exato-heatmap-90dias}, revela apenas um caso de instância com ponto $\varepsilon$ infactível. No entanto, observa-se a grande maioria dos pontos tiveram \textit{gap} de 100\%, o que indica que o limitante inferior reportado pelo \textit{solver} é próximo de zero ou da própria relaxação linear, e que as soluções factíveis encontradas estão distantes desse limitante, reforçando a percepção de que o método exato, sob esse regime, não consegue explorar de forma eficaz a região próxima à fronteira ótima global. A heterogeneidade entre as instâncias é acentuada, com algumas apresentando formas mais irregulares e amplitudes maiores, o que sugere que os perfis sazonais de demanda em 2022 e 2023 geram variações significativas na dificuldade de resolução para esse horizonte.

\begin{figure}[h]
	\centering
	\includegraphics[width=0.7\textwidth]{images/exato_mipgap_heatmap_90days.png}
	\caption{Mapa de calor do MIP gap no horizonte de 90 dias.}
	\label{fig:exato-heatmap-90dias}
	\fautor
\end{figure}

Novamente, o diagrama de caixa do tempo de execução por ponto $\varepsilon$, apresentado na Figura~\ref{fig:exato-runtime-90dias}, mostra que os tempos de execução atingem o limite de 3.600~s para todas as instâncias em todos os pontos, o que reforça ainda mais a complexidade do problema, mesmo em um horizonte de planejamento de 90 dias. A presença dos \textit{outliers} no gráfico dizem respeito ao \textit{overhead} computacional significativo oriundo do tamanho do modelo em memória, que impedia o processo do \textit{solver} de terminar a tarefa exatamente em 3.600~s.

\newpage

\begin{figure}[t]
	\centering
	\includegraphics[width=0.7\textwidth]{images/exato_runtime_boxplot_90days.png}
	\caption{Tempo de execução por ponto $\varepsilon$ no horizonte de 90 dias.}
	\label{fig:exato-runtime-90dias}
	\fautor
\end{figure}

A Tabela~\ref{tab:exato-triplet-90dias}, ao final desta seção, apresenta os resultados completos por instância e ponto $\varepsilon$ no horizonte de 90 dias.

% Tabela
\clearpage

\begin{sidewaystable}[p]
	\centering
	\caption{Resultados por ponto $\varepsilon$ e instância, horizonte de 90 dias.}
	\label{tab:exato-triplet-90dias}
	\resizebox{\textheight}{!}{{\footnotesize
\begin{tabular}{@{}M{2.5cm}ccccccccccccccccccccc@{}}
    \toprule
    \multirow{2}{*}{\textbf{Inst\^ancia}} & \multicolumn{3}{c}{$\varepsilon_{1}$} & \multicolumn{3}{c}{$\varepsilon_{2}$} & \multicolumn{3}{c}{$\varepsilon_{3}$} & \multicolumn{3}{c}{$\varepsilon_{4}$} & \multicolumn{3}{c}{$\varepsilon_{5}$} & \multicolumn{3}{c}{$\varepsilon_{6}$} & \multicolumn{3}{c}{$\varepsilon_{7}$}                                                                                                                                                                                                                                                                                                                                                                             \\
    \cmidrule(lr){2-4}\cmidrule(lr){5-7}\cmidrule(lr){8-10}\cmidrule(lr){11-13}\cmidrule(lr){14-16}\cmidrule(lr){17-19}\cmidrule(lr){20-22}
                                          & $Z_1$                                 & gap\%                                 & t(s)                                  & $Z_1$                                 & gap\%                                 & t(s)                                  & $Z_1$                                 & gap\%        & t(s)                              & $Z_1$                    & gap\%        & t(s)                              & $Z_1$                    & gap\%        & t(s)                              & $Z_1$                    & gap\%        & t(s)                              & $Z_1$                    & gap\%         & t(s)                              \\
    \toprule
    2022jan02\newline a\newline 2022apr01 & 5149                                  & 96.5                                  & 3600\textsuperscript{\textdagger}     & 6875                                  & 97.0                                  & 3661\textsuperscript{\textdagger}     & 8303                                  & 96.2         & 3600\textsuperscript{\textdagger} & 9984                     & 99.5         & 3600\textsuperscript{\textdagger} & 11160                    & 91.7         & 3603\textsuperscript{\textdagger} & 12334                    & 90.2         & 3653\textsuperscript{\textdagger} & 13771                    & 88.2          & 3601\textsuperscript{\textdagger} \\
    \midrule

    2022feb01\newline a\newline 2022may01 & 5690                                  & 96.9                                  & 3624\textsuperscript{\textdagger}     & 7289                                  & 97.2                                  & 3600\textsuperscript{\textdagger}     & 8842                                  & 100.0        & 3600\textsuperscript{\textdagger} & 10053                    & 96.3         & 3600\textsuperscript{\textdagger} & 11860                    & 100.0        & 3600\textsuperscript{\textdagger} & 13369                    & 100.0        & 3600\textsuperscript{\textdagger} & 14670                    & 74.7          & 3645\textsuperscript{\textdagger} \\
    \midrule

    2022apr01\newline a\newline 2022jun29 & 5235                                  & 96.8                                  & 3600\textsuperscript{\textdagger}     & 6618                                  & 97.3                                  & 3600\textsuperscript{\textdagger}     & 8297                                  & 100.0        & 3600\textsuperscript{\textdagger} & 9694                     & 100.0        & 3600\textsuperscript{\textdagger} & 11108                    & 100.0        & 3600\textsuperscript{\textdagger} & 12457                    & 79.4         & 3600\textsuperscript{\textdagger} & 13782                    & 76.8          & 3600\textsuperscript{\textdagger} \\
    \midrule

    2022jun01\newline a\newline 2022aug29 & 6650                                  & 97.4                                  & 3600\textsuperscript{\textdagger}     & 8189                                  & 97.6                                  & 3608\textsuperscript{\textdagger}     & 9645                                  & 96.0         & 3600\textsuperscript{\textdagger} & 11245                    & 100.0        & 3600\textsuperscript{\textdagger} & 10976                    & 96.0         & 3647\textsuperscript{\textdagger} & 14185                    & 100.0        & 3600\textsuperscript{\textdagger} & 15538                    & 100.0         & 3600\textsuperscript{\textdagger} \\
    \midrule

    2022aug01\newline a\newline 2022oct29 & 5690                                  & 99.4                                  & 3600\textsuperscript{\textdagger}     & 7327                                  & 100.0                                 & 3600\textsuperscript{\textdagger}     & 8965                                  & 100.0        & 3600\textsuperscript{\textdagger} & 10534                    & 93.9         & 3601\textsuperscript{\textdagger} & 10313                    & 96.8         & 3659\textsuperscript{\textdagger} & 13878                    & 100.0        & 3600\textsuperscript{\textdagger} & 15516                    & 100.0         & 3600\textsuperscript{\textdagger} \\
    \midrule

    2022sep01\newline a\newline 2022nov29 & 6484                                  & 97.1                                  & 3600\textsuperscript{\textdagger}     & 7630                                  & 97.3                                  & 3600\textsuperscript{\textdagger}     & 6634                                  & 99.9         & 3600\textsuperscript{\textdagger} & 10425                    & 82.2         & 3600\textsuperscript{\textdagger} & 11750                    & 99.9         & 3600\textsuperscript{\textdagger} & 12884                    & 79.0         & 3600\textsuperscript{\textdagger} & 14337                    & 99.9          & 3600\textsuperscript{\textdagger} \\
    \midrule

    2022nov01\newline a\newline 2023jan29 & 4965                                  & 95.1                                  & 3616\textsuperscript{\textdagger}     & 6152                                  & 99.2                                  & 3601\textsuperscript{\textdagger}     & 7104                                  & 100.0        & 3600\textsuperscript{\textdagger} & 8708                     & 100.0        & 3600\textsuperscript{\textdagger} & 10620                    & 100.0        & 3600\textsuperscript{\textdagger} & 11247                    & 100.0        & 3600\textsuperscript{\textdagger} & 13665                    & 100.0         & 3600\textsuperscript{\textdagger} \\
    \midrule

    2023jan01\newline a\newline 2023mar31 & 12885                                 & 100.0                                 & 3600\textsuperscript{\textdagger}     & 14267                                 & 100.0                                 & 3600\textsuperscript{\textdagger}     & 15650                                 & 100.0        & 3600\textsuperscript{\textdagger} & 17033                    & 100.0        & 3600\textsuperscript{\textdagger} & 18416                    & 100.0        & 3600\textsuperscript{\textdagger} & 19799                    & 100.0        & 3600\textsuperscript{\textdagger} & 21057                    & 78.6          & 3600\textsuperscript{\textdagger} \\
    \midrule

    2023mar01\newline a\newline 2023may29 & ---                                   & ---                                   & 3600\textsuperscript{\textdagger}     & 13120                                 & 100.0                                 & 3600\textsuperscript{\textdagger}     & 15215                                 & 100.0        & 3600\textsuperscript{\textdagger} & 17310                    & 100.0        & 3600\textsuperscript{\textdagger} & 18578                    & 100.0        & 3600\textsuperscript{\textdagger} & 21499                    & 100.0        & 3600\textsuperscript{\textdagger} & 23594                    & 100.0         & 3600\textsuperscript{\textdagger} \\
    \midrule

    2023may01\newline a\newline 2023jul29 & 27133                                 & 100.0                                 & 3600\textsuperscript{\textdagger}     & 27133                                 & 100.0                                 & 3600\textsuperscript{\textdagger}     & 27133                                 & 100.0        & 3600\textsuperscript{\textdagger} & 27133                    & 100.0        & 3600\textsuperscript{\textdagger} & 27133                    & 100.0        & 3600\textsuperscript{\textdagger} & 27133                    & 100.0        & 3600\textsuperscript{\textdagger} & 27133                    & 100.0         & 3600\textsuperscript{\textdagger} \\
    \midrule

    \midrule
    \textbf{M\'edia (10 inst.)}           & $\overline{Z_1}{=}8876$               & 97.7$\pm$1.7                          & 3604                                  & $\overline{Z_1}{=}10460$              & 98.5$\pm$1.4                          & 3607                                  & $\overline{Z_1}{=}11579$              & 99.2$\pm$1.6 & 3600                              & $\overline{Z_1}{=}13212$ & 97.2$\pm$5.7 & 3600                              & $\overline{Z_1}{=}14191$ & 98.4$\pm$2.8 & 3611                              & $\overline{Z_1}{=}15878$ & 94.8$\pm$8.8 & 3605                              & $\overline{Z_1}{=}17306$ & 91.8$\pm$11.1 & 3605                              \\

    \bottomrule
\end{tabular}
}
}
	\fautor
	\label{LastContentPage}
\end{sidewaystable}

\clearpage

\end{apendicesenv}
% ---


% ----------------------------------------------------------
% Anexos
% ----------------------------------------------------------

% ---
% Inicia os anexos
% ---
\begin{anexosenv}

    % \chapter{Páginas interessantes na Internet} 
    % \label{chapter:paginas-interessantes}
    % \input{tex/5-anexos/paginas-interessantes}

\end{anexosenv}
% ---

\end{document}